% La notion de paix — ECM 1re Tech — Cours signé OnBuch+
% Programme officiel MINESEC. Compiler avec : tectonic N05-notion-paix.tex
\newcommand{\DOCMATIERE}{ECM}
\newcommand{\DOCNIVEAU}{1re Tech}
\newcommand{\DOCMODULE}{ECM – classe de Première technique}
\newcommand{\DOCLECON}{N05}
\newcommand{\DOCTITRE}{La notion de paix}
% OnBuch+ — préambule « cours signés » ENRICHI (compilé avec tectonic / XeLaTeX)
% Système visuel « Pop » d'OnBuch+ : fond crème (jamais blanc), bordures encre
% épaisses, ombres « dures » décalées, titres Archivo Black en capitales.
% Version MATHÉMATIQUES : graphes (pgfplots/tikz), boîtes pédagogiques
% (définition, propriété, méthode, attention, le savais-tu, exercice résolu,
% exercice, TP, bilan), page de garde et signature OnBuch+ ancrée en bas.
%
% Macros attendues dans main.tex AVANT \input{preamble} :
%   \DOCMATIERE  \DOCNIVEAU  \DOCMODULE  \DOCLECON (numéro)  \DOCTITRE
\documentclass[11pt]{article}

\usepackage{fontspec}
\usepackage[french]{babel}
\usepackage[a4paper,margin=2cm,top=3.4cm,bottom=2.8cm,headheight=1.4cm,footskip=1.3cm]{geometry}
\usepackage{amsmath,amssymb,amsfonts}
\usepackage{mathtools} % \xmapsto, \coloneqq... souvent utilisés par les modèles
\usepackage{cancel} % simplifications barrées (bilans de forces)
% \abs{z} (module, valeur absolue) et \norm{u} (norme), utilisés par les modèles
\providecommand{\abs}{}\let\abs\relax\DeclarePairedDelimiter{\abs}{\lvert}{\rvert}
\providecommand{\norm}{}\let\norm\relax\DeclarePairedDelimiter{\norm}{\lVert}{\rVert}
\usepackage[table]{xcolor} % [table] : fournit aussi \rowcolors (alternance de lignes)
\usepackage{pagecolor}
\usepackage{tikz}
\usetikzlibrary{arrows.meta,positioning,calc,shapes.geometric,shapes.misc,shapes.arrows,decorations.pathmorphing,decorations.pathreplacing,decorations.markings,patterns,shadows,mindmap,trees,fit,backgrounds,matrix,intersections,angles,quotes}
\usepackage{pgfplots}
\usepackage[version=4]{mhchem} % équations nucléaires \ce{^{235}_{92}U}
\usepackage[european,siunitx]{circuitikz} % schémas électriques
\pgfplotsset{compat=1.18}
\usepgfplotslibrary{fillbetween,groupplots} % aires entre courbes (calcul intégral)
\usepackage{enumitem}
\usepackage{fancyhdr}
% [most] charge toutes les bibliothèques tcolorbox (listings, minted, raster,
% documentation...) et gonfle inutilement la mémoire TeX sur un document long
% (~300 boîtes) : on ne charge que ce qui est réellement utilisé.
\usepackage[skins,breakable]{tcolorbox}
\usepackage{graphicx}
\usepackage{colortbl}
\usepackage{booktabs}
\usepackage{tabularx}
\usepackage{multirow}
\usepackage{diagbox} % case d'en-tête diagonale des tableaux à double entrée
% Colonnes extensibles centrée / gauche / droite (C, L, R), souvent utilisées par les modèles
\newcolumntype{C}{>{\centering\arraybackslash}X}
\newcolumntype{L}{>{\raggedright\arraybackslash}X}
\newcolumntype{R}{>{\raggedleft\arraybackslash}X}
\usepackage{array}
\usepackage{multicol}
\usepackage{siunitx}
% parse-numbers=false : accepte aussi \SI{2,5 \times 10^{-3}}{...}
\sisetup{locale=FR,output-decimal-marker={,},group-minimum-digits=4,parse-numbers=false}
\DeclareSIUnit\ohm{\ensuremath{\symup{\Omega}}} % U+2126 absent de la police
\DeclareSIUnit\division{div} % réglages d'oscilloscope (ms/div, V/div)
\DeclareSIUnit\tr{tr} % tours (vitesse de rotation, ex. \SI{1500}{\tr\per\minute})
\DeclareSIUnit\molar{M} % concentration molaire (ex. \SI{3}{\molar}, \SI{1}{\micro\molar})
\DeclareSIUnit\an{an} % année (le modèle confond parfois avec \year, un compteur TeX) — ex. \SI{5}{\kilo\meter\per\an}
\DeclareSIUnit\FCFA{FCFA} % franc CFA (ex. \SI{500}{\FCFA\per\kilo\gram})
\DeclareSIUnit\degreeBrix{\textdegree Brix} % teneur en sucre d'un sirop/jus (réfractométrie)
\DeclareSIUnit\month{mois} % le modèle utilise parfois \month, un compteur TeX, comme unité de durée
\DeclareSIUnit\microliter{\micro\liter} % le modèle écrit parfois \microliter en un seul token
\DeclareSIUnit\million{millions} % dénombrement (ex. \SI{15}{\million\per\mL} spermatozoïdes)
\usepackage{qrcode}
% \degree : commande fréquemment utilisée par les modèles pour un angle ou une
% température en dehors de \SI{}{} (non fournie par les paquets chargés).
\providecommand{\degree}{^{\circ}}
% \qed (fin de démonstration), utilisé par les modèles sans amsthm
\providecommand{\qed}{\hfill$\square$}
% \barre : double barre des valeurs interdites dans un tableau de variations (tkz-tab)
\providecommand{\barre}{\ensuremath{\|}}
% environnement proof (amsthm non chargé) utilisé par les modèles
\ifdefined\proof\else\newenvironment{proof}[1][Démonstration]{\par\noindent\textit{#1.}\ }{\qed\par}\fi
\usepackage{lastpage}
\usepackage{hyperref}
\hypersetup{hidelinks}

% ============================================================
% Palette « Pop » OnBuch+ — mêmes hex que la classe P (plus_ui.dart)
% ============================================================
\definecolor{poporange}{HTML}{F59321}
\definecolor{popdark}{HTML}{E07A0C}
\definecolor{popcream}{HTML}{FFF6E8}
\definecolor{popink}{HTML}{1C1714}
\definecolor{poppurple}{HTML}{7A5AE0}
\definecolor{popgreen}{HTML}{2E9E6A}
\definecolor{poppink}{HTML}{E0457B}
\definecolor{popblue}{HTML}{2F7DE1}
\definecolor{popgold}{HTML}{C98A12}
\definecolor{popmuted}{HTML}{8A7F76}
\definecolor{popline}{HTML}{EADFD2}
\definecolor{popgray}{HTML}{8A7F76}
\colorlet{popgrey}{popgray}
% Alias tolérants (le modèle nomme parfois les couleurs différemment)
\colorlet{popwhite}{white}
\colorlet{popblack}{popink}
\colorlet{popred}{poppink}
\colorlet{popyellow}{popgold}
% Teintes claires (fonds de tableaux/encadrés) — le modèle nomme parfois
% les couleurs avec un suffixe L (ex. popblueL) sans les avoir définies.
\colorlet{poporangeL}{poporange!20}
\colorlet{popdarkL}{popdark!20}
\colorlet{poppurpleL}{poppurple!20}
\colorlet{popgreenL}{popgreen!20}
\colorlet{poppinkL}{poppink!20}
\colorlet{popblueL}{popblue!20}
\colorlet{popgoldL}{popgold!20}
\colorlet{popredL}{popred!20}
\colorlet{popgrayL}{popgray!20}
% Teintes claires pour fonds de tableaux / zones
\colorlet{popblueL}{popblue!10!white}
\colorlet{poporangeL}{poporange!14!white}
\colorlet{popgreenL}{popgreen!12!white}
\colorlet{poppurpleL}{poppurple!12!white}
\colorlet{poppinkL}{poppink!12!white}
\colorlet{popgoldL}{popgold!14!white}

\pagecolor{popcream}

% ============================================================
% Typographie
% ============================================================
\defaultfontfeatures{Ligatures=TeX}
\setmainfont{PlusJakartaSans-Medium.ttf}[Path=../../../fonts/,BoldFont=PlusJakartaSans-Bold.ttf]
% Maths en sans-serif (Fira Math) avec chiffres et lettres droites de Plus
% Jakarta, pour que les formules se fondent dans le texte.
\usepackage[math-style=ISO,bold-style=ISO]{unicode-math}
\setmathfont{FiraMath-Regular.otf}[Path=../../../fonts/]
\setmathfont{PlusJakartaSans-Medium.ttf}[Path=../../../fonts/,range={up/num,up/{Latin,latin},\mathup}]
\usepackage{icomma} % virgule décimale sans espace en mode math
% Caractères mathématiques Unicode absents des polices de texte (Plus Jakarta,
% Archivo Black) : rendus par leur équivalent mathématique, en texte comme en
% formule — sinon ils s'affichent en carré vide (ex. « Arithmétique dans ℤ »).
\usepackage{newunicodechar}
\newunicodechar{ℝ}{\ensuremath{\mathbb{R}}}
\newunicodechar{ℤ}{\ensuremath{\mathbb{Z}}}
\newunicodechar{ℂ}{\ensuremath{\mathbb{C}}}
\newunicodechar{ℕ}{\ensuremath{\mathbb{N}}}
\newunicodechar{ℚ}{\ensuremath{\mathbb{Q}}}
\newunicodechar{𝔻}{\ensuremath{\mathbb{D}}}
\newunicodechar{α}{\ensuremath{\alpha}}
\newunicodechar{β}{\ensuremath{\beta}}
\newunicodechar{γ}{\ensuremath{\gamma}}
\newunicodechar{δ}{\ensuremath{\delta}}
\newunicodechar{ε}{\ensuremath{\varepsilon}}
\newunicodechar{θ}{\ensuremath{\theta}}
\newunicodechar{λ}{\ensuremath{\lambda}}
\newunicodechar{μ}{\ensuremath{\mu}}
\newunicodechar{σ}{\ensuremath{\sigma}}
\newunicodechar{τ}{\ensuremath{\tau}}
\newunicodechar{φ}{\ensuremath{\varphi}}
\newunicodechar{ψ}{\ensuremath{\psi}}
\newunicodechar{ω}{\ensuremath{\omega}}
\newunicodechar{Σ}{\ensuremath{\Sigma}}
% majuscules produites par \MakeUppercase dans les titres (α-aminés -> Α)
\newunicodechar{Α}{\ensuremath{\alpha}}
\newunicodechar{Β}{\ensuremath{\beta}}
\newunicodechar{Γ}{\ensuremath{\Gamma}}
\newunicodechar{Δ}{\ensuremath{\Delta}}
\newunicodechar{Ε}{\ensuremath{\varepsilon}}
\newunicodechar{Θ}{\ensuremath{\Theta}}
\newunicodechar{Λ}{\ensuremath{\Lambda}}
\newunicodechar{Μ}{\ensuremath{\mu}}
\newunicodechar{Τ}{\ensuremath{\tau}}
\newunicodechar{Φ}{\ensuremath{\Phi}}
\newunicodechar{Ψ}{\ensuremath{\Psi}}
\newunicodechar{Ω}{\ensuremath{\Omega}}
\newunicodechar{↦}{\ensuremath{\mapsto}}
\newunicodechar{∈}{\ensuremath{\in}}
\newunicodechar{∉}{\ensuremath{\notin}}
\newunicodechar{∧}{\ensuremath{\wedge}}
\newunicodechar{⇔}{\ensuremath{\Leftrightarrow}}
\newunicodechar{⟺}{\ensuremath{\Longleftrightarrow}}
\newunicodechar{⇒}{\ensuremath{\Rightarrow}}
\newunicodechar{⟹}{\ensuremath{\Longrightarrow}}
\newunicodechar{∘}{\ensuremath{\circ}}
\newunicodechar{∪}{\ensuremath{\cup}}
\newunicodechar{∩}{\ensuremath{\cap}}
\newunicodechar{≡}{\ensuremath{\equiv}}
\newunicodechar{⊂}{\ensuremath{\subset}}
\newunicodechar{⊥}{\ensuremath{\perp}}
\newunicodechar{∃}{\ensuremath{\exists}}
\newunicodechar{∀}{\ensuremath{\forall}}
\newunicodechar{∛}{\ensuremath{\sqrt[3]{\,}}}
\newunicodechar{∜}{\ensuremath{\sqrt[4]{\,}}}
\newunicodechar{⃗}{\ensuremath{{}^{\to}}}
\newunicodechar{ⁿ}{\ifmmode^{n}\else\textsuperscript{\ensuremath{n}}\fi}
\newunicodechar{ˣ}{\ifmmode^{x}\else\textsuperscript{\ensuremath{x}}\fi}
\newunicodechar{ᵇ}{\ifmmode^{b}\else\textsuperscript{\ensuremath{b}}\fi}
\newunicodechar{ʳ}{\ifmmode^{r}\else\textsuperscript{\ensuremath{r}}\fi}
\newunicodechar{ᵘ}{\ifmmode^{u}\else\textsuperscript{\ensuremath{u}}\fi}
\newunicodechar{ʸ}{\ifmmode^{y}\else\textsuperscript{\ensuremath{y}}\fi}
\newunicodechar{ᵃ}{\ifmmode^{a}\else\textsuperscript{\ensuremath{a}}\fi}
\newunicodechar{ᵉ}{\ifmmode^{e}\else\textsuperscript{\ensuremath{e}}\fi}
\newunicodechar{ᵏ}{\ifmmode^{k}\else\textsuperscript{\ensuremath{k}}\fi}
\newunicodechar{ᵖ}{\ifmmode^{p}\else\textsuperscript{\ensuremath{p}}\fi}
\newunicodechar{ᴺ}{\ifmmode^{N}\else\textsuperscript{\ensuremath{N}}\fi}
\newunicodechar{ᶜ}{\ifmmode^{c}\else\textsuperscript{\ensuremath{c}}\fi}
\newunicodechar{ʰ}{\ifmmode^{h}\else\textsuperscript{\ensuremath{h}}\fi}
\newunicodechar{ᵠ}{\ifmmode^{q}\else\textsuperscript{\ensuremath{q}}\fi}
\newunicodechar{ₙ}{\ifmmode_{n}\else\textsubscript{\ensuremath{n}}\fi}
\newunicodechar{ₐ}{\ifmmode_{a}\else\textsubscript{\ensuremath{a}}\fi}
\newunicodechar{ᵢ}{\ifmmode_{i}\else\textsubscript{\ensuremath{i}}\fi}
\newunicodechar{ᵣ}{\ifmmode_{r}\else\textsubscript{\ensuremath{r}}\fi}
\newunicodechar{ₖ}{\ifmmode_{k}\else\textsubscript{\ensuremath{k}}\fi}
\newunicodechar{ᵦ}{\ifmmode_{\beta}\else\textsubscript{\ensuremath{\beta}}\fi}
\newunicodechar{ₓ}{\ifmmode_{x}\else\textsubscript{\ensuremath{x}}\fi}
\newfontfamily\popdisplay{ArchivoBlack-Regular.ttf}[Path=../../../fonts/]
\newfontfamily\poplogo{SpaceGrotesk-Bold.ttf}[Path=../../../fonts/]
\newfontfamily\popsemi{PlusJakartaSans-SemiBold.ttf}[Path=../../../fonts/]
\newfontfamily\popxbold{PlusJakartaSans-ExtraBold.ttf}[Path=../../../fonts/]
\newfontfamily\popxboldls{PlusJakartaSans-ExtraBold.ttf}[Path=../../../fonts/,LetterSpace=3.0]

\setlist[enumerate]{leftmargin=1.7em,itemsep=5pt,topsep=4pt}
\setlist[itemize]{leftmargin=1.7em,itemsep=4pt,topsep=4pt,label={\color{poporange}\textbullet}}
\setlength{\parindent}{0pt}
\setlength{\parskip}{5pt}
\renewcommand{\arraystretch}{1.25}

% pgfplots : style Pop par défaut
\pgfplotsset{
  every axis/.append style={axis line style={popink,line width=1pt},tick style={popink},
    grid style={popline,line width=0.5pt},label style={font=\small},tick label style={font=\footnotesize}},
  popcurve/.style={poporange,line width=1.8pt,no markers,smooth},
}
% Même style disponible dans un \draw TikZ simple (hors pgfplots)
\tikzset{popcurve/.style={poporange,line width=1.8pt}}
\tikzset{popgrid/.style={popline,very thin}}
\colorlet{popcurve}{poporange} % le nom du style est parfois utilisé comme couleur

% ============================================================
% Pill / squiggle
% ============================================================
\newcommand{\poppill}[3]{%
  \tikz[baseline=-0.55ex]\node[draw=popink,line width=1.2pt,rounded corners=2.6pt,%
    fill=#2,inner xsep=6.5pt,inner ysep=3pt]{%
    \popxboldls\fontsize{7}{7}\selectfont\color{#3}\MakeUppercase{#1}};%
}
\newcommand{\popsquiggle}[1]{%
  \tikz[baseline=-0.4ex]{
    \draw[#1,line width=2.6pt,line cap=round]
      plot [smooth] coordinates {(0,0.09) (0.34,0.01) (0.68,0.21) (1.02,0.01) (1.36,0.10)};
  }%
}
% Mot-clé surligné (marqueur orange sous le texte)
\newcommand{\cle}[1]{{\popxbold\color{popdark}#1}}

% ============================================================
% En-tête / pied
% ============================================================
\pagestyle{fancy}
\fancyhf{}
\renewcommand{\headrulewidth}{0pt}
\renewcommand{\footrulewidth}{0pt}
\lhead{\raisebox{-0.15ex}{\poplogo\fontsize{13}{13}\selectfont\color{popink}OnBuch\color{poporange}+}}
\rhead{\poppill{\DOCMATIERE\ \textbullet\ \DOCNIVEAU\ \textbullet\ Leçon \DOCLECON}{white}{popink}}
\lfoot{\popxbold\fontsize{7.6}{7.6}\selectfont\color{popmuted}\MakeUppercase{Cours signé OnBuch+}\ \textbullet\ {\popsemi\DOCTITRE}}
\rfoot{\tikz[baseline=-0.6ex]\node[rounded corners=8.5pt,fill=popink,minimum height=17pt,minimum width=17pt,inner xsep=5pt,inner sep=0]{\popxbold\fontsize{8}{8}\selectfont\color{popcream}\thepage\,/\,\pageref*{LastPage}};}

% ============================================================
% Titres
% ============================================================
\newcommand{\chaptitre}[2]{%
  \begin{tcolorbox}[enhanced,colback=poporange,colframe=popink,boxrule=2.6pt,arc=11pt,
      drop shadow=popink,left=15pt,right=15pt,top=12pt,bottom=11pt,before skip=3pt,after skip=18pt]
    \poppill{#2}{popink}{popcream}\par\vspace{9pt}
    {\raggedright\hyphenpenalty=10000\exhyphenpenalty=10000\popdisplay\fontsize{22}{24}\selectfont\color{popcream}\MakeUppercase{#1}\par}
  \end{tcolorbox}
}
% Section numérotée : pastille orange avec le numéro + titre Archivo
\newcounter{popsec}
\newcommand{\coursec}[1]{%
  \stepcounter{popsec}\setcounter{popsub}{0}%
  \par\vspace{16pt}\noindent
  \begin{minipage}{\linewidth}
  \tikz[baseline=-0.7ex]\node[draw=popink,line width=1.6pt,fill=poporange,rounded corners=4pt,minimum size=22pt,inner sep=0]{\popdisplay\fontsize{12}{12}\selectfont\color{popcream}\thepopsec};%
  \hspace{8pt}{\popdisplay\fontsize{15}{17}\selectfont\color{popink}\MakeUppercase{#1}}\par
  \vspace{4pt}\noindent{\color{poporange}\rule{36pt}{3.6pt}}
  \end{minipage}\par\nopagebreak\vspace{8pt}
}
\newcounter{popsub}
\newcommand{\courssub}[1]{%
  \stepcounter{popsub}%
  \par\vspace{9pt}\noindent{\popxbold\fontsize{11.5}{13}\selectfont\color{popink}{\color{poporange}\thepopsec.\thepopsub}\hspace{6pt}#1}\par\nopagebreak\vspace{4pt}
}

% ============================================================
% Boîtes pédagogiques — même recette : fond blanc, bordure épaisse colorée,
% ombre dure encre, badge plein. Titre optionnel [#1].
% ============================================================
\tcbset{popbase/.style={breakable,enhanced,colback=white,boxrule=2.3pt,arc=9pt,
  shadow={3pt}{-3pt}{0pt}{popink},left=14pt,right=14pt,top=11pt,bottom=11pt,before skip=11pt,after skip=15pt,
  fontupper=\popsemi\fontsize{10.6}{14.5}\selectfont\color{popink},
  fontlower=\popsemi\fontsize{10.6}{14.5}\selectfont\color{popink}}}
% Coche dessinée (le glyphe ✓ n'existe pas dans les polices du document)
\newcommand{\popcheck}[1]{\tikz[baseline=-0.5ex]\draw[#1,line width=1.6pt,line cap=round,line join=round](-0.13,0.0)--(-0.04,-0.09)--(0.13,0.1);}
\providecommand{\coche}{\popcheck{popgreen}}
\providecommand{\resultat}[1]{\par\smallskip\noindent\fcolorbox{popgreen}{popgreenL}{\parbox{\dimexpr\linewidth-2\fboxsep-2\fboxrule\relax}{\textbf{Résultat.} #1}}\par\smallskip}
\newcommand{\popboxhead}[3]{\poppill{#1}{#2}{white}\ifx\relax#3\relax\else\hspace{6pt}{\popxbold\fontsize{10.6}{12}\selectfont\color{popink}#3}\fi\par\vspace{7pt}}

\newtcolorbox{aretenir}[1][]{popbase,colframe=popgreen,before upper={\popboxhead{À retenir}{popgreen}{#1}}}
\newtcolorbox{exemplebox}[1][]{popbase,colframe=poporange,before upper={\popboxhead{Exemple}{poporange}{#1}}}
\newtcolorbox{definition}[1][]{popbase,colframe=popblue,before upper={\popboxhead{Définition}{popblue}{#1}}}
\newtcolorbox{propriete}[1][]{popbase,colframe=poppurple,before upper={\popboxhead{Propriété}{poppurple}{#1}}}
\newtcolorbox{methode}[1][]{popbase,colframe=popgold,before upper={\popboxhead{Méthode}{popgold}{#1}}}
\newtcolorbox{attention}[1][]{popbase,colframe=poppink,before upper={\popboxhead{Attention}{poppink}{#1}}}
\newtcolorbox{experience}[1][]{popbase,colframe=popblue,colback=popblueL,before upper={\popboxhead{Expérience}{popblue}{#1}}}
% « Le savais-tu ? » : carte violette pleine (contexte, culture, Cameroun)
\newtcolorbox{savaistu}[1][]{popbase,colback=poppurple,colframe=popink,
  fontupper=\popsemi\fontsize{10.4}{14.2}\selectfont\color{white},
  before upper={\poppill{Le savais-tu\,?}{white}{poppurple}\ifx\relax#1\relax\else\hspace{6pt}{\popxbold\color{white}#1}\fi\par\vspace{7pt}}}
% Exercice résolu : énoncé en haut, \tcblower puis la solution sur fond crème
\newcounter{popexo}\newcounter{popexr}
\newtcolorbox{exoresolu}[1][]{popbase,colframe=poporange,colbacklower=popcream,
  segmentation style={popink,line width=1.2pt,dashed},
  before upper={\stepcounter{popexr}\popboxhead{Exercice résolu \thepopexr}{poporange}{#1}},
  before lower={\poppill{Solution}{popink}{popcream}\par\vspace{6pt}}}
% Exercice (non résolu en place) — la correction est dans la partie Corrigés
\newtcolorbox{exercice}[1][]{popbase,colframe=popink,
  before upper={\stepcounter{popexo}\popboxhead{Exercice \thepopexo}{popink}{#1}}}
% Corrigé
\newtcolorbox{corrige}[1][]{popbase,colframe=popgreen,colback=popgreenL,
  before upper={\popboxhead{Corrigé}{popgreen}{#1}}}
% Objectifs / prérequis / bilan
\newtcolorbox{objectifs}{popbase,colframe=popink,colback=poporangeL,
  before upper={\popboxhead{Objectifs de la leçon}{popink}{}}}
\newtcolorbox{prerequis}{popbase,colframe=popink,colback=popblueL,
  before upper={\popboxhead{Prérequis}{popblue}{}}}
\newtcolorbox{bilan}{popbase,colframe=popink,colback=white,boxrule=2.8pt,
  before upper={\popboxhead{Fiche bilan}{poporange}{L'essentiel à connaître pour l'examen}}}
% Figure encadrée avec légende
\newtcolorbox{popfigure}[1][]{enhanced,breakable=false,colback=white,colframe=popline,boxrule=1.4pt,arc=8pt,
  left=10pt,right=10pt,top=10pt,bottom=8pt,before skip=10pt,after skip=12pt,halign=center,
  lower separated=false,halign lower=center,
  fontlower=\popsemi\fontsize{9}{11.5}\selectfont\color{popmuted},
  #1}
\newcommand{\legende}[1]{\par\vspace{6pt}{\popsemi\fontsize{9}{11.5}\selectfont\color{popmuted}#1}}

% ============================================================
% Signature OnBuch+
% ============================================================
\newcommand{\poplogoinline}[1]{{\poplogo\fontsize{#1}{#1}\selectfont\color{popink}OnBuch\color{poporange}+}}
\newcommand{\signatureonbuch}{%
  \begin{tcolorbox}[enhanced,colback=popink,colframe=popink,boxrule=2.2pt,arc=10pt,
      drop shadow=poporange,left=14pt,right=14pt,top=10pt,bottom=10pt,before skip=0pt,after skip=0pt]
    \tikz[baseline=-0.7ex]\node[circle,fill=popgreen,draw=popcream,line width=1.3pt,minimum size=22pt,inner sep=0]{\popcheck{white}};%
    \hspace{9pt}\begin{minipage}[c]{0.62\linewidth}
      {\popxbold\fontsize{12}{13}\selectfont\color{popcream}Signé\ \textbullet\ {\poplogo OnBuch\color{poporange}+}}\par\vspace{2pt}
      {\raggedright\popsemi\fontsize{8}{10.5}\selectfont\color{popline}\MakeUppercase{Équipe pédagogique OnBuch+}\\\MakeUppercase{Conforme au programme officiel MINESEC}\par}
    \end{minipage}\hfill
    {\poplogo\fontsize{22}{22}\selectfont\color{popcream}OnBuch\color{poporange}+}
  \end{tcolorbox}%
}

% ============================================================
% Page de garde
% #1 titre · #2 matière · #3 niveau · #4 module · #5 n° leçon · #6 durée ·
% #7 description / objectifs (texte ou itemize)
% ============================================================
\newcommand{\pagegarde}[7]{%
  \thispagestyle{empty}
  \newgeometry{a4paper,margin=1.7cm,top=1.6cm,bottom=1.4cm}%
  \begin{tikzpicture}[remember picture,overlay]
    \fill[poporange!18] ([xshift=-3.2cm,yshift=-3.6cm]current page.north east) circle (4.6cm);
    \fill[poppurple!14] ([xshift=2.4cm,yshift=7.6cm]current page.south west) circle (3.8cm);
  \end{tikzpicture}%
  {\poplogo\fontsize{30}{30}\selectfont\color{popink}OnBuch\color{poporange}+}\hfill
  \raisebox{0.6ex}{\poppill{Cours signé \textbullet\ Premium}{popink}{popcream}}\par
  \vspace{1pt}\popsquiggle{poppurple}\par
  \vspace{8pt}
  \begin{tcolorbox}[enhanced,colback=poporange,colframe=popink,boxrule=2.8pt,arc=13pt,
      drop shadow=popink,left=18pt,right=18pt,top=12pt,bottom=13pt,after skip=10pt]
    \poppill{#2 \textbullet\ #3}{popink}{popcream}\hspace{5pt}\poppill{#4}{white}{popink}\par\vspace{8pt}
    {\popdisplay\fontsize{14}{14}\selectfont\color{popink}LEÇON #5}\par\vspace{4pt}
    {\raggedright\hyphenpenalty=10000\exhyphenpenalty=10000\popdisplay\fontsize{25}{27}\selectfont\color{popcream}\MakeUppercase{#1}\par}
    \vspace{8pt}
    {\popxbold\fontsize{9.5}{9.5}\selectfont\color{popink}\MakeUppercase{Durée conseillée : #6}}
  \end{tcolorbox}
  \begin{tcolorbox}[enhanced,colback=white,colframe=popink,boxrule=2.2pt,arc=10pt,
      drop shadow=popink,left=16pt,right=16pt,top=10pt,bottom=10pt,after skip=8pt]
    \poppill{Dans cette leçon}{poppurple}{white}\par\vspace{5pt}
    \popsemi\fontsize{9.3}{12.6}\selectfont\color{popink}#7
  \end{tcolorbox}
  \noindent\begin{minipage}[t]{0.487\linewidth}
    \begin{tcolorbox}[enhanced,colback=popgreen,colframe=popink,boxrule=2.2pt,arc=10pt,
        drop shadow=popink,left=11pt,right=11pt,top=10pt,bottom=10pt,height=3.1cm,valign=center]
      \tcbox[colback=white,colframe=popink,boxrule=1.3pt,arc=5pt,boxsep=0pt,nobeforeafter,
        left=3pt,right=3pt,top=3pt,bottom=3pt,valign=center]{\qrcode[height=1.9cm]{https://play.google.com/store/apps/details?id=com.luvvix.onbuch&hl=fr}}%
      \hspace{8pt}\begin{minipage}[c]{0.5\linewidth}
        {\popdisplay\fontsize{11}{12.5}\selectfont\color{white}\MakeUppercase{Télécharge OnBuch}}\par\vspace{4pt}
        {\raggedright\popsemi\fontsize{8.4}{11}\selectfont\color{white}Gratuit \textbullet\ Google Play\\Tuteur IA, annales, concours, résultats\par}
      \end{minipage}
    \end{tcolorbox}
  \end{minipage}\hfill
  \begin{minipage}[t]{0.487\linewidth}
    \begin{tcolorbox}[enhanced,colback=poppurple,colframe=popink,boxrule=2.2pt,arc=10pt,
        drop shadow=popink,left=11pt,right=11pt,top=10pt,bottom=10pt,height=3.1cm,valign=center]
      \tikz[baseline=-0.7ex]\node[circle,fill=white,draw=popink,line width=1.3pt,minimum size=1.6cm,inner sep=0]{\popdisplay\fontsize{24}{24}\selectfont\color{poppurple}+};%
      \hspace{8pt}\begin{minipage}[c]{0.6\linewidth}
        {\popdisplay\fontsize{11}{12.5}\selectfont\color{white}\MakeUppercase{Passe à OnBuch+}}\par\vspace{4pt}
        {\raggedright\popsemi\fontsize{8.4}{11}\selectfont\color{white}Tous les cours signés, Tuteur IA illimité, fiches PDF — onglet OnBuch+ de l'app\par}
      \end{minipage}
    \end{tcolorbox}
  \end{minipage}
  \vspace{8pt}
  \popcontenu
  \vfill
  \signatureonbuch
  \restoregeometry
  \newpage
}

% Rangée « Ce cours contient » (page de garde)
\newcommand{\poptile}[3]{% #1 couleur #2 gros texte #3 légende
  \begin{tcolorbox}[enhanced,colback=#1,colframe=popink,boxrule=2pt,arc=9pt,shadow={2.5pt}{-2.5pt}{0pt}{popink},
      left=6pt,right=6pt,top=9pt,bottom=9pt,halign=center,valign=center,height=1.75cm,width=\linewidth,nobeforeafter]
    {\popdisplay\fontsize{12.5}{13}\selectfont\color{white}\MakeUppercase{#2}}\par\vspace{3pt}
    {\centering\popsemi\fontsize{7.8}{9.5}\selectfont\color{white}#3\par}
  \end{tcolorbox}}
\newcommand{\popcontenu}{%
  \par\noindent{\popxboldls\fontsize{7.5}{7.5}\selectfont\color{popmuted}CE COURS SIGNÉ CONTIENT}\par\vspace{7pt}
  \noindent\begin{minipage}[t]{0.235\linewidth}\poptile{popblue}{Cours}{complet, illustré, pas à pas}\end{minipage}\hfill
  \begin{minipage}[t]{0.235\linewidth}\poptile{poporange}{Méthodes}{et exercices résolus}\end{minipage}\hfill
  \begin{minipage}[t]{0.235\linewidth}\poptile{poppink}{Exercices}{type examen + corrigés}\end{minipage}\hfill
  \begin{minipage}[t]{0.235\linewidth}\poptile{popgreen}{Fiche bilan}{l'essentiel à retenir}\end{minipage}\par}

% Sommaire
\newtcolorbox{sommaire}{enhanced,colback=white,colframe=popink,boxrule=2.2pt,arc=10pt,
  drop shadow=popink,left=18pt,right=18pt,top=13pt,bottom=13pt,before skip=6pt,after skip=20pt,
  before upper={\poppill{Sommaire}{poppurple}{white}\par\vspace{9pt}\popsemi\fontsize{10.6}{14.5}\selectfont\color{popink}}}

% Signature de fin de cours — poussée tout en bas de la dernière page
\newcommand{\findecours}{%
  \par\vfill\nopagebreak
  \begin{center}\popsquiggle{poporange}\end{center}\vspace{4pt}
  \signatureonbuch
}
\AtBeginDocument{\def\llbracket{\mathopen{[\mkern-3.5mu[}}\def\rrbracket{\mathclose{]\mkern-3.5mu]}}} % intervalles d'entiers (glyphe ⟦ absent de la police)
% Glyphes absents de Fira Math : redéfinis à partir de glyphes disponibles
\AtBeginDocument{%
  \def\setminus{\mathbin{\backslash}}\let\smallsetminus\setminus
  \def\vdots{\mathord{\vbox{\baselineskip3.2pt\lineskiplimit0pt\kern4pt\hbox{.}\hbox{.}\hbox{.}}}}%
  \def\ddots{\mathinner{\mkern1mu\raise6.5pt\hbox{.}\mkern2mu\raise3.6pt\hbox{.}\mkern2mu\raise0.7pt\hbox{.}\mkern1mu}}%
  \def\triangle{\mathord{\scalebox{0.9}{$\symup{\Delta}$}}}%
  \def\nabla{\mathord{\raisebox{\depth}{\scalebox{1}[-1]{$\symup{\Delta}$}}}}%
  \def\checkmark{\popcheck{popdark}}%
  \def\star{\mathbin{\ast}}\def\bigstar{\mathord{\ast}}%
  \def\diamond{\mathbin{\scalebox{0.6}{$\Diamond$}}}%
  \def\bigcup{\mathop{\mathchoice{\raisebox{-0.3em}{\scalebox{1.7}{$\cup$}}}{\scalebox{1.25}{$\cup$}}{\cup}{\cup}}\displaylimits}%
  \def\bigcap{\mathop{\mathchoice{\raisebox{-0.3em}{\scalebox{1.7}{$\cap$}}}{\scalebox{1.25}{$\cap$}}{\cap}{\cap}}\displaylimits}%
  \def\lesseqqgtr{\mathrel{\lessgtr}}\def\bigoplus{\mathop{\oplus}\displaylimits}%
}
\providecommand{\ssi}{si et seulement si} % abréviation fréquente
\AtBeginDocument{\providecommand{\wideparen}{\overparen}} % arc de cercle

\begin{document}

\pagegarde{La notion de paix}{ECM}{1re Tech}{ECM – classe de Première technique}{N05}{7 séances (fiche de progression harmonisée)}{Cette leçon définit la notion de paix et explore ses symboles, ses stratégies de promotion et les mécanismes de règlement pacifique des différends au Cameroun. Elle montre, à travers des documents et des enquêtes, comment la lutte contre l'extrémisme violent et le vivre ensemble (bilinguisme, multiculturalisme) consolident l'unité nationale.
\vspace{4pt}
\begin{multicols}{2}\small
\begin{itemize}
\item À la fin de cette leçon, je sais définir la paix et distinguer paix négative et paix positive
\item Identifier et expliquer les symboles de la paix au Cameroun à partir d'une enquête documentaire
\item Décrire les stratégies de promotion de la paix mises en œuvre au Cameroun
\item Expliquer les mécanismes de règlement pacifique des différends (négociation, médiation, arbitrage, justice)
\item Analyser les causes et les conséquences de l'extrémisme violent au Cameroun et les réponses de l'État
\item Justifier l'importance du vivre ensemble pour l'unité nationale
\end{itemize}\end{multicols}}

\begin{sommaire}
\begin{multicols}{2}
\begin{enumerate}
\item La notion de paix
\item TD4 : Les symboles de la paix au Cameroun (enquêtes)
\item Les stratégies de promotion de la paix
\item Le règlement pacifique des différends
\item Dossier 2 : La lutte contre l'extrémisme violent au Cameroun
\item Le vivre ensemble et le bilinguisme (Dossier 3)
\item Activité d'intégration
\item Méthodes et exercices résolus
\item Exercices
\item Corrigés des exercices
\item Fiche bilan
\end{enumerate}
\end{multicols}
\end{sommaire}

\begin{objectifs}
\begin{itemize}
\item À la fin de cette leçon, je sais définir la paix et distinguer paix négative et paix positive
\item Identifier et expliquer les symboles de la paix au Cameroun à partir d'une enquête documentaire
\item Décrire les stratégies de promotion de la paix mises en œuvre au Cameroun
\item Expliquer les mécanismes de règlement pacifique des différends (négociation, médiation, arbitrage, justice)
\item Analyser les causes et les conséquences de l'extrémisme violent au Cameroun et les réponses de l'État
\item Justifier l'importance du vivre ensemble pour l'unité nationale
\item Présenter le rôle de la Commission Nationale pour la Promotion du Bilinguisme et du Multiculturalisme (CNPBM)
\item Appliquer les notions de paix et d'unité à des situations concrètes de la vie courante
\item Rédiger et justifier une démarche, une réponse ou une conclusion argumentée
\end{itemize}
\end{objectifs}

\begin{prerequis}
\begin{itemize}
\item La notion d'État et les institutions de la République (classe de Seconde)
\item Les droits et devoirs du citoyen (ECM, Seconde)
\item La notion d'unité nationale et de diversité culturelle du Cameroun
\item Les grandes régions du Cameroun et leur diversité linguistique (Géographie, 4e-3e)
\end{itemize}
\end{prerequis}

\newpage

\coursec{La notion de paix}

En 2019, le Grand Dialogue National est organisé à Yaoundé pour répondre à la crise dans les régions du Nord-Ouest et du Sud-Ouest. La même année, des élèves d'un lycée de Buea reprennent le chemin de l'école sous escorte, tandis qu'à Maroua, des comités de vigilance collaborent avec les forces de défense contre Boko Haram. Comment un pays de plus de 250 groupes ethniques et deux langues officielles peut-il construire et préserver la paix ? Quels symboles, quelles institutions et quelles stratégies le Cameroun mobilise-t-il pour faire vivre le « vivre ensemble » ? C'est à ces questions que cette leçon apporte des réponses.

\courssub{Définition et dimensions de la paix}

Au sens le plus courant, la paix évoque la fin des combats, le silence des armes. Pourtant, l'histoire et la sociologie nous enseignent qu'un cessez-le-feu ne suffit pas à garantir une société stable. Une paix durable repose sur des piliers invisibles : la justice, l'équité, la confiance entre groupes humains.

\begin{definition}[La paix : définition globale]
La \cle{paix} est un état de concorde et de sécurité caractérisé par l'absence de conflit violent et la présence effective de la \cle{justice sociale}. Elle ne se réduit pas à l'absence de guerre ; elle suppose l'instauration d'un ordre juste où les droits fondamentaux sont respectés, où les ressources sont équitablement réparties et où les mécanismes de règlement des différends fonctionnent sans violence.
\end{definition}

Le sociologue norvégien \cle{Johan Galtung}, fondateur des études sur la paix, a distingué deux dimensions complémentaires qui structurent toute analyse rigoureuse :

\begin{itemize}
    \item La \cle{paix négative} : c'est l'absence de violence directe, de guerre ouverte, de combats armés. C'est un « silence des armes ». Elle est nécessaire mais insuffisante ; une dictature sans opposition armée peut connaître une paix négative tout en opprimant son peuple.
    \item La \cle{paix positive} : c'est la présence de structures justes, équitables et participatives. Elle intègre la justice sociale, le respect des droits humains, le développement équitable, la participation citoyenne et l'harmonie culturelle. C'est la paix « structurelle » qui prévient les causes profondes du conflit.
\end{itemize}

\begin{savaistu}[Johan Galtung, père de la distinction paix négative / paix positive]
Né en 1930, Johan Galtung est un mathématicien et sociologue norvégien, fondateur de l'Institut international de recherche sur la paix (PRIO) et du Journal of Peace Research. Dans son ouvrage fondateur \emph{Violence, Peace and Peace Research} (1969), il théorise la \cle{violence structurelle} (pauvreté, inégalités, discrimination) et la \cle{violence culturelle} (idéologies justifiant la domination). Pour lui, la paix positive exige de s'attaquer à ces violences invisibles. Il a médité des conflits majeurs (Israël/Palestine, Yougoslavie, Corée) et conseillé l'ONU. Sa grille de lecture est aujourd'hui la référence mondiale en ECM et en relations internationales.
\end{savaistu}

\begin{attention}[Piège fréquent : réduire la paix à l'absence de guerre]
L'erreur classique, y compris dans les copies d'examen, consiste à définir la paix uniquement comme « l'absence de conflit armé ». Cette définition correspond à la \textbf{paix négative} seulement. Elle ignore la \textbf{paix positive} (justice, équité, développement, participation). Au Probatoire, une définition complète doit mentionner les deux dimensions et citer Galtung. Exemple de formulation attendue : « La paix est un état dynamique combinant l'absence de violence directe (paix négative) et la présence de conditions structurelles justes (paix positive) selon Johan Galtung. »
\end{attention}

\courssub{Les niveaux de la paix}

La paix ne se décrète pas seulement au niveau des États ; elle se construit et se vit à des échelles emboîtées, de l'intime au global. Une rupture à un niveau rejaillit sur les autres : un conflit conjugal peut déstabiliser un quartier ; une guerre régionale crée des flux de réfugiés qui tendent les relations internationales.

\begin{popfigure}
\begin{tikzpicture}[scale=0.9, every node/.style={font=\small}]
    % Couleurs pour chaque niveau
    \colorlet{lvl1}{popblue!80}
    \colorlet{lvl2}{popgreen!70}
    \colorlet{lvl3}{poporange!70}
    \colorlet{lvl4}{poppurple!70}
    \colorlet{lvl5}{poppink!70}
    
    % Cercles concentriques (du plus grand au plus petit pour le remplissage)
    \fill[lvl5] (0,0) circle (5.5cm);
    \fill[lvl4] (0,0) circle (4.4cm);
    \fill[lvl3] (0,0) circle (3.3cm);
    \fill[lvl2] (0,0) circle (2.2cm);
    \fill[lvl1] (0,0) circle (1.1cm);
    
    % Contours
    \draw[popdark, thick] (0,0) circle (5.5cm);
    \draw[popdark, thick] (0,0) circle (4.4cm);
    \draw[popdark, thick] (0,0) circle (3.3cm);
    \draw[popdark, thick] (0,0) circle (2.2cm);
    \draw[popdark, thick] (0,0) circle (1.1cm);
    
    % Textes au centre (niveau 1)
    \node[popdark, align=center, font=\bfseries] at (0,0) {Paix\\intérieure};
    
    % Textes sur les anneaux (placement radial)
    % Niveau 2 : Familiale
    \node[lvl2!80!black, align=center, font=\bfseries] at (0, 1.65) {Paix familiale};
    \node[lvl2!80!black, align=center, font=\bfseries] at (0, -1.65) {et communautaire};
    
    % Niveau 3 : Locale / Régionale
    \node[lvl3!80!black, align=center, font=\bfseries] at (2.75, 1.2) {Paix locale};
    \node[lvl3!80!black, align=center, font=\bfseries] at (-2.75, 1.2) {et régionale};
    \node[lvl3!80!black, align=center, font=\bfseries] at (2.75, -1.2) {Cohésion};
    \node[lvl3!80!black, align=center, font=\bfseries] at (-2.75, -1.2) {sociale};
    
    % Niveau 4 : Nationale
    \node[lvl4!80!black, align=center, font=\bfseries] at (0, 3.85) {Paix nationale};
    \node[lvl4!80!black, align=center, font=\bfseries] at (0, -3.85) {Unité, justice, vivre ensemble};
    
    % Niveau 5 : Internationale
    \node[lvl5!80!black, align=center, font=\bfseries] at (0, 5.0) {Paix internationale};
    \node[lvl5!80!black, align=center, font=\bfseries] at (0, -5.0) {Coopération, Droit, ONU, UA};
    
    % Flèches d'interdépendance
    \draw[popdark, -{Stealth[length=3mm]}, thick] (0, 1.1) -- (0, 1.3);
    \draw[popdark, -{Stealth[length=3mm]}, thick] (0, -1.1) -- (0, -1.3);
    \draw[popdark, -{Stealth[length=3mm]}, thick] (0, 2.2) -- (0, 2.5);
    \draw[popdark, -{Stealth[length=3mm]}, thick] (0, -2.2) -- (0, -2.5);
    \draw[popdark, -{Stealth[length=3mm]}, thick] (0, 3.3) -- (0, 3.6);
    \draw[popdark, -{Stealth[length=3mm]}, thick] (0, -3.3) -- (0, -3.6);
    \draw[popdark, -{Stealth[length=3mm]}, thick] (0, 4.4) -- (0, 4.7);
    \draw[popdark, -{Stealth[length=3mm]}, thick] (0, -4.4) -- (0, -4.7);
\end{tikzpicture}
\legende{Les cinq niveaux emboîtés de la paix : de l'intériorité de la personne (centre) à l'ordre mondial (périphérie). Chaque niveau conditionne les suivants ; la stabilité globale repose sur la solidité des fondations locales et individuelles.}
\end{popfigure}

\begin{itemize}
    \item \textbf{Paix intérieure} (individu) : équilibre psychologique, estime de soi, gestion non-violente des émotions. C'est le socle : un individu en guerre contre lui-même peine à être artisan de paix.
    \item \textbf{Paix familiale et communautaire} : dialogue intergénérationnel, résolution des disputes sans violence, solidarité de voisinage, respect des coutumes pacifiques (palabres, médiation des notables).
    \item \textbf{Paix locale et régionale} : cohésion entre groupes ethniques, religieux, professionnels au sein d'une commune, d'un arrondissement, d'une région. Gestion concertée des ressources (eau, terres, marchés).
    \item \textbf{Paix nationale} : unité dans la diversité, État de droit, justice équitable, institutions légitimes, intégration des minorités, lutte contre la corruption et les inégalités territoriales. C'est le niveau où s'appliquent la Constitution et les lois.
    \item \textbf{Paix internationale} : respect de la Charte de l'ONU, de l'Acte constitutif de l'Union africaine, règlement pacifique des différends frontaliers, coopération économique, diplomatie préventive.
\end{itemize}

\courssub{Les fondements de la paix}

La paix positive ne tombe pas du ciel ; elle se construit sur des piliers concrets que toute politique publique, toute association, tout citoyen doit consolider.

\begin{propriete}[Les cinq piliers de la paix positive (référentiel ECM)]
Pour qu'une société soit en paix positive, cinq conditions sont conjointement nécessaires :
\begin{enumerate}
    \item \textbf{La justice} : égalité devant la loi, lutte contre l'impunité, accès effectif aux tribunaux, justice transitionnelle après crises.
    \item \textbf{La tolérance} : acceptation de la différence (ethnique, religieuse, d'opinion, de genre), refus de l'exclusion et du discours de haine.
    \item \textbf{Le dialogue} : culture de la concertation, espaces de parole (parlements, conseils municipaux, comités de vigilance, clubs scolaires), médiation.
    \item \textbf{Le respect des droits humains} : droits civils et politiques (liberté d'expression, vote) \textbf{et} droits économiques, sociaux, culturels (santé, éducation, travail décent).
    \item \textbf{Le développement équitable} : répartition juste des richesses, aménagement du territoire, lutte contre la pauvreté, emploi des jeunes, infrastructures de base.
\end{enumerate}
L'absence d'un seul pilier fragilise l'édifice entier.
\end{propriete}

\begin{methode}[Analyser une situation de paix ou de conflit]
Pour traiter un cas d'étude (exercice, dossier, oral), suivez cette grille en quatre temps :
\begin{enumerate}
    \item \textbf{Identifier les acteurs} : qui sont les parties prenantes ? (individus, groupes, État, ONG, entreprises, puissances étrangères). Distinguer acteurs directs et indirects (bénéficiaires, spoilers).
    \item \textbf{Analyser les causes} : causes profondes (structurelles : inégalités, marginalisation, mauvaise gouvernance) et causes immédiates (déclencheurs : élection contestée, meurtre, pénurie). Utiliser l'arbre à problèmes (racines = causes, tronc = problème central, branches = effets).
    \item \textbf{Évaluer les conséquences} : humaines (morts, déplacés, traumatismes), matérielles (écoles brûlées, routes coupées), économiques (chômage, fuite des investisseurs), sociales (méfiance, radicalisation).
    \item \textbf{Proposer des solutions} : court terme (cessez-le-feu, aide humanitaire, médiation), moyen terme (réconciliation, justice transitionnelle, reconstruction), long terme (réformes institutionnelles, développement inclusif, éducation à la paix). Hiérarchiser selon urgence et faisabilité.
\end{enumerate}
\end{methode}

\courssub{TD4 : Les symboles de la paix au Cameroun (enquête guidée)}

Au-delà des concepts, la paix s'incarne dans des symboles forts, vecteurs d'identité commune et de mémoire partagée. Ce travail d'enquête (réalisable en groupe ou en individuel) vise à les identifier et à en analyser la portée.

\begin{methode}[Démarche d'enquête sur les symboles de la paix]
\begin{enumerate}
    \item \textbf{Recenser} : Lister les symboles officiels (drapeau, hymne, devise, sceau, armoiries), les monuments (Monument de la Réunification à Yaoundé, Statue de la Liberté à Buea), les lieux de mémoire (Mausolée des Martyrs), les journées commémoratives (20 mai, 11 février, 1er mai), les distinctions honorifiques (Ordre de la Valeur, Ordre du Mérite).
    \item \textbf{Documenter} : Pour chaque symbole, noter sa description, son origine historique, sa signification explicite (paix, unité, travail, progrès) et son usage protocolaire.
    \item \textbf{Analyser la réception} : Interroger des élèves, des enseignants, des anciens combattants ou des notables : « Que vous évoque ce symbole ? Unit-il ou divise-t-il ? » Noter les convergences et les divergences (ex. : perception du 20 mai dans les régions anglophones).
    \item \textbf{Synthétiser} : Produire une fiche synthétique ou un panneau d'exposition : « Les symboles de la paix au Cameroun : entre consensus national et appropriations locales ».
\end{enumerate}
\end{methode}

\begin{exemplebox}[Principaux symboles de la paix et de l'unité nationale]
\begin{center}
\begin{tabularx}{\linewidth}{|l|X|X|}\hline
\rowcolor{popblueL}
\textbf{Symbole} & \textbf{Description / Signification} & \textbf{Valeur de paix portée} \\ \hline
Drapeau national & Vert (forêt/Sud), Rouge (unité/sang versé), Jaune (soleil/Savanes), Étoile jaune (unité) & Unité dans la diversité territoriale et humaine \\ \hline
Devise & « Paix – Travail – Patrie » (ordre constitutionnel) & La paix comme condition première du travail et de l'amour de la patrie \\ \hline
Hymne national & « Ô Cameroun, berceau de nos ancêtres... » & Héritage commun, devoir de transmission, fidélité \\ \hline
Monument de la Réunification (Yaoundé) & Deux spirales qui fusionnent (Ahmadou Ahidjo / John Ngu Foncha) & Réconciliation des deux héritages coloniaux, unité nationale \\ \hline
Journée de la Jeunesse (11 fév.) & Défilés, activités socio-éducatives, message du Chef de l'État & Investissement dans la jeunesse, relais de la paix future \\ \hline
Ordre de la Valeur & Distinction honorifique suprême & Reconnaissance du mérite au service de la nation \\ \hline
\end{tabularx}
\end{center}
\end{exemplebox}

\courssub{Les stratégies de promotion de la paix}

La promotion de la paix ne se limite pas à la gestion de crise ; elle relève d'une politique publique transversale, portée par l'État, la société civile et les partenaires au développement.

\begin{propriete}[Stratégies majeures de promotion de la paix au Cameroun]
\begin{enumerate}
    \item \textbf{Cadre institutionnel et juridique} : Constitution de 1996 (préambule, art. 1er, 65), Loi n°2019/014 du 19 déc. 2019 relative aux discours de haine et au tribalisme, Loi n°2005/007 du 27 juil. 2005 portant création de la Commission Nationale des Droits de l'Homme et des Libertés (CNDHL), Décret n°2017/013 du 23 mars 2017 créant le Comité National de Pilotage du Plan d'Urgence Humanitaire (CNPUH).
    \item \textbf{Dialogue national et concertation} : Grand Dialogue National (GDN, sept.-oct. 2019) → recommandations : statut spécial NO/SO, commission bilinguisme, décentralisation. Mise en œuvre : Loi n°2019/024 du 24 déc. 2019 (Code général des collectivités territoriales décentralisées), création des régions.
    \item \textbf{Justice transitionnelle et réconciliation} : Centres de Désarmement, Démobilisation et Réintégration (DDR) pour ex-combattants (Boko Haram, séparatistes), Commissions Vérité-Justice-Réconciliation (locales), médiation des notables et leaders religieux (Cadre de concertation des confessions religieuses).
    \item \textbf{Éducation à la paix et citoyenneté} : Intégration ECM dans les curricula (Primaire, Secondaire, Technique), Clubs UNESCO / Paix / Droits de l'Homme dans les lycées, Formation des enseignants (ENS, ENIEG), Campagnes « École sans violence ».
    \item \textbf{Développement inclusif et lutte contre la pauvreté} : SND30 (2020-2030) : Pilar « Capital humain et bien-être », Plan d'Urgence Triennal (PLANUT) pour le Nord, NO, SO, Adamaoua, Est. Projets : PAJER-U, PIAASI, insertion socio-économique des jeunes et femmes.
    \item \textbf{Sécurité humaine et gouvernance locale} : Police de proximité, Comités de vigilance encadrés, Comités locaux de sécurité (CLS), Opérations « Épervier » (sécurité routière), Réforme du secteur de la sécurité (RSS).
\end{enumerate}
\end{propriete}

\courssub{Le règlement pacifique des différends}

Le règlement pacifique est l'outil opérationnel de la paix positive. Il s'appuie sur l'article 33 de la Charte de l'ONU (négociation, enquête, médiation, conciliation, arbitrage, règlement judiciaire, recours aux organes régionaux) et sur l'article 4 de l'Acte constitutif de l'UA.

\begin{definition}[Mécanismes de règlement pacifique (graduation)]
\begin{itemize}
    \item \textbf{Négociation} : Pourparlers directs entre parties sans tiers (ex. : accords d'entreprise, conventions collectives).
    \item \textbf{Bons offices / Médiation} : Intervention d'un tiers facilitateur (personnalité, État, organisation) qui propose mais n'impose pas (ex. : médiation de la CEMAC, du Centre Carter, de l'Église catholique au Cameroun).
    \item \textbf{Conciliation} : Tiers qui formule une proposition de règlement après enquête (ex. : Commission de conciliation CIRDI, commissions administratives locales).
    \item \textbf{Arbitrage} : Les parties confient le différend à un tribunal arbitral dont la sentence s'impose (ex. : différends commerciaux, investissements, OHADA).
    \item \textbf{Règlement judiciaire} : Saisine d'une cour permanente dont l'arrêt est contraignant (CIJ pour les États, CEDH, Cour africaine des droits de l'homme, Cour de justice de la CEMAC/CEEAC).
    \item \textbf{Recours aux organes régionaux} : UA (Conseil de paix et de sécurité), CEEAC (COPAX), CEMAC.
\end{itemize}
\end{definition}

\begin{exemplebox}[Résolution pacifique du différend frontalier du Bakassi (Cameroun - Nigéria)]
Ce cas emblématique illustre le passage de la paix négative (risque de guerre) à la paix positive (droit, coopération).
\begin{itemize}
    \item \textbf{Contexte} : Péninsule pétrolifère et halieutique revendiquée par les deux pays. Incidents armés (1981, 1994 - mort de soldats camerounais).
    \item \textbf{Choix du droit} : Le Cameroun saisit la \cle{Cour internationale de Justice (CIJ)} en 1994. Le Nigéria accepte la compétence.
    \item \textbf{Arrêt de la CIJ (10 octobre 2002)} : La Cour attribue la souveraineté du Bakassi au Cameroun sur la base des traités coloniaux (traité germano-britannique de 1913) et de l'\emph{effectivité} administrative camerounaise.
    \item \textbf{Mise en œuvre} : Accords de \cle{Greentree} (juin 2006) sous égide de l'ONU (Kofi Annan). Transfert progressif, protection des populations nigérianes restées sur place, démilitarisation, commission mixte de suivi.
    \item \textbf{Enseignement} : Le règlement pacifique par le droit international (paix positive) a évité une guerre régionale, sauvé des vies et permis une coopération ultérieure (pipeline Tchad-Cameroun, lutte anti-Boko Haram).
\end{itemize}
\end{exemplebox}

\begin{exemplebox}[Médiation locale : Conflit agro-pastoral dans le Logone-et-Chari (2021-2023)]
\begin{itemize}
    \item \textbf{Mécanisme} : Médiation conduite par le Gouverneur de l'Extrême-Nord, le Sultan de Logone-Birni, le Lamido de Kousseri, avec l'appui du PNUD et de l'ONG Search for Common Ground.
    \item \textbf{Processus} : Cessez-le-feu $\rightarrow$ Enquête conjointe $\rightarrow$ Identification des couloirs de transhumance $\rightarrow$ Signature de conventions locales (pactes de non-agression, calendriers de passage) $\rightarrow$ Mise en place de comités de suivi mixtes (éleveurs/agriculteurs/administration).
    \item \textbf{Résultat} : Baisse significative des affrontements en 2023-2024, retour progressif des déplacés, réouverture des marchés hebdomadaires.
\end{itemize}
\end{exemplebox}

\courssub{Dossier 2 : La lutte contre l'extrémisme violent au Cameroun}

L'extrémisme violent (EV) désigne l'adhésion à une idéologie radicale justifiant l'usage de la violence pour imposer un projet politique, religieux ou identitaire. Au Cameroun, il prend deux visages principaux.

\begin{propriete}[Les deux foyers d'extrémisme violent au Cameroun]
\begin{itemize}
    \item \textbf{Boko Haram / ISWAP (Extrême-Nord, depuis 2014)} : Idéologie salafiste djihadiste, rejet de l'école occidentale, de la démocratie, de l'État laïc. Mode opératoire : attentats suicides (souvent femmes/enfants enlevés), embuscades, enlèvements, pillages. Zones : Mayo-Sava, Mayo-Tsanaga, Logone-et-Chari.
    \item \textbf{Séparatisme armé (Nord-Ouest / Sud-Ouest, depuis 2017)} : Revendication d'un État indépendant « Ambazonie ». Groupes armés multiples (ADF, Red Dragons, Tigers, etc.). Mode opératoire : « Ghost towns » (villes mortes), attaques ciblées (forces de l'ordre, administration, écoles), enlèvements contre rançon, destruction d'infrastructures.
\end{itemize}
\end{propriete}

\begin{aretenir}[Réponse nationale : Approche globale « Sécurité - Développement - Prévention »]
\begin{enumerate}
    \item \textbf{Volet militaire / sécuritaire} : Force Multinationale Mixte (FMM - Tchad, Niger, Nigeria, Cameroun, Bénin), Opération « Émergence 4 » (BIR, armée régulière), Comités de vigilance (CV) encadrés par la loi (Loi n°2013/001), Désarmement volontaire (centres DDR de Mora, Mémé, Buea, Bamenda).
    \item \textbf{Volet judiciaire} : Loi n°2014/028 du 23 déc. 2014 réprimant le terrorisme, Tribunaux militaires, Cellule de traitement des dossiers terroristes, Respect du droit humanitaire (formation FDS avec CICR).
    \item \textbf{Volet prévention / déradicalisation (P/CVE)} : Plan d'Action National de Prévention de l'Extrémisme Violent (PAN-PEV, 2018), Rôle des leaders religieux (Imams, Pasteurs, Chefs traditionnels) pour contrer le discours radical, Éducation à la tolérance (ECM, clubs paix), Surveillance des lieux de culte et écoles coraniques.
    \item \textbf{Volet humanitaire / développement} : Plan de Réponse Humanitaire (OCHA), Plan d'Urgence Triennal (PLANUT, 370 milliards FCFA), Reconstruction d'écoles/santé, Appui aux activités génératrices de revenus (AGR) pour jeunes à risque, Réinsertion des ex-combattants (formation professionnelle, kits d'installation).
\end{enumerate}
\end{aretenir}

\begin{attention}[Distinction cruciale : Extrémisme violent vs Extrémisme (idées)]
L'\textbf{extrémisme} (au sens large) est le fait de porter des idées radicales, hors du consensus démocratique, sans forcément passer à l'acte violent. L'\textbf{extrémisme violent} (EV) implique le passage à l'acte (terrorisme, insurrection). La loi camerounaise (2014/028) sanctionne les \emph{actes} de terrorisme et l'\emph{apologie} (discours de haine, recrutement), non la simple pensée. La prévention (P/CVE) vise à empêcher la bascule de l'extrémisme idéologique vers l'extrémisme violent.
\end{attention}

\courssub{Le vivre ensemble et le bilinguisme (Dossier 3 : La CNPBM)}

Le « vivre ensemble » n'est pas un slogan ; c'est un contrat social quotidien dans un pays de \cle{diversité culturelle} (250+ groupes, 2 langues officielles, 3 grandes zones écologiques, laïcité). La \cle{Commission Nationale pour la Promotion du Bilinguisme et du Multiculturalisme (CNPBM)} en est l'institution garante.

\begin{definition}[La CNPBM : Création, Mission, Composition]
\begin{itemize}
    \item \textbf{Base légale} : Décret n°2017/013 du 23 mars 2017 (suite recommandation Grand Dialogue National), modifié par Décret n°2019/002 du 4 janv. 2019.
    \item \textbf{Mission} : Promouvoir le bilinguisme officiel (français/anglais) et le multiculturalisme comme facteurs d'unité nationale, de cohésion sociale et de paix. Veiller au respect des droits des minorités linguistiques et culturelles.
    \item \textbf{Composition} : Président (nommé par le PR), Vice-Président, Membres (représentants administrations, société civile, universités, confessions religieuses, médias, diaspora). Secrétariat permanent à Yaoundé. Antennes régionales en cours de déploiement.
\end{itemize}
\end{definition}

\begin{propriete}[Champs d'action et réalisations de la CNPBM (2017-2024)]
\begin{enumerate}
    \item \textbf{Surveillance et évaluation du bilinguisme} : Enquêtes dans les administrations, tribunaux, médias publics, écoles, aéroports. Rapports annuels au Président de la République. Sanctions administratives recommandées pour non-respect (affichage unilingue, accueil unilingue).
    \item \textbf{Promotion du multiculturalisme} : Inventaires du patrimoine culturel immatériel, soutien aux festivals (Ngondo, Nyem-Nyem, Medumba, Mpo'o), atlas des langues nationales (plus de 270 recensées), appui à l'enseignement des langues nationales (programme ELAN, manuels).
    \item \textbf{Education et sensibilisation} : Concours nationaux de bilinguisme (écoles, universités), Journée du bilinguisme (janvier), Journée de la diversité culturelle (21 mai), Campagnes médias « Le bilinguisme, c'est mon droit », Formation des formateurs (bilinguisme fonctionnel).
    \item \textbf{Réception et traitement des plaintes} : Guichet plaintes (physique, numérique : site web, application mobile). Exemples : plainte pour absence de version anglaise d'un arrêté ministériel, plainte pour refus d'inscription en section anglophone, plainte pour discrimination culturelle.
    \item \textbf{Conseil et proposition législative} : Avis sur projets de lois (ex. Code général des CTD, Loi sur la décentralisation), Proposition de loi sur la protection des minorités, Plaidoyer pour la ratification de la Charte africaine des langues.
\end{enumerate}
\end{propriete}

\begin{exemplebox}[Cas pratique : Plainte pour non-bilinguisme dans un hôpital de district (2022)]
\begin{itemize}
    \item \textbf{Faits} : Un patient anglophone se plaint de ne trouver aucune signalétique, aucun formulaire, aucun personnel capable de le servir en anglais dans un hôpital de district de la région du Centre.
    \item \textbf{Saisine} : Dépôt plainte via l'application mobile CNPBM / Courrier au Secrétariat permanent.
    \item \textbf{Instruction} : Enquête sur place par antenne régionale. Constat : affichage 100\% francophone, personnel non bilingue, formulaires unilingues.
    \item \textbf{Résolution} : Recommandation au MINSANTÉ et au Directeur de l'hôpital : traduction affichage/formulaires (délai 3 mois), affectation personnel bilingue / formation continue, mise en place « Guichet bilingue ».
    \item \textbf{Suivi} : Rapport de mise en œuvre transmis au Président. Sanction administrative (avertissement) au Directeur pour négligence.
\end{itemize}
\end{exemplebox}

\begin{methode}[Construire le vivre ensemble au quotidien (niveau local / établissement)]
\begin{enumerate}
    \item \textbf{Connaître l'autre} : Organiser des « journées culturelles croisées » (gastronomie, tenue, danse, contes) entre élèves de différentes origines.
    \item \textbf{Pratiquer le bilinguisme fonctionnel} : Alterner langues d'affichage, de communication orale (assemblées, clubs), de correspondance administrative dans l'établissement.
    \item \textbf{Institutionaliser le dialogue} : Mettre en place un « Conseil de la vie lycéenne » paritaire (francophones/anglophones, filles/garçons, toutes spécialités) avec pouvoir de proposition.
    \item \textbf{Sanctionner la haine} : Appliquer le règlement intérieur (inspiré Loi 2019/014) : toute insulte ethnique, religieuse, linguistique = conseil de discipline + médiation réparatrice.
    \item \textbf{S'engager solidaire} : Projets communs inter-classes / inter-spécialités (reboisement, nettoyage quartier, collecte pour déplacés, tutorat pair-à-pair).
\end{enumerate}
\end{methode}

\courssub{La paix, condition du développement}

Il n'y a pas de développement durable sans paix, pas de paix durable sans développement. Ce cercle vertueux (ou vicieux) est au cœur de l'Agenda 2030 de l'ONU (ODD 16 : Paix, Justice, Institutions efficaces) et de la Stratégie Nationale de Développement 2020-2030 (SND30) du Cameroun.

\begin{itemize}
    \item \textbf{Stabilité et investissement} : les investisseurs (publics et privés) fuient l'incertitude. Une route non entretenue par insécurité isole une zone agricole ; une usine fermée par des grèves violentes détruit des emplois. La paix permet la planification à long terme.
    \item \textbf{Éducation et capital humain} : écoles fermées (stratégie « ghost towns » au NO/SO), enseignants en fuite, enfants recrutés par groupes armés = génération sacrifiée. La paix rouvre les classes, forme les compétences.
    \item \textbf{Cohésion sociale et capital de confiance} : la méfiance interethnique augmente les coûts de transaction (contrats, garanties, médiation). La paix reconstruite (dialogue, justice transitionnelle) restaure la confiance, lubrifiant de l'économie.
    \item \textbf{État de droit et redistribution} : un État fort et juste perçoit l'impôt, finance la santé, les routes, l'électricité. La corruption et la prédation (symptômes de l'absence de paix positive) détournent les ressources du développement.
\end{itemize}

\begin{exoresolu}[Analyse d'un conflit agro-pastoral dans le Mayo-Danay]
\textbf{Énoncé} : Dans le département du Mayo-Danay (Extrême-Nord), des affrontements opposent éleveurs peuls et agriculteurs mousgoum en décembre 2023. Bilan : 12 morts, 3 villages brûlés, 2 000 déplacés. Le sous-préfet convoque une réunion de réconciliation. En utilisant la méthode d'analyse, identifiez acteurs, causes, conséquences et proposez deux solutions (court et long terme).

\textbf{Solution détaillée} :
\begin{enumerate}
    \item \textbf{Acteurs} :
    \begin{itemize}
        \item Directs : Éleveurs (transhumants ou sédentaires), Agriculteurs (autochtones), Autorités administratives (sous-préfet, préfet), Forces de l'ordre, Notables traditionnels (Lamibe, Chefs de canton), ONG humanitaires (CICR, Croix-Rouge).
        \item Indirects : Marchands de bétail, Élus locaux (maires, députés), Groupes d'autodéfense, Services de l'élevage et de l'agriculture.
    \end{itemize}
    \item \textbf{Causes} :
    \begin{itemize}
        \item \textit{Profondes} : Pression démographique (+3\%/an), changement climatique (sécheresse, raréfaction eau/pâturages), absence de couloirs de transhumance balisés, loi foncière inadaptée (droit coutumier vs droit moderne), impunité antérieure, prolifération armes légères.
        \item \textit{Immédiates} : Destruction de champs par troupeaux (début saison sèche), vol de bétail, rumeur de meurtre d'un berger, absence de médiation précoce.
    \end{itemize}
    \item \textbf{Conséquences} :
    \begin{itemize}
        \item Humaines : 12 morts, blessés, traumatismes, 2 000 déplacés (femmes/enfants majoritaires).
        \item Matérielles : 3 villages brûlés, greniers détruits, bétail volé/abattu.
        \item Économiques : Marché hebdomadaire fermé, perte de revenus, hausse prix céréales/viande.
        \item Sociales : Haine interethnique, discours de haine sur WhatsApp, risque de radicalisation jeunes.
    \end{itemize}
    \item \textbf{Solutions} :
    \begin{itemize}
        \item \textit{Court terme (urgence)} : Cessez-le-feu imposé par forces de l'ordre, secours humanitaires (abris, vivres, soins), médiation immédiate par notables/autorités, arrestation des auteurs de violences (justice), sécurisation retour déplacés.
        \item \textit{Long terme (structurel)} : Délimitation et balisage couloirs de transhumance + aires de pâturage, conventions locales écrites (éleveurs/agriculteurs/autorités), commissions de conciliation permanentes (chefs traditionnels, administration), investissements eau (forages, barrages), éducation à la paix à l'école, désarmement volontaire, développement économique local (emploi jeunes).
    \end{itemize}
\end{enumerate}
\end{exoresolu}

\begin{aretenir}[Synthèse : La notion de paix]
\begin{itemize}
    \item La paix n'est pas l'absence de guerre ; c'est un \textbf{processus dynamique} combinant \textbf{paix négative} (absence violence directe) et \textbf{paix positive} (justice, équité, droits humains, développement) — distinction de \textbf{Johan Galtung}.
    \item Elle s'incarne à \textbf{cinq niveaux emboîtés} : intérieure $\rightarrow$ familiale/communautaire $\rightarrow$ locale/régionale $\rightarrow$ nationale $\rightarrow$ internationale. La rupture à un niveau contamine les autres.
    \item Ses \textbf{fondements} sont la justice, la tolérance, le dialogue, les droits humains, le développement équitable.
    \item Au Cameroun, les \textbf{menaces} sont multiples : crise anglophone (NO/SO), insurrection Boko Haram (Extrême-Nord), conflits agro-pastoraux (Nord), criminalité urbaine, inégalités, discours de haine.
    \item La paix est la \textbf{condition sine qua non du développement} : elle attire l'investissement, permet l'éducation, restaure la confiance, légitime l'État.
    \item Le \textbf{règlement pacifique des différends} (négociation, médiation, arbitrage, justice ; ex. Bakassi : CIJ + Greentree) est le modèle à suivre : droit, dialogue, mise en œuvre concertée.
    \item Les \textbf{symboles} (drapeau, devise, hymne, monuments, distinctions) et les \textbf{institutions} (CNPBM, CNDHL, CDR) ancrent la paix dans le concret du quotidien.
    \item La \textbf{lutte contre l'extrémisme violent} exige une approche globale (sécurité, justice, prévention P/CVE, développement, réinsertion DDR) respectueuse des droits humains.
    \item Le \textbf{vivre ensemble} se construit par le bilinguisme effectif, la valorisation du multiculturalisme et l'engagement citoyen de chacun (école, quartier, entreprise, réseaux sociaux).
\end{itemize}
\end{aretenir}

\begin{exercice}[Entraînement Probatoire - Sujet type]
\begin{enumerate}
    \item Définir la paix positive en citant son théoricien. La distinguer de la paix négative. (4 pts)
    \item Citer trois piliers de la paix positive et illustrer chacun par une action concrète au Cameroun. (6 pts)
    \item Présenter les deux principaux foyers d'extrémisme violent au Cameroun et les quatre volets de la réponse nationale. (6 pts)
    \item Expliquer le rôle de la CNPBM dans la consolidation du vivre ensemble. Donner deux exemples de ses réalisations. (4 pts)
\end{enumerate}
\end{exercice}

\begin{corrige}[Exercice n°1 - Éléments de réponse]
\begin{enumerate}
    \item \textbf{Paix positive (J. Galtung)} : Présence de structures justes, équitables, participatives (justice sociale, droits humains, développement). \textbf{Paix négative} : Simple absence de violence directe/guerre. La première prévient les causes, la seconde constate l'arrêt des effets.
    \item \textbf{Justice} : Création tribunaux de proximité, aide juridictionnelle. \textbf{Tolérance} : Loi 2019/014 contre discours de haine, Journée diversité culturelle. \textbf{Développement équitable} : PLANUT (zones en crise), PAJER-U (emploi jeunes).
    \item \textbf{Foyers} : (1) Boko Haram/ISWAP (Extrême-Nord, djihadisme). (2) Séparatisme armé (NO/SO, identité). \textbf{Volets réponse} : (1) Militaire (FMM, BIR, CV). (2) Judiciaire (Loi antiterrorisme, tribunaux militaires). (3) Prévention/P-CVE (PAN-PEV, leaders religieux, éducation). (4) Humanitaire/Développement (PLANUT, DDR, réinsertion).
    \item \textbf{Rôle CNPBM} : Promouvoir bilinguisme/multiculturalisme comme ciment national. \textbf{Réalisations} : (1) Enquêtes bilinguisme administrations + rapports au PR + recommandations sanctions. (2) Traitement plaintes citoyens (non-respect bilinguisme services publics) + résolutions concrètes (traduction, formation). (3) Concours bilinguisme écoles, Atlas langues, Journées officielles.
\end{enumerate}
\end{corrige}

\coursec{TD4 : Les symboles de la paix au Cameroun (enquêtes)}

Dans la section précédente, nous avons découvert que la paix n'est pas seulement l'absence de guerre : c'est un état de concorde, de sécurité et de justice qui permet à chacun de vivre dignement. Mais comment un peuple \emph{rend-il visible} cette paix ? Comment la rappelle-t-il chaque jour à ses citoyens ? C'est là qu'interviennent les \cle{symboles de la paix} : des signes que tout le monde peut voir, toucher, chanter ou visiter. Ce travail dirigé (TD) nous apprendra à les reconnaître autour de nous et à mener une véritable \cle{enquête} dans notre localité, comme le ferait un enquêteur professionnel.

\courssub{Qu'est-ce qu'un symbole de paix ?}

Commençons par une situation simple. Dans ton quartier ou ton village, lorsque deux familles se réconcilient après un conflit, on organise souvent une cérémonie : on partage la kola, on se serre la main devant les anciens, on plante parfois un arbre. Ces gestes et ces objets ne sont pas la paix elle-même : ils la \emph{représentent} et la \emph{rappellent}. C'est exactement le rôle d'un symbole.

\begin{definition}[Symbole de paix]
Un \cle{symbole de paix} est un signe matériel (objet, monument, lieu, drapeau) ou immatériel (chant, devise, geste, personnage) qui représente la \cle{concorde} et l'\cle{unité} d'un peuple, et qui rappelle aux citoyens l'importance de vivre ensemble sans violence.
\end{definition}

Un symbole de paix possède trois caractéristiques essentielles :
\begin{itemize}
\item il est \textbf{visible ou perceptible} par tous (on le voit, on l'entend, on le visite) ;
\item il a une \textbf{signification partagée} : la communauté reconnaît ce qu'il représente ;
\item il a une \textbf{fonction de rappel} : il invite chacun à préserver la paix et à rejeter la violence.
\end{itemize}

\begin{attention}[Symbole de paix ou symbole de souveraineté ?]
Ne confonds pas \cle{symbole de paix} et \cle{symbole de souveraineté}. Un symbole de souveraineté (sceau de l'État, portrait du Président de la République, passeport) affirme l'autorité et l'indépendance de l'État. Un symbole de paix rappelle la concorde et l'unité. \textbf{Attention au piège} : certains symboles comme le drapeau ou l'hymne national \emph{cumulent les deux fonctions}. Dans une copie d'examen, précise toujours laquelle des deux fonctions est mise en avant dans la question posée.
\end{attention}

\courssub{Les symboles nationaux de la paix au Cameroun}

\textbf{Le drapeau national.} Le drapeau du Cameroun comporte trois bandes verticales : le vert, le rouge et le jaune, avec une \cle{étoile} au centre de la bande rouge. Cette étoile est appelée \emph{l'étoile de l'unité} : elle symbolise l'unité indivisible du peuple camerounais, c'est-à-dire la volonté de vivre ensemble malgré la diversité des régions, des langues et des cultures. Le vert évoque la forêt du Sud et l'espérance, le rouge l'unité, le jaune la savane du Nord et le soleil.

\textbf{L'hymne national.} « Ô Cameroun, berceau de nos ancêtres » est chanté lors des cérémonies officielles, dans les écoles et les lycées. Ses paroles appellent à l'unité, au travail et à la paix. Chanter ensemble le même hymne, c'est déjà un acte de vivre ensemble.

\textbf{La devise.} La devise officielle du Cameroun est « \cle{Paix-Travail-Patrie} ».

\begin{exemplebox}[La devise « Paix-Travail-Patrie »]
La devise du Cameroun est inscrite sur les \cle{armoiries} de la République. Le fait que le mot \emph{Paix} vienne en premier n'est pas un hasard : il rappelle que la paix est le \textbf{premier fondement de la nation}. Sans paix, pas de travail possible ; sans travail, pas de patrie prospère. Ainsi, chaque fois qu'un élève voit les armoiries sur un document officiel (bulletin, acte administratif, billet de banque), il reçoit un message de paix.
\end{exemplebox}

\textbf{Les armoiries.} Elles représentent un bouclier portant la carte du Cameroun, une épée et une balance (symboles de la justice), encadrées de deux faisceaux, avec la devise. La justice représentée sur les armoiries est une condition de la paix : là où règne l'injustice, la paix est menacée.

\courssub{Les monuments et lieux de paix}

\textbf{Le Monument de la Réunification (Yaoundé).} Érigé au cœur de la capitale, ce monument est l'un des plus importants symboles de paix du pays. Il se compose d'une structure en spirale (forme de coquille) surmontée d'une statue : un \cle{vieillard} tenant cinq \cle{enfants}, qui tendent vers lui une étoile. Le vieillard symbolise la sagesse et la tradition ; les enfants représentent les générations futures ; la spirale évoque la montée progressive vers l'unité.

\begin{popfigure}
\begin{center}
\begin{tikzpicture}[scale=1.0]
% spirale simplifiée du monument
\draw[fill=popblueL,draw=popblue,thick] (0,0) .. controls (1.6,0.2) and (1.8,1.4) .. (0.9,2.1) .. controls (0.2,2.7) and (-0.9,2.7) .. (-1.2,1.8) .. controls (-1.5,0.9) and (-0.9,0.1) .. (0,0);
\draw[fill=popblue!45,draw=popblue,thick] (0,0.15) .. controls (1.1,0.35) and (1.25,1.25) .. (0.65,1.75) .. controls (0.15,2.15) and (-0.6,2.1) .. (-0.8,1.5) .. controls (-1.0,0.9) and (-0.6,0.25) .. (0,0.15);
% colonne centrale
\draw[fill=popgoldL,draw=popgold,thick] (-0.12,0) rectangle (0.12,3.1);
% vieillard (tête + corps)
\draw[fill=poporangeL,draw=poporange,thick] (0,3.55) circle (0.22);
\draw[fill=poporangeL,draw=poporange,thick] (-0.28,3.3) -- (0.28,3.3) -- (0.2,2.55) -- (-0.2,2.55) -- cycle;
% bras du vieillard
\draw[poporange,very thick] (-0.25,3.15) -- (-0.95,3.45);
\draw[poporange,very thick] (0.25,3.15) -- (0.95,3.45);
% enfants
\foreach \x in {-1.15,-0.75,0.75,1.15}{
  \draw[fill=popgreenL,draw=popgreen] (\x,3.75) circle (0.13);
  \draw[popgreen,thick] (\x,3.62) -- (\x,3.35);
}
% étoile
\node[star,star points=5,fill=poppink,draw=poppink,minimum size=0.5cm] at (0,4.15) {};
% socle
\draw[fill=popmuted!40,draw=popmuted,thick] (-1.7,-0.35) rectangle (1.7,0);
\end{tikzpicture}
\end{center}
\legende{Schéma simplifié du Monument de la Réunification à Yaoundé : 1. structure en spirale (montée vers l'unité) ; 2. le vieillard (sagesse, tradition) ; 3. les enfants (avenir du pays) ; 4. l'étoile de l'unité.}
\end{popfigure}

\begin{savaistu}[Le Monument de la Réunification]
Le Monument de la Réunification à Yaoundé symbolise la \cle{réunification} des deux Camerouns : le Cameroun sous administration française (indépendant le 1\textsuperscript{er} janvier 1960) et le Cameroun méridional sous administration britannique, qui a choisi par le référendum du 11 février 1961 de rejoindre la République du Cameroun. Le 1\textsuperscript{er} octobre 1961, les deux entités forment la République fédérale du Cameroun. Ce monument rappelle donc que l'unité du pays est le fruit d'un \emph{choix pacifique} du peuple, et non de la guerre.
\end{savaistu}

\textbf{Les autres monuments et les places publiques.} D'autres monuments (stèles commémoratives, monuments de l'unité dans les chefs-lieux) et des lieux comme le Boulevard du 20 Mai à Yaoundé ou les places des fêtes dans les chefs-lieux de régions servent de cadres aux célébrations nationales. Ces lieux rassemblent les citoyens de toutes les origines : ils incarnent concrètement le vivre ensemble.

\courssub{Les personnalités, institutions et symboles internationaux}

\textbf{Les personnalités et institutions incarnant la paix.} Dans nos localités, la paix est incarnée par des personnes concrètes :
\begin{itemize}
\item les \cle{chefs traditionnels} (sultans, lamibe, fons, chefs de village), qui règlent les conflits par la palabre et la médiation ;
\item les \cle{autorités religieuses} (pasteurs, imams, prêtres), qui prêchent la réconciliation et le pardon ;
\item les \cle{médiateurs} et les lauréats de prix de la paix, dont l'exemple le plus célèbre est le prix Nobel de la paix, décerné chaque année à des personnes ou organisations ayant œuvré pour la paix dans le monde.
\end{itemize}

\textbf{Les symboles internationaux présents au Cameroun.} Certains symboles de paix sont universels et se retrouvent aussi au Cameroun :
\begin{itemize}
\item la \cle{colombe} blanche, souvent dessinée avec un rameau dans le bec, symbole universel de la paix ;
\item le \cle{rameau d'olivier}, hérité des traditions anciennes, signe de réconciliation ;
\item le \cle{drapeau blanc}, signe de trêve et de demande de paix pendant un conflit ;
\item la \cle{flamme de la paix}, allumée lors de certaines cérémonies de réconciliation.
\end{itemize}

\begin{aretenir}[L'essentiel sur les symboles de la paix]
Un symbole de paix est un signe matériel ou immatériel qui incarne la concorde et l'unité. Au Cameroun, les principaux symboles sont : le drapeau (étoile de l'unité), l'hymne national, la devise « Paix-Travail-Patrie », les armoiries, le Monument de la Réunification, les chefs traditionnels et religieux, ainsi que les symboles universels (colombe, rameau d'olivier, drapeau blanc).
\end{aretenir}

\courssub{La démarche d'enquête : comment enquêter sur les symboles de la paix ?}

Savoir reconnaître un symbole dans un manuel, c'est bien ; aller vérifier sur le terrain ce que les populations en pensent, c'est mieux. C'est le but de ce TD : te transformer en enquêteur. Une enquête est une \cle{démarche d'investigation} qui suit des étapes précises, comme en sciences expérimentales.

\begin{methode}[Mener une enquête sur les symboles de la paix]
\begin{enumerate}
\item \textbf{Définir l'objectif} : formuler clairement ce qu'on cherche. Exemple : « Quels sont les symboles de paix connus et cités par les habitants de ma localité ? »
\item \textbf{Élaborer les outils} : rédiger un \cle{questionnaire} (questions fermées : « Citez un symbole de paix » ; questions ouvertes : « Que représente pour vous ce monument ? »), préparer une grille d'\cle{observation} des monuments et un guide d'\cle{interview}.
\item \textbf{Choisir l'échantillon} : sélectionner les personnes à interroger (élèves, parents, commerçants, notables, autorités religieuses) en nombre suffisant et varié.
\item \textbf{Collecter les données} : distribuer le questionnaire, mener les interviews, observer et photographier les monuments et lieux de la localité.
\item \textbf{Dépouiller les données} : compter les réponses, classer les symboles cités, calculer des pourcentages.
\item \textbf{Interpréter et conclure} : analyser les résultats, les comparer, et rédiger une conclusion argumentée.
\end{enumerate}
\end{methode}

\begin{experience}[Enquête de terrain dans la localité]
\textbf{Matériel} : cahier, stylo, questionnaire imprimé, téléphone ou appareil photo (avec autorisation), grille d'observation.

\textbf{Protocole} : par groupes de quatre élèves, interrogez dix personnes de la localité (en respectant la politesse et en demandant l'accord avant toute photo). Notez chaque symbole cité. Visitez ensuite un monument ou une place publique et décrivez-le sur la grille d'observation (forme, inscriptions, état, fréquentation).

\textbf{Observations attendues} : la plupart des personnes citent le drapeau, la devise ou un monument ; certains notables citent des symboles locaux (arbre à palabres, case du chef).

\textbf{Interprétation} : les symboles nationaux sont les plus connus, mais les symboles locaux jouent un rôle important dans la paix de proximité. L'enquête montre que la paix se vit à deux niveaux : national et local.
\end{experience}

\courssub{Restitution et analyse des résultats}

Une fois les données collectées, il faut les présenter clairement. Le diagramme en barres est un outil efficace pour comparer les réponses. Voici l'exemple d'une enquête fictive menée auprès de 100 personnes sur la question : « Quel est, selon vous, le principal symbole de la paix au Cameroun ? »

\begin{popfigure}
\begin{center}
\begin{tikzpicture}
\begin{axis}[
    ybar, bar width=0.55cm,
    width=0.95\linewidth, height=6.2cm,
    ymin=0, ymax=35,
    ylabel={Pourcentage de citations (\%)},
    symbolic x coords={Drapeau, Monument, Devise, Colombe, Autres},
    xtick=data,
    x tick label style={font=\small},
    nodes near coords, nodes near coords align={vertical},
    every node near coord/.append style={font=\small},
    ymajorgrids=true, grid style={popline!40},
    axis line style={popdark},
    ]
\addplot[fill=popblue,draw=popblue!70!black] coordinates {(Drapeau,30) (Monument,25) (Devise,20) (Colombe,15) (Autres,10)};
\end{axis}
\end{tikzpicture}
\end{center}
\legende{Résultats fictifs de l'enquête « symbole de paix le plus cité » auprès de 100 personnes : drapeau 30 \%, monument 25 \%, devise 20 \%, colombe 15 \%, autres 10 \%.}
\end{popfigure}

\textbf{Analyse des résultats.} On observe que :
\begin{itemize}
\item le \textbf{drapeau} arrive en tête (30 \%) : c'est le symbole le plus visible, présent dans les écoles, les administrations et les cérémonies ;
\item le \textbf{monument} (25 \%) et la \textbf{devise} (20 \%) suivent : les symboles nationaux dominent largement (75 \% au total) ;
\item la \textbf{colombe} (15 \%) montre que les symboles internationaux sont aussi connus ;
\item les \textbf{autres} réponses (10 \%) regroupent les symboles locaux : arbre à palabres, chefs traditionnels, lieux de culte.
\end{itemize}

\textbf{Conclusion argumentée.} L'enquête révèle que les personnes interrogées identifient d'abord la paix à travers les symboles nationaux, ce qui traduit un fort attachement à l'unité du pays. Cependant, la présence de symboles locaux montre que la paix se construit aussi au quotidien, dans les villages et les quartiers, grâce aux institutions traditionnelles.

\begin{exoresolu}[Exploiter les résultats d'une enquête]
Énoncé. Un groupe d'élèves de Première technique a interrogé 80 personnes de son quartier. Résultats : 24 personnes citent le drapeau, 16 citent le Monument de la Réunification, 12 citent la devise, 8 citent la colombe et 20 citent d'autres symboles. 1) Calcule le pourcentage de chaque catégorie. 2) Rédige une conclusion de deux phrases.
\tcblower
Solution détaillée.
1) On divise chaque effectif par le total (80) et on multiplie par 100 :
\begin{align*}
\text{Drapeau} &: \frac{24}{80}\times 100 = 30\,\% \\
\text{Monument} &: \frac{16}{80}\times 100 = 20\,\% \\
\text{Devise} &: \frac{12}{80}\times 100 = 15\,\% \\
\text{Colombe} &: \frac{8}{80}\times 100 = 10\,\% \\
\text{Autres} &: \frac{20}{80}\times 100 = 25\,\%
\end{align*}
Vérification : $30 + 20 + 15 + 10 + 25 = 100\,\%$. Le compte est bon.

2) Conclusion : « Le drapeau est le symbole de paix le plus cité (30 \%), ce qui montre que les symboles nationaux marquent fortement les esprits. Toutefois, un quart des personnes interrogées citent d'autres symboles, souvent locaux, ce qui prouve que la paix se vit aussi au niveau de la communauté de proximité. »
\end{exoresolu}

\begin{exercice}[À toi d'enquêter]
1) Rédige un questionnaire de cinq questions pour enquêter sur les symboles de la paix dans ta localité. 2) Interroge dix personnes et présente tes résultats dans un tableau. 3) Rédige une conclusion argumentée de cinq lignes en précisant le symbole le plus cité et ton interprétation.
\end{exercice}

\begin{corrige}[Exercice : À toi d'enquêter]
Éléments de correction. 1) Le questionnaire doit comporter des questions claires et variées, par exemple : « Connaissez-vous un symbole de la paix au Cameroun ? Lequel ? », « Avez-vous déjà visité un monument de paix ? », « Que signifie pour vous la devise Paix-Travail-Patrie ? », « Qui incarne la paix dans votre quartier ? », « Que peut-on faire pour mieux faire connaître ces symboles ? ». 2) Le tableau doit présenter les symboles cités, les effectifs et les pourcentages (effectif divisé par 10, multiplié par 100). 3) La conclusion doit citer le symbole dominant, proposer une explication (visibilité, éducation, cérémonies) et souligner le rôle des symboles locaux. Toute réponse cohérente, avec des calculs de pourcentages exacts et une argumentation claire, est acceptée.
\end{corrige}

Ainsi, ce TD nous a appris à identifier les symboles de la paix et à mener une enquête rigoureuse pour comprendre ce qu'ils représentent pour les populations. Mais reconnaître les symboles ne suffit pas : il faut aussi agir. C'est pourquoi la section suivante étudiera les \emph{stratégies de promotion de la paix}, c'est-à-dire les actions concrètes menées pour construire et entretenir la concorde au Cameroun.

\coursec{Les stratégies de promotion de la paix}

Après avoir identifié les \cle{symboles de la paix} qui unissent les Camerounais autour d'un imaginaire commun (drapeau, hymne, sceau, devise), nous devons maintenant comprendre \emph{comment} la paix se construit concrètement au quotidien. La paix n'est pas un état statique que l'on décrète une fois pour toutes ; c'est un \cle{processus dynamique} qui exige des stratégies délibérées, portées par l'État, la société civile et les partenaires internationaux. Cette section analyse les leviers majeurs de la promotion de la paix au Cameroun.

\courssub{L'éducation à la paix et à la citoyenneté}

L'école est le premier laboratoire de la paix. L'\cle{Éducation à la Citoyenneté et à la Morale (ECM)}, enseignée de la primaire au secondaire, vise à former des citoyens conscients de leurs droits, de leurs devoirs et des valeurs du \cle{vivre ensemble} (tolérance, dialogue, respect de la différence).

\begin{exemplebox}[Clubs de paix et campagnes scolaires]
Dans de nombreux lycées techniques et établissements publics, des \cle{clubs de paix} organisent des journées thématiques : « Journée du vivre ensemble », concours d'éloquence sur la tolérance, fresques murales illustrant la diversité culturelle. Ces initiatives transforment les apprenants en acteurs de prévention des violences en milieu scolaire (harcèlement, tribalisme, radicalisation).
\end{exemplebox}

Au-delà du curriculum formel, l'éducation non formelle (associations de jeunesse, mouvements de scoutisme, organisations de femmes) renforce l'apprentissage de la résolution non violente des conflits. La \cle{pédagogie par l'exemple} reste primordiale : un enseignant qui gère une dispute entre élèves par la médiation plutôt que par la sanction pure enseigne la paix en acte.

\courssub{Le dialogue et la concertation}

Lorsque des tensions surgissent, le \cle{dialogue politique} et la \cle{concertation} sont les alternatives à l'affrontement. Le Cameroun a une tradition de concertations nationales : les lois sur les libertés de 1990 (loi n°90/052 du 19 décembre 1990) issues du grand débat national, et le \cle{Grand Dialogue National} de 2019.

\begin{exemplebox}[Le Grand Dialogue National (30 septembre -- 4 octobre 2019)]
Convoqué par le Président de la République face à la crise dans les régions du Nord-Ouest et du Sud-Ouest, ce dialogue a réuni les \cle{forces vives de la nation} (partis politiques, société civile, chefs traditionnels, religieux, diaspora, jeunes). Il a abouti à des recommandations majeures : statut spécial pour les deux régions anglophones, création de la Chambre des Chefs Traditionnels, relance de la décentralisation, amnistie pour certains actes. Cet exemple illustre la méthode \cle{inclusive} : associer toutes les parties prenantes pour légitimer les solutions.
\end{exemplebox}

La concertation s'opère aussi à l'échelle locale : conseils municipaux, comités de développement villageois, cadres de concertation intercommunautaires. Elle repose sur la \cle{recherche du consensus} plutôt que sur le rapport de force.

\courssub{La justice et l'équité}

La paix durable ne saurait exister sans \cle{justice}. L'impunité nourrit la frustration et la vengeance. Au Cameroun, plusieurs chantiers visent à renforcer l'État de droit :
\begin{itemize}
    \item \cle{Lutte contre la corruption} : Opération \cle{Épervier} (depuis 2006), Commission Nationale Anti-Corruption (CONAC), Tribunal Criminel Spécial (TCS). Sanctionner la prédation des deniers publics restaure la confiance citoyenne.
    \item \cle{Accès à la justice} : Maisons de la Justice, assistance judiciaire gratuite pour les justiciables démunis, juridictions de proximité (tribunaux d'arrondissement).
    \item \cle{Protection des droits humains} : Commission Nationale des Droits de l'Homme et des Libertés (CNDHL), ratification des instruments internationaux (Charte africaine des droits de l'homme, PIDCP).
\end{itemize}

\begin{attention}[Piège à éviter]
Réduire la justice à la seule répression pénale. La \cle{justice transitionnelle} (vérité, réparation, réconciliation) et la \cle{justice restaurative} (médiation pénale, composition pénale) sont tout aussi essentielles pour cicatriser les blessures sociales sans systématiquement incarcérer.
\end{attention}

\courssub{Le développement équitable}

Les inégalités territoriales et sociales sont des poudrières. La \cle{promotion de la paix} passe par une répartition plus juste des richesses et des infrastructures.

\begin{propriete}[Décentralisation et développement local]
La loi n°2019/024 du 24 décembre 2019 (Code général des collectivités territoriales décentralisées) renforce l'autonomie financière et administrative des communes et des régions. L'objectif : que les populations gèrent elles-mêmes leurs priorités (eau, santé, routes, marchés), réduisant le sentiment d'abandon qui alimente les revendications violentes.
\end{propriete}

Le \cle{Plan d'Urgence Triennal (PUT 2015-2017)} et le \cle{Plan National de Développement (PND 2020-2030)} ciblent spécifiquement les zones fragilisées (Extrême-Nord, Nord-Ouest, Sud-Ouest, Est) via le \cle{Plan de Reconstruction et de Développement (PRD)} pour rattraper le retard infrastructurel.

\courssub{Le rôle des médias et des réseaux sociaux}

Les médias sont des \cle{amplificateurs} : ils peuvent attiser la haine (discours de haine, \cle{fake news}, manipulation ethnique) ou construire la paix (journalisme de solutions, émissions de dialogue, vérification des faits).

\begin{savaistu}[Le Conseil National de la Communication (CNC)]
Le Cameroun abrite le siège du \cle{Conseil National de la Communication}, organe de régulation des médias. Il veille au respect de l'éthique et de la déontologie journalistiques. Il sanctionne les organes de presse diffusant des discours de haine ou des informations mensongères de nature à troubler l'ordre public. C'est un garde-fou institutionnel contre la médiatisation de la violence.
\end{savaistu}

Les \cle{réseaux sociaux} (Facebook, WhatsApp, X/Twitter, TikTok) sont devenus le principal vecteur de rumeurs et d'appels à la violence. La stratégie de paix inclut donc l'\cle{éducation aux médias et à l'information (EMI)} : apprendre aux jeunes à vérifier une source, à signaler un contenu haineux, à ne pas partager une information non vérifiée.

\courssub{Le rôle des acteurs non étatiques}

L'État ne peut pas tout. La \cle{société civile} est le relais indispensable :
\begin{itemize}
    \item \cle{ONG nationales et internationales} : Appui à la médiation communautaire, assistance humanitaire, plaidoyer pour les droits des déplacés internes.
    \item \cle{Confessions religieuses} : Église catholique, protestante, musulmane, etc. Leur autorité morale permet d'ouvrir des couloirs humanitaires, de médier entre groupes armés et forces de l'ordre, de prêcher le pardon.
    \item \cle{Chefs traditionnels} : Gardiens des coutumes, ils règlent une grande partie (estimation courante \textasciitilde 80\%) des conflits fonciers et coutumiers à la base (justice coutumière), désengorgeant les tribunaux et maintenant la cohésion villageoise.
    \item \cle{Organisations de jeunesse et de femmes} : Elles portent des initiatives de \cle{cohésion sociale} (mariages intercommunautaires, marchés de la paix, nettoyage commun des espaces publics).
\end{itemize}

\courssub{La coopération internationale}

Le Cameroun inscrit sa paix dans un cadre multilatéral :
\begin{itemize}
    \item \cle{ONU} : Appui aux élections, maintien de la paix (contingents camerounais en Centrafrique, Mali), agences humanitaires (HCR, PAM, UNICEF) pour les réfugiés et déplacés.
    \item \cle{Union Africaine (UA)} : Mécanisme de prévention des conflits, Charte africaine de la démocratie, des élections et de la gouvernance.
    \item \cle{CEEAC} (Communauté Économique des États de l'Afrique Centrale) : Force en attente d'Afrique centrale (FOMAC), libre circulation des personnes et des biens (facteur d'intégration pacifique).
    \item \cle{Accords de partenariat} : Coopération avec l'Union Européenne (Accord de Samoa 2023, successeur de Cotonou), coopération bilatérale (France, Chine, USA, Allemagne) incluant des volets « gouvernance et paix ».
\end{itemize}

\begin{center}
\begin{popfigure}
\begin{tikzpicture}[
    node distance=1.2cm and 1.5cm,
    actor/.style={rectangle, rounded corners, draw=popblue, thick, fill=popblueL, minimum width=3.5cm, minimum height=1.2cm, align=center, font=\small},
    action/.style={rectangle, rounded corners, draw=popgreen, thick, fill=popgreenL, minimum width=4cm, minimum height=1cm, align=center, font=\small\itshape},
    arrow/.style={-Stealth, thick, popdark},
    group/.style={rectangle, draw=poporange, thick, dashed, inner sep=5pt, fit=##1, label={[poporange,yshift=-2pt]above:##2}}
]
% Niveaux
\node[actor, align=center] (etat){ÉTAT\\(Gouvernement, Parlement,\\Justice, Forces de défense,\\Collectivités territoriales)};
\node[actor, below left=of etat, align=center] (civile){SOCIÉTÉ CIVILE\\(ONG, Confessions religieuses,\\Chefs traditionnels,\\Jeunes, Femmes, Médias)};
\node[actor, below right=of etat, align=center] (inter){COMMUNAUTÉ\\INTERNATIONALE\\(ONU, UA, CEEAC, UE,\\Partenaires bilatéraux)};
% Actions État
\node[action, right=of etat, xshift=1.5cm, align=center] (a1){Légiférer, budgeter,\\appliquer la loi};
\node[action, below right=of etat, xshift=1.5cm, align=center] (a2){Décentraliser, investir\\dans zones en crise};
\node[action, below=of etat, yshift=-1.5cm, align=center] (a3){Organiser le dialogue\\national (GDN)};
% Actions Société civile
\node[action, left=of civile, xshift=-1.5cm, align=center] (b1){Médiation communautaire,\\éducation à la paix};
\node[action, below left=of civile, xshift=-1.5cm, align=center] (b2){Plaidoyer, veille\\citoyenne, alerte};
\node[action, below=of civile, yshift=-1.5cm, align=center] (b3){Assistance humanitaire,\\réinsertion ex-combattants};
% Actions International
\node[action, right=of inter, xshift=1.5cm, align=center] (c1){Appui technique,\\financement, observation};
\node[action, below right=of inter, xshift=1.5cm, align=center] (c2){Médiation diplomatique,\\bon offices};
\node[action, below=of inter, yshift=-1.5cm, align=center] (c3){Aide humanitaire,\\renforcement capacités};
% Flèches
\draw[arrow] (etat) -- (a1);
\draw[arrow] (etat) -- (a2);
\draw[arrow] (etat) -- (a3);
\draw[arrow] (civile) -- (b1);
\draw[arrow] (civile) -- (b2);
\draw[arrow] (civile) -- (b3);
\draw[arrow] (inter) -- (c1);
\draw[arrow] (inter) -- (c2);
\draw[arrow] (inter) -- (c3);
% Interactions transverses
\draw[arrow, poppurple, dashed] (a3) to[bend left=20] node[midway, above, sloped, font=\tiny, poppurple] {Invite, facilite} (b1);
\draw[arrow, poppurple, dashed] (c1) to[bend right=20] node[midway, above, sloped, font=\tiny, poppurple] {Appuie} (a2);
\draw[arrow, poppurple, dashed] (b2) to[bend left=20] node[midway, below, sloped, font=\tiny, poppurple] {Interpelle} (a1);
% Groupes
\node[draw=poporange, thick, dashed, inner sep=5pt, fit=(etat)(a1)(a2)(a3), label={[poporange]above:Acteurs étatiques}] {};
\node[draw=poporange, thick, dashed, inner sep=5pt, fit=(civile)(b1)(b2)(b3), label={[poporange]above:Acteurs non étatiques}] {};
\node[draw=poporange, thick, dashed, inner sep=5pt, fit=(inter)(c1)(c2)(c3), label={[poporange]above:Partenaires internationaux}] {};
\end{tikzpicture}
\legende{Organigramme des acteurs de la promotion de la paix au Cameroun et de leurs actions principales. Les flèches pleines indiquent l'action directe ; les flèches pointillées violettes montrent les interactions et synergies entre acteurs.}
\end{popfigure}
\end{center}

\begin{center}
\begin{popfigure}
\begin{tikzpicture}[
    scale=0.85,
    event/.style={circle, fill=popblue, inner sep=2pt},
    date/.style={font=\small, text=popdark, anchor=west},
    desc/.style={font=\small, text=popmuted, anchor=west, text width=5.5cm},
    axis/.style={thick, popline},
    dot/.style={circle, fill=poporange, inner sep=3pt},
]
% Axe principal
\draw[axis] (0,0) -- (16,0);
\foreach \x in {0,2,4,6,8,10,12,14,16} \draw[axis] (\x,-0.1) -- (\x,0.1);
% Événements
% 1961
\node[dot] (e1961) at (0,0) {};
\node[date] at (0,0.5) {1961};
\node[desc] at (0.3,0.5) {Réunification : Cameroun Occidental + Cameroun Oriental = République Fédérale.};
% 1972
\node[dot] (e1972) at (2.5,0) {};
\node[date] at (2.5,0.5) {1972};
\node[desc] at (2.8,0.5) {Référendum : État unitaire. Fin du fédéralisme.};
% 1984
\node[dot] (e1984) at (4.5,0) {};
\node[date] at (4.5,-0.8) {1984};
\node[desc] at (4.8,-0.8) {Tentative de coup d'État. Renforcement de l'unité nationale.};
% 1990
\node[dot] (e1990) at (6.5,0) {};
\node[date] at (6.5,0.5) {1990};
\node[desc] at (6.8,0.5) {Lois sur les libertés (partis, presse, associations). Grand débat national.};
% 1996
\node[dot] (e1996) at (8.5,0) {};
\node[date] at (8.5,-0.8) {1996};
\node[desc] at (8.8,-0.8) {Constitution du 18 janv. : Décentralisation, bilinguisme, Cour suprême, Conseil constitutionnel.};
% 2006
\node[dot] (e2006) at (10.5,0) {};
\node[date] at (10.5,0.5) {2006};
\node[desc] at (10.8,0.5) {Opération Épervier (anti-corruption). Traité de Greentree (Bakassi).};
% 2014
\node[dot] (e2014) at (12,0) {};
\node[date] at (12,-0.8) {2014};
\node[desc] at (12.3,-0.8) {Loi antiterroriste. Début insurrection Boko Haram (Extrême-Nord).};
% 2016
\node[dot] (e2016) at (13.5,0) {};
\node[date] at (13.5,0.5) {2016};
\node[desc] at (13.8,0.5) {Crise anglophone : revendications corporatistes $\to$ conflit armé (NOSO).};
% 2019
\node[dot] (e2019) at (15,0) {};
\node[date] at (15,-0.8) {2019};
\node[desc] at (15.3,-0.8) {Grand Dialogue National. Statut spécial NOSO. Relance décentralisation.};
% 2020+
\node[dot] (e2020) at (16,0) {};
\node[date] at (16,0.5) {2020+};
\node[desc] at (16.3,0.5) {Mise en œuvre recommandations GDN. Élections régionales. DDR (Désarmement).};
% Flèche temps
\draw[->, thick, popdark] (16.5,0) -- (17,0) node[right, font=\small] {Temps};
\end{tikzpicture}
\legende{Frise chronologique : grandes étapes de la consolidation de la paix au Cameroun (1961--2020+). Les points orange marquent les jalons institutionnels et sécuritaires majeurs.}
\end{popfigure}
\end{center}

\courssub{Synthèse : articuler les stratégies}

Ces sept stratégies ne sont pas des tiroirs étanches ; elles s'\cle{imbriquent} et se \cle{renforcent mutuellement} :
\begin{itemize}
    \item L'\cle{éducation} forme les citoyens au dialogue.
    \item Le \cle{dialogue} légitime les réformes de la \cle{justice} et du \cle{développement}.
    \item La \cle{justice} protège les libertés médiatiques et les droits des acteurs non étatiques.
    \item Les \cle{médias} relayent les actions de la coopération internationale.
    \item Les \cle{acteurs non étatiques} comblent les vides de l'État et alertent la communauté internationale.
\end{itemize}

\begin{definition}[Promotion de la paix]
La \cle{promotion de la paix} est l'ensemble des actions structurelles et conjoncturelles, menées par l'État, la société civile et les partenaires internationaux, pour \emph{prévenir} l'éclatement des conflits violents, \emph{gérer} les crises en cours par des moyens non violents, et \emph{consolider} durablement la concorde nationale par la justice, l'équité et la participation citoyenne.
\end{definition}

\begin{methode}[Rédiger une conclusion argumentée (ECM / Épreuve écrite)]
Dans les épreuves du Probatoire, on vous demande souvent de « conclure » ou de « donner votre avis motivé ». Suivez cette démarche en 4 étapes :
\begin{enumerate}
    \item \textbf{Annoncez clairement votre position} (thèse) en une phrase : « En définitive, la promotion de la paix au Cameroun repose sur la complémentarité indispensable entre l'action étatique et l'engagement citoyen. »
    \item \textbf{Justifiez par au moins deux arguments distincts}, chacun illustré par un fait précis (loi, événement, institution, chiffre) : \textit{Argument 1 : Le cadre légal (Constitution 1996, Code collectivités 2019) pose les bases institutionnelles. Argument 2 : La société civile (GDN 2019, chefs traditionnels, ONG) assure l'ancrage local et la légitimité.}
    \item \textbf{Nuancez} (optionnel mais valorisant) : « Toutefois, la persistance des discours de haine en ligne montre que l'éducation aux médias reste un chantier prioritaire. »
    \item \textbf{Concluez par une ouverture} : « C'est à cette condition que le Cameroun relèvera le défi de l'émergence dans la paix et la stabilité. »
\end{enumerate}
\end{methode}

\begin{attention}[Erreur fréquente au Probatoire]
Confondre « \cle{pacification} » (action militaire/policière pour rétablir l'ordre) et « \cle{promotion de la paix} » (processus global, préventif, structurel, impliquant la société entière). Un exercice peut vous demander : « La paix est-elle seulement l'affaire de l'armée ? » $\rightarrow$ \textbf{Non}, argumenter avec les 7 stratégies ci-dessus.
\end{attention}

\begin{exoresolu}[Analyse d'une situation concrète]
\textbf{Énoncé :} Dans la commune de X (région de l'Extrême-Nord), des tensions récurrentes opposent éleveurs (peuls) et agriculteurs (massa, mouyeng) pour l'accès aux points d'eau et aux couloirs de transhumance. Le sous-préfet a réuni les chefs de village, l'imam, le curé, le président de la chambre d'agriculture, une ONG locale de médiation et les jeunes des deux communautés. Ensemble, ils ont balisé les couloirs, créé un comité de gestion des points d'eau paritaire, et organisé un tournoi de football intercommunautaire.

\textbf{Question :} Identifiez, dans ce récit, au moins \textbf{quatre stratégies de promotion de la paix} mises en œuvre et citez l'acteur principal pour chacune.

\tcblower
\textbf{Solution détaillée :}
\begin{enumerate}
    \item \textbf{Le dialogue et la concertation} : Acteur principal = \textbf{Sous-préfet} (représentant de l'État) qui convoque et anime la réunion inclusive.
    \item \textbf{Le rôle des acteurs non étatiques (chefs traditionnels \& confessions religieuses)} : Acteurs principaux = \textbf{Chefs de village, Imam, Curé} dont l'autorité morale légitime l'accord et facilite l'adhésion des communautés.
    \item \textbf{La justice et l'équité (règlement local)} : Acteur principal = \textbf{Comité de gestion paritaire} (instance locale de régulation) qui assure une application équitable des règles d'accès aux ressources (eau, pâture), prévenant l'arbitraire.
    \item \textbf{L'éducation à la paix / cohésion sociale par le sport} : Acteurs principaux = \textbf{Jeunes \& ONG locale} qui organisent le tournoi de football, vecteur de rencontre pacifique et de déconstruction des préjugés.
    \item (Bonus) \textbf{Développement équitable} : Le balisage des couloirs et l'aménagement des points d'eau constituent un investissement dans la prévention structurelle des conflits agro-pastoraux.
\end{enumerate}
\end{exoresolu}

\begin{exercice}[Entraînement Probatoire]
\textbf{Sujet :} « La promotion de la paix au Cameroun ne saurait être l'apanage du seul pouvoir public. » 
Appuyez-vous sur \textbf{trois stratégies distinctes} étudiées dans cette leçon pour argumenter cette affirmation. Votre réponse comportera une introduction, un développement structuré en trois paragraphes (une stratégie par paragraphe) et une conclusion argumentée (selon la méthode vue ci-dessus).
\end{exercice}

\begin{corrige}[Exercice n°1 -- Éléments de réponse attendus]
\textbf{Introduction :} Accroche (proverbe : « La main ne lave pas la main sans l'autre main »), définition de la promotion de la paix, thèse (nécessité d'une approche multi-acteurs), annonce du plan (3 stratégies).
\textbf{Développement :}
\begin{enumerate}
    \item \textbf{Dialogue et concertation} : Exemple GDN 2019 (forces vives, pas que gouvernement). Rôle des chefs traditionnels/religieux comme médiateurs légitimes.
    \item \textbf{Acteurs non étatiques (Société civile)} : ONG humanitaires (accès zones inaccessibles à l'État), médias communautaires (contre discours de haine), femmes/jeunes (cercles de paix).
    \item \textbf{Coopération internationale} : Appui technique/financier (UE, ONU) pour élections, DDR, décentralisation. Pression diplomatique pour respect DDH.
\end{enumerate}
\textbf{Conclusion :} Résumé des 3 arguments. Nuance : l'État reste le garant ultime (monopole violence légitime, cadre légal). Ouverture : « La paix est un chantier permanent où chaque Camerounais est à la fois ouvrier et architecte. »
\end{corrige}

\coursec{Le règlement pacifique des différends}

Dans la section précédente, nous avons étudié les \cle{stratégies de promotion de la paix} : l'éducation, la justice sociale, le dialogue, la culture de la tolérance. Mais promouvoir la paix ne suffit pas toujours : malgré toutes les précautions, des conflits naissent inévitablement entre les personnes, les groupes et même les États. Deux voisins se disputent une parcelle de terrain, deux élèves se reprochent la disparition d'un téléphone, deux pays revendiquent la même presqu'île. La question n'est donc pas d'éviter tout conflit — c'est impossible — mais de savoir \emph{comment} le résoudre sans violence. C'est l'objet de cette section : le \cle{règlement pacifique des différends}.

\courssub{Qu'est-ce qu'un différend ?}

\begin{definition}[Différend]
Un \cle{différend} est un désaccord ou un conflit d'intérêts entre des personnes, des groupes ou des États. Il oppose au moins deux parties qui revendiquent des positions incompatibles : un droit, un bien, un territoire, une reconnaissance.
\end{definition}

L'intuition est simple : dès qu'il y a vie en société, il y a des intérêts qui se croisent. Deux apprentis veulent utiliser la même machine à l'atelier ; deux familles revendiquent le même champ hérité d'un grand-père ; deux pays se disputent une zone de pêche. Le différend n'est pas en soi une catastrophe : il devient dangereux lorsqu'on choisit de le régler par la \cle{violence} (coups, destructions, guerre) au lieu de le régler par le \cle{droit} et le \cle{dialogue}.

\begin{definition}[Règlement pacifique des différends]
Le \cle{règlement pacifique des différends} est l'ensemble des procédures non violentes permettant de résoudre un conflit : la négociation, la médiation, la conciliation, l'arbitrage et la voie judiciaire. Il repose sur le dialogue, le droit et la recherche d'une solution acceptée par les parties.
\end{definition}

\begin{aretenir}[Les grands modes de règlement pacifique]
\begin{enumerate}
\item La \cle{négociation directe} : les parties discutent elles-mêmes.
\item La \cle{médiation} et la \cle{conciliation} : un tiers aide les parties à s'entendre.
\item L'\cle{arbitrage} : un arbitre tranche et sa décision s'impose.
\item La \cle{voie judiciaire} : un tribunal juge selon la loi.
\end{enumerate}
Ces modes vont du plus souple au plus contraignant : on commence en général par la négociation, et l'on ne saisit le juge qu'en cas d'échec.
\end{aretenir}

\courssub{La négociation directe}

\begin{definition}[Négociation]
La \cle{négociation} est une discussion directe entre les parties en conflit, en vue de trouver elles-mêmes un accord. Chacune expose ses prétentions, écoute celles de l'autre et accepte de faire des concessions réciproques.
\end{definition}

C'est le mode le plus simple, le plus rapide et le moins coûteux. Il ne fait intervenir personne d'autre que les parties. Son succès exige trois conditions : la \cle{bonne foi} (ne pas mentir), l'\cle{écoute} (comprendre les intérêts de l'autre et pas seulement ses positions) et la volonté de \cle{compromis} (accepter de ne pas tout obtenir).

\begin{exemplebox}[Négociation à l'atelier]
Deux apprentis d'un atelier de menuiserie de Douala veulent utiliser la scie circulaire au même moment. Au lieu de se disputer l'outil, ils discutent : l'un a une commande urgente pour le soir même, l'autre un travail pour le lendemain. Ils conviennent que le plus pressé utilise la machine le matin, et l'autre l'après-midi. Le différend est réglé par négociation directe, sans chef, sans juge, sans violence.
\end{exemplebox}

\courssub{La médiation et la conciliation}

Lorsque la négociation directe échoue — les parties ne se parlent plus, ou campent sur leurs positions — on fait appel à un \cle{tiers}.

\begin{definition}[Médiation et conciliation]
La \cle{médiation} est l'intervention d'un tiers neutre, le \cle{médiateur}, qui aide les parties à renouer le dialogue et leur propose des solutions. La \cle{conciliation} est une forme voisine où le tiers, le conciliateur, rapproche activement les positions en formulant des propositions concrètes. Dans les deux cas, la solution n'est pas imposée : elle n'a de valeur que si les parties l'acceptent librement.
\end{definition}

Au Cameroun, les médiateurs sont nombreux et proches des citoyens :
\begin{itemize}
\item le \cle{chef traditionnel} ou le notable du quartier, pour les conflits de voisinage, de terres ou de famille ;
\item le \cle{conseil de famille}, pour les différends entre époux ou entre héritiers ;
\item le \cle{Médiateur de la République}, institution nationale chargée de recevoir les plaintes des citoyens contre les administrations et de rechercher des solutions amiables ;
\item les comités de sages, les responsables religieux, les délégués de classe ou de personnel dans les écoles et les entreprises.
\end{itemize}

\begin{attention}[Limite de la justice traditionnelle de conciliation]
La justice traditionnelle de conciliation (palabre, conseil de famille, sentence des notables) est utile et reconnue dans nos coutumes, mais elle \textbf{ne doit jamais violer les droits humains ni contredire la loi républicaine}. Une « sentence » coutumière qui humilierait une femme, priverait un enfant de ses droits successoraux ou imposerait un châtiment corporel serait illégale : la Constitution et les lois du Cameroun priment sur la coutume.
\end{attention}

\courssub{Les mécanismes traditionnels : la palabre et le conseil de famille}

Bien avant les tribunaux modernes, les sociétés camerounaises ont développé leurs propres modes de règlement pacifique. La \cle{palabre} est une assemblée publique où la parole circule librement : chacun expose son point de vue devant les anciens, qui recherchent un consensus. Le \cle{conseil de famille} réunit les aînés pour régler les conflits successoraux, matrimoniaux ou de voisinage. La \cle{justice coutumière de conciliation} vise moins à punir qu'à \emph{réconcilier} : le coupable reconnaît sa faute, répare le tort causé (compensation, excuses publiques, partage d'un repas de réconciliation), et l'harmonie du groupe est restaurée.

\begin{savaistu}[La réconciliation plutôt que la punition]
Dans de nombreuses traditions camerounaises, le conflit se termine par un geste symbolique : le partage de la noix de cola, un repas commun, une poignée de main devant témoins. L'objectif n'est pas de désigner un vainqueur et un vaincu, mais de rétablir le \cle{vivre ensemble} : les deux parties continueront à se croiser chaque jour au village ou au quartier. Cette sagesse africaine inspire aujourd'hui la « justice réparatrice » pratiquée dans plusieurs pays du monde.
\end{savaistu}

\courssub{L'arbitrage}

\begin{definition}[Arbitrage]
L'\cle{arbitrage} est le recours à un tiers appelé \cle{arbitre}, choisi ou accepté par les parties, dont la décision — la \emph{sentence arbitrale} — s'impose aux parties comme un jugement. Contrairement au médiateur, l'arbitre \textbf{tranche} le différend.
\end{definition}

L'arbitrage est très utilisé dans les conflits commerciaux et du travail (au Cameroun, il existe par exemple le Centre d'arbitrage du GICAM pour les litiges entre entreprises) et dans les conflits internationaux, lorsque deux États conviennent de soumettre leur différend à des arbitres. Sa force : une décision rapide, par des spécialistes, et définitive. Sa condition : les parties doivent avoir accepté à l'avance de se soumettre à la sentence.

\courssub{La voie judiciaire : les tribunaux}

Quand tout le reste échoue, ou lorsque le conflit porte sur un droit que la loi protège, il reste la \cle{voie judiciaire} : saisir un tribunal qui jugera selon la loi.

\begin{definition}[Voie judiciaire]
La \cle{voie judiciaire} consiste à porter le différend devant un juge, qui entend les parties, examine les preuves et rend une décision obligatoire, fondée sur la loi. Au plan national, ce sont les tribunaux camerounais (tribunaux de première instance, cours d'appel, Cour suprême). Au plan international, les États peuvent saisir la \cle{Cour Internationale de Justice} (CIJ), organe judiciaire principal de l'ONU, qui siège à La Haye et règle les différends entre États.
\end{definition}

\begin{attention}[Ne pas confondre]
Le \cle{médiateur} propose, l'\cle{arbitre} et le \cle{juge} tranchent. La médiation exige l'accord final des deux parties ; la sentence arbitrale et le jugement s'imposent à elles. Autre piège : la CIJ ne juge \textbf{que des États}, jamais des individus, et seulement si les États concernés ont accepté sa compétence.
\end{attention}

\begin{popfigure}
\begin{center}
\begin{tikzpicture}[
font=\small,
node distance=0.55cm and 1.1cm,
etape/.style={rectangle, rounded corners, draw=popblue, fill=popblueL, thick, text width=3.6cm, align=center, minimum height=1.05cm},
decision/.style={rectangle, rounded corners, draw=poporange, fill=poporangeL, thick, text width=2.5cm, align=center, minimum height=0.9cm},
fleche/.style={-{Stealth[length=2.5mm]}, thick, popdark},
lbl/.style={font=\scriptsize\itshape, popmuted}
]
\node[etape, align=center] (neg){\textbf{1. Négociation directe}\\Les parties discutent seules};
\node[etape, right=of neg, align=center] (med){\textbf{2. Médiation / conciliation}\\Un tiers aide au dialogue};
\node[etape, right=of med, align=center] (arb){\textbf{3. Arbitrage}\\L'arbitre tranche};
\node[etape, right=of arb, align=center] (jug){\textbf{4. Voie judiciaire}\\Le tribunal juge};
\node[decision, below=1.1cm of neg] (ok1) {Accord ?};
\node[decision, below=1.1cm of med] (ok2) {Accord ?};
\node[decision, below=1.1cm of arb] (ok3) {Sentence};
\node[decision, below=1.1cm of jug] (ok4) {Jugement};
\draw[fleche] (neg) -- (med) node[midway, above, lbl]{échec};
\draw[fleche] (med) -- (arb) node[midway, above, lbl]{échec};
\draw[fleche] (arb) -- (jug) node[midway, above, lbl]{refus / échec};
\draw[fleche] (neg) -- (ok1);
\draw[fleche] (med) -- (ok2);
\draw[fleche] (arb) -- (ok3);
\draw[fleche] (jug) -- (ok4);
\node[below=0.15cm of ok1, font=\scriptsize\itshape, popgreen, text width=2.6cm, align=center] {si oui : conflit réglé};
\node[below=0.15cm of ok2, font=\scriptsize\itshape, popgreen, text width=2.6cm, align=center] {si oui : conflit réglé};
\node[below=0.15cm of ok3, font=\scriptsize\itshape, popgreen, text width=2.6cm, align=center] {décision obligatoire};
\node[below=0.15cm of ok4, font=\scriptsize\itshape, popgreen, text width=2.6cm, align=center] {décision obligatoire};
\end{tikzpicture}
\end{center}
\legende{Arbre de décision : les modes de règlement pacifique d'un différend, du plus souple (négociation) au plus contraignant (jugement). On ne passe au niveau supérieur qu'en cas d'échec du niveau précédent.}
\end{popfigure}

\courssub{Méthode : résoudre un conflit à l'école ou en famille}

\begin{methode}[Résoudre pacifiquement un conflit quotidien]
\begin{enumerate}
\item \textbf{Se calmer et écouter les deux parties} : laisser chacun raconter les faits sans l'interrompre, sans insulter, sans crier.
\item \textbf{Identifier les intérêts réels} : derrière les positions (« je veux ceci »), chercher les besoins (« j'ai besoin de respect, de sécurité, de reconnaissance »).
\item \textbf{Proposer plusieurs solutions} : imaginer ensemble au moins deux ou trois issues possibles, sans en rejeter aucune d'emblée.
\item \textbf{Rechercher un compromis acceptable} : choisir la solution que chacun peut accepter, même si elle ne satisfait personne à 100 \%.
\item \textbf{Sceller l'accord et réparer} : excuses, réparation du tort, poignée de main ; si le conflit persiste, faire appel à un tiers (délégué, censeur, notable, médiateur).
\end{enumerate}
\end{methode}

\begin{exoresolu}[Un conflit au lycée technique]
\textbf{Énoncé.} Dans un atelier d'électrotechnique, un outil de mesure coûteux disparaît. L'apprenti A accuse publiquement l'apprenti B de vol. B nie et menace A de le frapper après les cours. Le chef d'atelier intervient. Décris la démarche pacifique à suivre et indique quel mode de règlement est adapté à chaque étape.
\tcblower
\textbf{Solution détaillée.}
\begin{enumerate}
\item \textbf{Stopper l'escalade} : le chef d'atelier sépare les deux apprentis et interdit toute menace. La violence transformerait un simple différend en faute grave.
\item \textbf{Écouter les parties séparément} : A explique pourquoi il soupçonne B ; B donne sa version. On découvre que l'outil a été rangé par erreur dans une autre armoire par un troisième élève.
\item \textbf{Identifier les intérêts} : A voulait retrouver l'outil ; B veut que son honneur soit rétabli après l'accusation publique.
\item \textbf{Réparer} : A présente des excuses publiques à B devant l'atelier ; B retire ses menaces. C'est une \cle{négociation} facilitée par une \cle{médiation} (le chef d'atelier).
\item \textbf{Si la médiation avait échoué} : le différend serait remonté au conseil de discipline (voie quasi judiciaire interne), et en cas de vol avéré ou de coups, devant les tribunaux.
\end{enumerate}
\textbf{Conclusion} : la plupart des conflits quotidiens se règlent par l'écoute, la vérité des faits et la réparation, sans juge ni violence.
\end{exoresolu}

\courssub{Étude de cas : le différend de Bakassi (Cameroun -- Nigéria)}

L'exemple le plus célèbre de règlement pacifique d'un différend en Afrique concerne notre pays : l'affaire de la péninsule de \cle{Bakassi}.

\textbf{Les faits.} Bakassi est une presqu'île marécageuse d'environ \SI{665}{km^2}, riche en poissons et en pétrole, située dans le golfe de Guinée, à la frontière entre le Cameroun et le Nigéria. Le Cameroun s'appuie sur des traités coloniaux (accord anglo-allemand de 1913) ; le Nigéria revendique la presqu'île et y installe des populations et des soldats. Des affrontements armés éclatent, notamment en 1981 et en 1994 : le différend menace de devenir une guerre entre les deux plus grands pays d'Afrique centrale et de l'Ouest.

\textbf{La saisine de la CIJ (1994).} Plutôt que la guerre, le Cameroun choisit le droit : le 29 mars 1994, il saisit la \cle{Cour Internationale de Justice} pour faire trancher la frontière. C'est un acte de courage civique et de confiance dans le droit international.

\textbf{L'arrêt du 10 octobre 2002.} Après huit ans de procédure (mémoires, plaidoiries, expertise), la CIJ rend son arrêt : elle attribue la péninsule de Bakassi au \cle{Cameroun}, en se fondant sur les traités coloniaux et le principe de l'intangibilité des frontières héritées de la colonisation. Le Nigéria, d'abord réticent, finit par accepter l'arrêt.

\textbf{Les accords de Greentree (12 juin 2006).} Sous l'égide du Secrétaire général de l'ONU, Kofi Annan, les présidents Paul Biya et Olusegun Obasanjo signent à Greentree (États-Unis) les accords organisant le retrait nigérian, le transfert progressif de l'autorité et la protection des populations nigérianes restant sur place.

\textbf{La rétrocession pacifique (2006 -- 2013).} Le retrait des troupes nigérianes commence dès 2006 ; le transfert de l'autorité civile est effectif en 2008 ; le processus est achevé le 14 août 2013. Pas un seul coup de feu n'a été tiré pour appliquer l'arrêt : le différend de Bakassi est réglé \emph{intégralement par des moyens pacifiques}.

\begin{exemplebox}[Chronologie chiffrée du règlement de Bakassi]
\begin{itemize}
\item \textbf{1994} : saisine de la CIJ par le Cameroun (29 mars).
\item \textbf{2002} : arrêt de la CIJ du 10 octobre, attribuant Bakassi au Cameroun.
\item \textbf{2006} : accords de Greentree (12 juin), sous l'égide de l'ONU.
\item \textbf{2008} : retrait complet des troupes nigérianes et transfert de l'autorité civile.
\item \textbf{2013} : achèvement du processus (14 août), sans guerre.
\end{itemize}
L'arrêt de la CIJ du 10 octobre 2002 a donc été appliqué sans guerre grâce aux accords de Greentree de juin 2006 : preuve que le droit peut l'emporter sur les armes.
\end{exemplebox}

\begin{popfigure}
\begin{center}
\begin{tikzpicture}[scale=1.0, font=\small]
% Mer
\fill[popblueL] (0,0) rectangle (12,3.2);
\node[popblue, font=\itshape] at (9.8,0.7) {Golfe de Guinée (océan Atlantique)};
% Nigeria
\fill[popgreenL] (0,3.2) -- (0,7) -- (5.2,7) -- (5.6,5.6) -- (5.2,4.4) -- (4.6,3.2) -- cycle;
\node[popdark, font=\bfseries] at (2.4,6.3) {NIGERIA};
% Cameroun
\fill[poporangeL] (4.6,3.2) -- (5.2,4.4) -- (5.6,5.6) -- (5.2,7) -- (12,7) -- (12,3.2) -- cycle;
\node[popdark, font=\bfseries] at (9.2,6.3) {CAMEROUN};
% Bakassi
\fill[poppinkL, draw=poppink, thick] (4.6,3.2) -- (5.2,4.4) -- (5.0,3.9) -- (4.2,3.4) -- (3.6,3.2) -- cycle;
\draw[poppink, thick, -{Stealth}] (1.6,1.2) -- (4.3,3.1);
\node[poppink, font=\bfseries, align=center] at (1.7,0.9) {Péninsule de\\Bakassi (665 km$^2$)};
% Frontière
\draw[popdark, very thick, dashed] (4.6,3.2) -- (5.2,4.4) -- (5.6,5.6) -- (5.2,7);
\node[popdark, font=\scriptsize, rotate=55] at (5.75,4.6) {frontière (CIJ, 2002)};
% Villes
\fill[popdark] (7.6,4.1) circle (1.6pt); \node[popdark, right, font=\scriptsize] at (7.6,4.1) {Douala};
\fill[popdark] (6.6,4.6) circle (1.6pt); \node[popdark, above left, font=\scriptsize] at (6.6,4.6) {Limbe};
\fill[popdark] (2.6,4.4) circle (1.6pt); \node[popdark, left, font=\scriptsize] at (2.6,4.4) {Calabar};
% Nord
\draw[-{Stealth}, very thick, popdark] (11.3,5.2) -- (11.3,6.6);
\node[popdark, above, font=\bfseries] at (11.3,6.6) {N};
% Échelle
\draw[popdark, thick] (0.4,2.6) -- (2.4,2.6);
\node[popdark, below, font=\scriptsize] at (1.4,2.6) {0 \quad $\approx$ 100 km};
\end{tikzpicture}
\end{center}
\legende{Carte très simplifiée de la péninsule de Bakassi : la frontière Cameroun -- Nigéria (pointillés) suit l'arrêt de la CIJ de 2002 ; la presqu'île (zone rose), riche en pêche et en pétrole, a été rétrocédée pacifiquement au Cameroun entre 2006 et 2013.}
\end{popfigure}

\begin{savaistu}[Un modèle pour l'Afrique]
Le règlement du différend de Bakassi est cité par l'\cle{ONU} comme un \textbf{modèle de résolution pacifique des conflits frontaliers en Afrique}. Alors que d'autres conflits territoriaux sur le continent ont dégénéré en guerres longues et meurtrières, le Cameroun et le Nigéria ont montré qu'un État peut défendre ses droits par le droit, et qu'un État qui perd un procès peut se retirer avec dignité. Cet exemple est enseigné dans les écoles de diplomatie du monde entier.
\end{savaistu}

\begin{aretenir}[L'essentiel de la section]
\begin{itemize}
\item Un \cle{différend} est un désaccord ou un conflit d'intérêts ; son \cle{règlement pacifique} utilise le dialogue et le droit, jamais la violence.
\item Les modes de règlement vont du plus souple au plus contraignant : \cle{négociation} $\rightarrow$ \cle{médiation}/conciliation $\rightarrow$ \cle{arbitrage} $\rightarrow$ \cle{voie judiciaire}.
\item Les mécanismes traditionnels (palabre, conseil de famille) visent la réconciliation, mais doivent respecter les droits humains et la loi.
\item Le différend de \cle{Bakassi} (saisine de la CIJ en 1994, arrêt de 2002, accords de Greentree en 2006, rétrocession achevée en 2013) prouve qu'un conflit entre États peut se régler sans guerre.
\end{itemize}
\end{aretenir}

\begin{exercice}[Pour s'entraîner]
\begin{enumerate}
\item Définis les termes suivants : différend, médiation, arbitrage.
\item Classe ces situations selon le mode de règlement le plus adapté : (a) deux élèves se disputent une place ; (b) deux entreprises se disputent un marché ; (c) deux États se disputent une frontière ; (d) deux époux se disputent la garde d'un enfant.
\item Explique en cinq lignes pourquoi l'affaire de Bakassi est un modèle de règlement pacifique.
\item Pourquoi une sentence coutumière contraire à la loi républicaine est-elle inacceptable ?
\end{enumerate}
\end{exercice}

\begin{corrige}[Exercice : Pour s'entraîner]
\begin{enumerate}
\item \textbf{Différend} : désaccord ou conflit d'intérêts entre personnes, groupes ou États. \textbf{Médiation} : intervention d'un tiers neutre qui aide les parties à trouver elles-mêmes un accord. \textbf{Arbitrage} : recours à un arbitre dont la décision s'impose aux parties.
\item (a) négociation directe, éventuellement médiation du délégué ; (b) arbitrage commercial ; (c) voie judiciaire internationale (CIJ) ; (d) médiation familiale puis tribunal en cas d'échec.
\item Le Cameroun et le Nigéria ont choisi le droit plutôt que la guerre : saisine de la CIJ (1994), arrêt accepté (2002), accords de Greentree (2006) et rétrocession sans effusion de sang (2008-2013). L'ONU le cite comme modèle africain.
\item Parce que la Constitution et les lois de la République priment sur la coutume : aucune tradition ne peut justifier une violation des droits humains (humiliation, privation des droits d'un enfant, châtiment corporel).
\end{enumerate}
\end{corrige}

\coursec{Dossier 2 : La lutte contre l'extrémisme violent au Cameroun}

Dans la section précédente, nous avons étudié le \cle{règlement pacifique des différends} : le dialogue, la médiation, l'arbitrage et la justice permettent de résoudre les conflits sans violence. Mais que se passe-t-il lorsqu'un groupe refuse par principe tout dialogue et choisit délibérément la violence pour imposer ses idées ? C'est exactement le problème posé par l'\cle{extrémisme violent}, un phénomène qui a frappé durement la région de l'Extrême-Nord du Cameroun depuis 2014. Ce dossier montre comment un État et ses citoyens peuvent lutter contre ce fléau tout en préservant la paix et le vivre ensemble.

\courssub{Comprendre l'extrémisme et l'extrémisme violent}

\textbf{Intuition.} Imagine un débat en classe : deux élèves ne sont pas d'accord sur un sujet. Le désaccord est normal et même utile. Mais si l'un des deux décide que l'autre n'a \emph{pas le droit} de penser différemment, qu'il doit être puni, humilié, voire frappé pour ses idées, on quitte le désaccord pour entrer dans l'extrémisme. L'extrémisme commence quand on refuse à l'autre le droit d'être différent.

\begin{definition}[Extrémisme]
L'\cle{extrémisme} est une attitude de \textbf{rejet radical de l'autre} : celui qui ne partage pas mes idées, ma religion ou mon mode de vie est considéré comme un ennemi à éliminer ou à exclure. L'extrémiste refuse le débat, le compromis et la diversité.
\end{definition}

\begin{definition}[Extrémisme violent et terrorisme]
L'\cle{extrémisme violent} est une \textbf{idéologie qui justifie l'usage de la violence} (attaques, enlèvements, destructions) pour imposer ses vues et rejeter toute différence. Lorsque cette violence vise délibérément des populations civiles pour semer la terreur et faire plier un État ou une société, on parle de \cle{terrorisme}.
\end{definition}

\begin{aretenir}[Trois notions à ne pas confondre]
\begin{itemize}
\item \textbf{Extrémisme} : rejet radical de l'autre et de la différence (attitude, idéologie).
\item \textbf{Extrémisme violent} : passage à l'acte violent au nom de cette idéologie.
\item \textbf{Terrorisme} : violence organisée contre des civils pour semer la peur et imposer un changement politique ou social.
\end{itemize}
\end{aretenir}

\begin{attention}[Ne pas assimiler une religion à l'extrémisme]
Piège majeur : dire « les terroristes se réclament de telle religion, donc cette religion est violente ». C'est \textbf{faux}. Les groupes terroristes \textbf{manipulent} les textes religieux, les sortent de leur contexte et les déforment pour recruter des jeunes. La grande majorité des croyants, de toutes les religions, rejette la violence. Au Cameroun, les premières victimes de Boko Haram dans l'Extrême-Nord sont d'ailleurs très souvent des musulmans qui refusaient l'idéologie du groupe. Confondre religion et terrorisme, c'est faire le jeu des extrémistes, qui veulent diviser les communautés.
\end{attention}

\courssub{Les manifestations de l'extrémisme violent au Cameroun}

Depuis \textbf{2014}, la secte terroriste \cle{Boko Haram} (officiellement \emph{Jama'atu Ahlis Sunna Lidda'awati wal-Jihad}, « Groupe du peuple de la Sunna pour la prédication et le jihad »), née au Nigeria voisin, étend ses exactions à la région de l'\textbf{Extrême-Nord} du Cameroun, notamment dans les départements du \textbf{Logone-et-Chari}, du \textbf{Mayo-Sava} et du \textbf{Mayo-Tsanaga}, frontaliers du Nigeria et proches du lac Tchad. Le groupe se réclame de l'État islamique (Daech) depuis 2015 sous le nom d'\cle{ISWAP} (Islamic State West Africa Province).

Les principales manifestations sont :
\begin{itemize}
\item des \textbf{attaques} contre des villages, des marchés, des églises, des mosquées et des positions militaires ;
\item des \textbf{enlèvements} de civils, y compris des femmes et des enfants (certains utilisés comme boucliers humains ou kamikazes) ;
\item des \textbf{attentats-suicides} commis par des kamikazes, souvent de jeunes filles manipulées, dans des lieux publics (marchés, mosquées, arrêts de bus) ;
\item des \textbf{pillages} de récoltes et de bétail, des incendies de maisons, d'écoles et de centres de santé ;
\item le \textbf{recrutement forcé ou par séduction} de jeunes déscolarisés, sans emploi, par endoctrinement et promesses financières.
\end{itemize}

\begin{popfigure}
\begin{center}
\begin{tikzpicture}[scale=1.0, font=\small]
% Contour simplifié de l'Extrême-Nord
\draw[thick, popdark, fill=popcream] (0,0) -- (0,4.6) -- (1.2,5.4) -- (2.6,5.8) -- (3.8,5.2) -- (4.6,4.2) -- (4.4,2.6) -- (3.6,1.2) -- (2.2,0) -- cycle;
% Frontière Nigeria (ouest)
\draw[very thick, poppurple, dashed] (0,0) -- (0,4.6) -- (1.2,5.4);
\node[popdark, rotate=78] at (-0.35,2.6) {\footnotesize Nigeria};
% Lac Tchad (nord-est)
\draw[fill=popblue!40, draw=popblue] (2.6,5.8) -- (3.8,5.2) -- (4.2,5.9) -- (3.0,6.3) -- cycle;
\node[popblue] at (3.5,6.0) {\footnotesize Lac Tchad};
% Départements
\draw[fill=poporange!45, draw=poporange, thick] (0.1,3.4) -- (1.1,4.5) -- (2.5,5.6) -- (3.6,5.1) -- (3.2,3.6) -- (1.4,3.2) -- cycle;
\node[popdark] at (1.9,4.4) {\footnotesize \textbf{Logone-et-Chari}};
\draw[fill=popgold!45, draw=popgold, thick] (0.1,3.3) -- (1.4,3.1) -- (3.1,3.5) -- (3.0,2.2) -- (1.2,1.9) -- (0.1,2.2) -- cycle;
\node[popdark] at (1.6,2.7) {\footnotesize \textbf{Mayo-Sava}};
\draw[fill=popgreen!35, draw=popgreen, thick] (0.1,2.1) -- (1.2,1.8) -- (2.9,2.1) -- (3.5,1.3) -- (2.2,0.2) -- (0.1,0.1) -- cycle;
\node[popdark] at (1.5,1.1) {\footnotesize \textbf{Mayo-Tsanaga}};
% Villes
\fill[popink] (2.7,5.0) circle (0.07); \node[right, popdark] at (2.8,5.0) {\footnotesize Fotokol};
\fill[popink] (1.8,3.9) circle (0.07); \node[right, popdark] at (1.9,3.9) {\footnotesize Kousséri};
\fill[popink] (2.2,2.6) circle (0.07); \node[right, popdark] at (2.3,2.6) {\footnotesize Mora};
\fill[popink] (1.6,0.8) circle (0.07); \node[right, popdark] at (1.7,0.8) {\footnotesize Mokolo};
\fill[popink] (3.4,1.6) circle (0.09); \node[right, popdark] at (3.5,1.6) {\footnotesize \textbf{Maroua}};
% Nord
\draw[-{Stealth}, thick, popdark] (5.0,4.6) -- (5.0,5.6);
\node[popdark] at (5.0,5.85) {\footnotesize N};
\end{tikzpicture}
\end{center}
\legende{Carte schématique de l'Extrême-Nord du Cameroun : les trois départements les plus touchés par les exactions de Boko Haram / ISWAP (Logone-et-Chari, Mayo-Sava, Mayo-Tsanaga), situés le long de la frontière nigériane et du lac Tchad. Croquis simplifié, sans échelle exacte.}
\end{popfigure}

\courssub{Les causes de l'extrémisme violent}

Pourquoi des jeunes camerounais rejoignent-ils un groupe terroriste ? L'enquête sociale (interviews d'anciens membres, rapports des ONG, du PNUD, de l'UNICEF et des autorités) fait ressortir plusieurs causes combinées, souvent résumées par l'approche \textbf{« pousse-tire »} (push-pull factors) :

\begin{itemize}
\item \textbf{Facteurs de poussée (push) :}
  \begin{itemize}
  \item La \textbf{pauvreté} : l'Extrême-Nord est l'une des régions les plus pauvres du pays (taux de pauvreté supérieur à 70\% selon les enquêtes ECAM) ; sans ressources, certains jeunes sont attirés par l'argent promis par les recruteurs.
  \item Le \textbf{chômage et le sous-emploi des jeunes} : sans emploi décent ni perspective, un jeune devient une proie facile pour ceux qui lui offrent un « rôle », un revenu et une identité.
  \item La \textbf{faiblesse de la présence de l'État} : dans certaines zones reculées, l'absence d'écoles, de routes, de services de santé, de sécurité et d'administration laisse le champ libre aux groupes armés.
  \item La \textbf{marginalisation et le sentiment d'injustice} : le sentiment d'être oublié, méprisé ou exclu nourrit la rancune et le désir de vengeance.
  \end{itemize}
\item \textbf{Facteurs d'attraction (pull) :}
  \begin{itemize}
  \item La \textbf{manipulation religieuse et idéologique} : les recruteurs déforment les textes sacrés pour faire croire que la violence est un devoir sacré et promettent un paradis aux « martyrs ».
  \item L'\textbf{appartenance et la protection} : le groupe offre une « famille », une structure, des armes et un statut social à des jeunes en quête de reconnaissance.
  \item Les \textbf{revenus immédiats} : butin de guerre, rançons, salaire mensuel versé par le groupe.
  \end{itemize}
\end{itemize}

\begin{aretenir}[Une cause n'explique jamais tout]
L'extrémisme violent naît d'un \textbf{faisceau de causes} : misère économique + manipulation idéologique + sentiment d'abandon + absence de perspectives. C'est pourquoi la réponse doit être à la fois sécuritaire, sociale, éducative et économique.
\end{aretenir}

\courssub{Les conséquences pour les populations et le pays}

Les exactions de Boko Haram / ISWAP ont des effets lourds et durables :

\begin{itemize}
\item \textbf{Coût humain :} des milliers de morts et de blessés parmi les civils et les forces de défense ; des veuves, des orphelins, des traumatismes psychologiques durables.
\item \textbf{Crise de déplacement :} plus de \textbf{300 000 déplacés internes} (chiffres OIM/HCR 2023) et des dizaines de milliers de réfugiés nigérians ; des familles entières fuient leurs villages et vivent dans des camps (Minawao, Zamai, etc.) ou chez des parents, dans la précarité.
\item \textbf{Effondrement éducatif :} des centaines d'écoles fermées ou détruites, privant des dizaines de milliers d'enfants d'instruction ; recrutement facilité par la déscolarisation.
\item \textbf{Choc économique :} les marchés frontaliers (Kousseri, Fotokol, Amchidé), l'agriculture, l'élevage et la pêche sur le lac Tchad sont perturbés ; le commerce transfrontalier avec le Nigeria diminue drastiquement.
\item \textbf{Fracture sociale :} psychose collective, méfiance intercommunautaire, stigmatisation de certaines populations (Kanouri, Peuls), tensions entre communautés d'accueil et déplacés.
\item \textbf{Atteinte aux services de base :} centres de santé fermés, personnel soignant en fuite, pénurie de médicaments, rupture des campagnes de vaccination.
\end{itemize}

\begin{popfigure}
\begin{center}
\begin{tikzpicture}
\begin{axis}[
  ybar, bar width=0.55cm,
  width=0.9\linewidth, height=6.5cm,
  xlabel={Année}, ylabel={Nombre d'incidents violents (estimation ACLED/ONU)},
  symbolic x coords={2014,2015,2016,2017,2018,2019,2020,2021,2022,2023},
  xtick=data, x tick label style={rotate=35, anchor=east, font=\footnotesize},
  ymin=0, ymax=350,
  nodes near coords, nodes near coords style={font=\tiny},
  axis lines=left, grid=major, grid style={popline!40},
  title style={font=\small}, title={Évolution des incidents violents attribués à Boko Haram/ISWAP dans l'Extrême-Nord}
]
\addplot[fill=poporange, draw=popdark] coordinates {
  (2014,50) (2015,180) (2016,220) (2017,160) (2018,130) (2019,110) (2020,95) (2021,85) (2022,70) (2023,65)
};
\end{axis}
\end{tikzpicture}
\end{center}
\legende{Évolution indicative du nombre d'incidents violents (attaques, affrontements, attentats) attribués à Boko Haram / ISWAP dans l'Extrême-Nord (données synthétisées ACLED/ONU, arrondies, à visée pédagogique). On observe un pic vers 2015-2016, puis une baisse progressive liée à la riposte militaire, à la coopération régionale et à la mobilisation citoyenne, sans disparition totale de la menace (résurgence ponctuelle en 2023-2024).}
\end{popfigure}

\textbf{Commentaire du diagramme.} La courbe en barres montre une montée rapide des incidents à partir de 2014, un maximum autour de 2015-2016, puis un recul net et soutenu. Ce recul ne signifie pas la fin du danger : il traduit l'efficacité combinée des opérations militaires, de la coopération internationale (FMM), des comités de vigilance et des programmes de déradicalisation. La menace a muté : moins d'attaques massives, mais plus d'embuscades, d'engins explosifs improvisés (EEI) et de recrutement sournois. La vigilance reste donc nécessaire.

\courssub{La réponse sécuritaire de l'État}

Face à la violence, l'État camerounais a d'abord répondu par des moyens de défense, dans le cadre de la \textbf{Loi n°2014/028 du 23 décembre 2014 relative à la lutte contre le terrorisme} (modifiée par la Loi n°2017/012) :

\begin{itemize}
\item des \textbf{opérations militaires} menées par les Forces de Défense (BIR, Bataillon d'Intervention Rapide ; Gendarmerie ; Police) dans les zones frontalières : opérations \emph{Alpha}, \emph{Épervier}, sécurisation des axes routiers, démantèlement de camps terroristes, libération d'otages ;
\item la participation active à la \cle{Force Multinationale Mixte (FMM)}, créée en 1994 (réactivée en 2012, opérationnelle en 2015), qui associe le Cameroun, le Nigeria, le Tchad, le Niger et le Bénin sous commandement unifié (quartier général à N'Djamena) pour combattre ensemble Boko Haram / ISWAP dans le bassin du lac Tchad ;
\item l'appui aux \cle{comités de vigilance} (CV), groupes d'habitants volontaires organisés par les autorités administratives (préfets, sous-préfets) qui surveillent les villages, contrôlent les entrées/sorties, signalent les mouvements suspects et coopèrent avec les forces de défense ; ils sont encadrés par la \textbf{Loi n°2019/020 du 24 décembre 2019} portant organisation de la coopération entre les forces de défense et les populations.
\end{itemize}

\begin{exemplebox}[Les comités de vigilance : la citoyenneté en action]
Dans le Mayo-Sava et le Logone-et-Chari, des jeunes volontaires (souvent d'anciens chasseurs traditionnels, cultivateurs, élèves) patrouillent la nuit, contrôlent les entrées des villages et informent les autorités via talkies-walkies ou téléphones. Les \textbf{comités de vigilance} illustrent la \textbf{participation citoyenne à la sécurité}, aux côtés des forces de défense : la sécurité n'est pas seulement l'affaire de l'armée, elle concerne chaque citoyen. Ces comités ont permis de déjouer de nombreuses attaques, d'identifier des kamikazes et de rassurer les populations. Leur action est encadrée par la loi pour éviter les dérives (justice privée, abus).
\end{exemplebox}

\courssub{La réponse préventive, sociale et juridique}

La force seule ne suffit pas : on ne tue pas une idéologie avec des armes. L'État, avec l'appui de partenaires (PNUD, UNICEF, UNESCO, UE, ONG nationales/internationales) et la société civile, mène une action de fond structurée autour de la \textbf{Stratégie Nationale de Lutte contre l'Extrémisme Violent (SNLEV)} et du \textbf{Plan d'Urgence Triennal d'Adaptation du Secteur Éducation (PLANUT)} pour l'Extrême-Nord :

\begin{itemize}
\item \textbf{Cadre juridique et judiciaire :} application de la loi antiterroriste (poursuites, jugements, peines) ; création de sections spécialisées au Tribunal Militaire de Yaoundé ; respect des droits humains dans la procédure pénale (formation des magistrats, avocats).
\item \textbf{Déradicalisation et désengagement :} centres de prise en charge (ex. centre de Mérina à Maroua, centre de Mokolo) offrant un accompagnement psychologique, théologique (avec des imams et prêtres formés), social et juridique pour déconstruire l'idéologie extrémiste.
\item \textbf{Réinsertion socio-économique des repentis :} formation professionnelle (couture, menuiserie, soudure, agriculture, élevage, informatique), kits de réinsertion, appui à la création d'activités génératrices de revenus (AGR), médiation communautaire pour l'acceptation par les familles et villages.
\item \textbf{Éducation et résilience :} réouverture et construction d'écoles (classes temporaires, écoles sous tentes), cantines scolaires, distribution de kits scolaires, enseignement de la citoyenneté, de l'esprit critique et de la résolution non-violente des conflits (programme ECM) ; formation des enseignants à la pédagogie protectrice.
\item \textbf{Développement intégré de l'Extrême-Nord :} \textbf{Plan d'Urgence Triennal (PLANUT 2022-2024)} et \textbf{Plan Régional de Développement} : routes (axe Kousseri-Mora-Mokolo), eau potable, électricité, centres de santé, marchés modernes, appui à l'agriculture résiliente (semences améliorées, irrigation), emploi des jeunes (Programme Spécial Jeunes, PAJER-U, PIAASI).
\item \textbf{Dialogue interreligieux et interculturel :} plateformes de concertation (Cadre de Concertation des Confessions Religieuses du Cameroun), prêches conjoints pour la paix, dénonciation commune de la manipulation religieuse, célébrations communes (fêtes nationales, journées internationales).
\item \textbf{Lutte contre la propagande en ligne :} surveillance du cyberespace, contre-narratifs sur les réseaux sociaux, éducation aux médias et à l'information (EMI) dans les lycées et centres de formation.
\end{itemize}

\begin{savaistu}[Des centres de déradicalisation au Cameroun]
Le Cameroun a mis en place des \textbf{centres de déradicalisation et de réinsertion} pour les ex-membres repentis de Boko Haram / ISWAP, notamment dans l'Extrême-Nord (Mérina à Maroua, Mokolo, Kolofata). Ces centres accueillent d'anciens combattants (souvent recrutés de force, enrôlés enfants, ou trompés par l'idéologie), leur offrent un encadrement psychologique (prise en charge du trauma), un enseignement religieux authentique (avec des théologiens musulmans et chrétiens agréés) et une formation professionnelle certifiante (couture, menuiserie, mécanique, agriculture, coiffure, informatique) afin de les aider à redevenir des citoyens utiles et pacifiques. Le taux de récidive observé reste faible (< 5\%), validant l'approche « justice transitionnelle » combinant vérité, réparation et réconciliation. Cette approche montre qu'on peut combattre l'extrémisme autrement que par la seule répression.
\end{savaistu}

\courssub{Le rôle du citoyen dans la lutte contre l'extrémisme}

Chaque citoyen, y compris l'élève que tu es, a un rôle à jouer. La \cle{sécurité humaine} est l'affaire de tous (Constitution, Préambule : « Tout citoyen a le devoir de défendre la Patrie et de protéger la paix ») :

\begin{enumerate}
\item \textbf{La vigilance citoyenne :} être attentif aux comportements suspects (personnes inconnues rôdant, colis abandonnés, véhicules inhabituels), aux objets suspects (colis, sacs, véhicules garés longtemps) dans son quartier, son école, son atelier, les lieux de culte, les marchés.
\item \textbf{La dénonciation responsable :} signaler aux autorités compétentes (chef de quartier, chef de village, gendarmerie, police, numéro vert 1501 ou 117) toute tentative de recrutement, toute menace, tout discours de haine, sans se faire justice soi-même (pas de lynchage, pas de vengeance).
\item \textbf{Le refus de la manipulation et l'esprit critique :} vérifier les informations avant de les partager (lutte contre les \emph{fake news}), se méfier des discours de haine et des appels à la violence sur les réseaux sociaux (WhatsApp, Facebook, TikTok), développer son esprit critique face aux prédicateurs de la violence.
\item \textbf{La promotion active de la tolérance et du vivre ensemble :} respecter les différences de religion, d'ethnie, d'opinion, de langue ; refuser les rumeurs et les amalgames qui montent les communautés les unes contre les autres ; participer aux activités intercommunautaires (sport, culture, nettoyage, fêtes).
\item \textbf{L'engagement dans la vie locale :} adhérer aux associations de jeunesse, aux clubs UNESCO, aux comités de vigilance (si majeur et apte), aux conseils municipaux des jeunes ; être un relais d'information fiable.
\end{enumerate}

\begin{methode}[Analyser un document sur l'extrémisme violent]
Face à un article de presse, un témoignage, un rapport d'ONG, une image ou une vidéo sur l'extrémisme violent :
\begin{enumerate}
\item \textbf{Identifier la source et sa fiabilité :} qui parle ? (journaliste accrédité, témoin direct, autorité administrative, ONG reconnue, réseau social anonyme) Le document est-il daté, sourcé, recoupable ?
\item \textbf{Dégager les faits rapportés :} que s'est-il passé exactement ? Où ? Quand ? Qui sont les victimes ? Qui sont les auteurs présumés ? Quel est le bilan (humain, matériel) ?
\item \textbf{Repérer les causes explicites ou implicites :} pauvreté, chômage, manipulation religieuse/idéologique, sentiment d'abandon, vengeance, enrôlement forcé, appât du gain ?
\item \textbf{Relever les réponses apportées ou proposées :} sécuritaires (militaires, police, FMM, CV), judiciaires (arrestations, procès), préventives (éducation, développement, dialogue), réinsertion (centres, formation) ?
\item \textbf{Conclure avec esprit critique :} le document est-il équilibré ? Évite-t-il les amalgames (religion = terrorisme, ethnie = complicité) ? Propose-t-il des pistes durables ? Quelle leçon de citoyenneté en tirer ?
\end{enumerate}
\end{methode}

\begin{exoresolu}[Étude de cas : le témoignage d'Amadou]
Énoncé. Amadou, 19 ans, habitant du Mayo-Sava, raconte : « Après la mort de mon père, je ne pouvais plus aller à l'école. Un homme du village voisin m'a promis de l'argent et m'a dit que Dieu voulait que je combatte les "injustes". Un ancien du village, qui avait refusé de les suivre, m'a ouvert les yeux. J'ai alerté le comité de vigilance. »
1) Relever dans ce témoignage deux causes de l'embrigadement des jeunes.
2) Identifier deux acteurs de la lutte contre l'extrémisme qui apparaissent dans le récit.
3) Proposer une mesure de prévention adaptée à la situation d'Amadou.
\tcblower
Solution détaillée.
1) \textbf{Causes de l'embrigadement} : (a) la \textbf{pauvreté et la déscolarisation} (Amadou ne peut plus aller à l'école après la mort de son père, il perd son filet de protection social et devient vulnérable) ; (b) la \textbf{manipulation religieuse et l'appât financier} (le recruteur prétend que « Dieu veut » le combat, tout en promettant de l'argent : mélange d'endoctrinement idéologique et d'exploitation de la précarité économique).
2) \textbf{Acteurs de la lutte} : (a) \textbf{l'ancien du village}, figure d'autorité morale, citoyen éclairé qui déconstruit le discours extrémiste par la parole, l'exemple et le mentorat (rôle de la société civile) ; (b) \textbf{le comité de vigilance}, structure de participation citoyenne à la sécurité de proximité, alertée par Amadou (interface population-forces de défense).
3) \textbf{Mesure de prévention} : faciliter la \textbf{rescolarisation ou la formation professionnelle qualifiante} des jeunes orphelins, déscolarisés et vulnérables (bourses d'État, cantines scolaires, centres de formation professionnelle ruraux, apprentissage en atelier artisanat subventionné), car un jeune occupé, formé et inséré socialement est beaucoup moins vulnérable aux recruteurs. On peut aussi citer le renforcement du dialogue interreligieux au niveau local pour corriger les faux enseignements et protéger les jeunes de la manipulation.
\end{exoresolu}

\begin{attention}[Pièges à éviter dans une copie d'examen (Probatoire / Baccalauréat Technique)]
\begin{itemize}
\item Ne pas écrire « Boko Haram = l'islam » ou « Les musulmans = terroristes » : c'est un \textbf{amalgame faux et dangereux} (voir la boîte Attention plus haut). Préciser : « Boko Haram / ISWAP est un groupe terroriste qui \emph{se réclame} de l'islam mais en donne une lecture déviée, rejetée par l'immense majorité des musulmans camerounais. »
\item Ne pas limiter la lutte à l'armée : une bonne copie distingue toujours la \textbf{réponse sécuritaire} (militaire, FMM, comités de vigilance, cadre juridique) de la \textbf{réponse préventive et structurelle} (éducation, développement, déradicalisation, réinsertion, dialogue interreligieux, EMI).
\item Ne pas oublier le \textbf{rôle du citoyen} : vigilance, dénonciation, esprit critique, tolérance, engagement associatif — c'est souvent ce qui est demandé dans la dernière question d'un sujet d'ECM (compétence « Participer à la vie citoyenne »).
\item Ne pas confondre \textbf{causes profondes} (structurelles : pauvreté, exclusion, gouvernance) et \textbf{facteurs déclenchants} (proches : recruteur, promesses, événement traumatique).
\item Ne pas présenter la \textbf{déradicalisation} comme une « récompense » pour les terroristes : c'est une mesure de \textbf{sécurité publique} (réduire la menace) et de \textbf{justice restaurative} (réparer le lien social).
\end{itemize}
\end{attention}

\begin{exercice}[À toi de jouer]
1) Définis l'extrémisme violent en une phrase complète, en employant les mots « idéologie », « violence » et « différence ».
2) Cite trois conséquences des exactions de Boko Haram / ISWAP dans l'Extrême-Nord et explique en une phrase comment chacune menace la paix et le vivre ensemble.
3) Ton petit frère te dit : « Tous les membres de cette religion sont des terroristes. » Rédige un court paragraphe (5 à 8 lignes) pour lui répondre avec des arguments juridiques, factuels et moraux.
4) En t'appuyant sur la carte et le diagramme, explique pourquoi la coopération régionale (FMM) est indispensable pour lutter contre ce groupe terroriste.
\end{exercice}

\begin{corrige}[Exercice — À toi de jouer]
1) L'extrémisme violent est une \textbf{idéologie} qui justifie l'usage de la \textbf{violence} pour imposer ses vues et rejeter toute \textbf{différence} (religieuse, ethnique, politique, culturelle).
2) Exemples de conséquences :
   \begin{itemize}
   \item Les \textbf{morts, blessés et déplacés internes} : ils détruisent des familles, vident les villages, créent du ressentiment et brisent la cohésion sociale, fondement de la paix.
   \item Les \textbf{fermetures d'écoles et la déscolarisation massive} : elles privent les enfants d'instruction, fabriquent une génération vulnérable au recrutement et compromettent l'avenir du pays.
   \item Le \textbf{ralentissement économique et la destruction des moyens d'existence} (marchés, agriculture, élevage) : ils aggravent la pauvreté, donc le terreau de l'extrémisme, et entravent le développement national.
   \end{itemize}
   Chacune affaiblit la confiance entre communautés et entre citoyens et État, base du vivre ensemble.
3) Réponse attendue : Le raisonnement de ton frère est un \textbf{amalgame} (confusion entre un tout et une partie). Juridiquement, la \textbf{Constitution camerounaise (Préambule, Art. 1er)} garantit la liberté de religion et l'égalité de tous les citoyens devant la loi, sans distinction de croyance. Factuellement, les terroristes \textbf{manipulent} la religion en déformant ses textes ; la grande majorité des croyants rejette la violence et vit en paix avec les autres. D'ailleurs, dans l'Extrême-Nord, de nombreuses victimes de Boko Haram / ISWAP sont elles-mêmes des fidèles de cette religion (musulmans modérés, imams opposés au groupe). Juger tout un groupe à partir d'une minorité violente, c'est \textbf{injuste}, \textbf{inexact} et \textbf{dangereux} : cela divise la société, stigmatise des innocents et fait le jeu des extrémistes qui cherchent justement à opposer les communautés.
4) La carte montre que les zones d'action de Boko Haram / ISWAP (Logone-et-Chari, Mayo-Sava, Mayo-Tsanaga) sont \textbf{transfrontalières} : elles longent la frontière nigériane et le lac Tchad, zone partagée par 4 pays (Cameroun, Nigeria, Tchad, Niger). Les terroristes traversent facilement les frontières pour attaquer, se replier, se ravitailler. Aucun pays seul ne peut contrôler cet espace vaste, poreux et complexe (îles du lac, zones humides, forêts). La \textbf{Force Multinationale Mixte (FMM)} permet des opérations conjointes, le partage de renseignements, la poursuite à chaud au-delà des frontières et la mutualisation des moyens (air, renseignement, logistique). Sans elle, le groupe se reconstituerait sans cesse depuis les pays voisins.
\end{corrige}

\begin{aretenir}[L'essentiel du dossier — Fiche de révision Probatoire]
\begin{itemize}
\item \textbf{Phénomène :} Extrémisme violent (Boko Haram / ISWAP) dans l'Extrême-Nord depuis 2014. Attaques, enlèvements, kamikazes, pillages, recrutement de jeunes.
\item \textbf{Causes (faisceau) :} Pauvreté + chômage + déscolarisation + manipulation religieuse/idéologique + sentiment d'abandon/absence État + porosité frontières.
\item \textbf{Conséquences :} Morts, déplacés (300 000+), écoles fermées, économie locale ruinée, traumatismes, tensions intercommunautaires.
\item \textbf{Réponse globale (approche holistique) :}
  \begin{itemize}
  \item \textbf{Sécuritaire :} Armée (BIR, etc.), FMM (5 pays), Comités de vigilance (loi 2019), Cadre juridique (Loi antiterroriste 2014/2017).
  \item \textbf{Préventive/Sociale :} Déradicalisation (centres), Réinsertion (formation, AGR), Éducation (réouverture écoles, ECM, EMI), Développement (PLANUT, routes, eau, emploi jeunes), Dialogue interreligieux.
  \end{itemize}
\item \textbf{Rôle du citoyen :} Vigilance, Dénonciation (numéros verts), Esprit critique (anti-fake news), Tolérance/Vivre ensemble, Engagement associatif.
\item \textbf{Principes directeurs :} Respect des droits humains, Distinction religion/terrorisme, Justice transitionnelle (vérité, réparation, réconciliation), Approche genre (filles/femmes victimes spécifiques).
\end{itemize}
\end{aretenir}

\coursec{Le vivre ensemble et le bilinguisme (Dossier 3)}

Après avoir analysé les mécanismes de l'extrémisme violent et les réponses institutionnelles pour y faire face (Dossier 2), il apparaît clairement que la répression seule ne saurait suffire. La prévention durable repose sur la construction active d'un socle commun : le \cle{vivre ensemble}. Au Cameroun, ce vivre ensemble trouve son expression juridique et symbolique la plus forte dans le \cle{bilinguisme officiel} (français-anglais) et le \cle{multiculturalisme}, garantis par la Constitution et pilotés par une institution dédiée : la Commission Nationale pour la Promotion du Bilinguisme et du Multiculturalisme (CNPBM). Cette section explore les fondements, les enjeux, les outils institutionnels et le rôle de chaque citoyen pour faire de la diversité une force et non une fracture.

\courssub{Le vivre ensemble : définition, fondements et enjeux}

\begin{definition}[Le vivre ensemble]
Le \textbf{vivre ensemble} est la capacité de personnes différentes (origines ethniques, langues, cultures, religions, opinions) à cohabiter harmonieusement dans un même espace politique, dans le respect mutuel de la dignité, des droits et des devoirs de chacun, en partageant un destin commun et des règles du jeu acceptées par tous. Ce n'est pas une simple tolérance passive, mais une \emph{construction active} quotidienne.
\end{definition}

\begin{savaistu}[Le socle constitutionnel]
Le préambule de la Constitution du 18 janvier 1996 affirme l'attachement du peuple camerounais au bilinguisme et à la diversité culturelle, et l'article 1er, alinéa 3 dispose : « Les langues officielles de la République du Cameroun sont le français et l'anglais. Elles ont le même statut. L'État assure la promotion du bilinguisme sur l'ensemble du territoire. » Cet article est la \cle{charte fondatrice} du vivre ensemble camerounais : il élève la diversité linguistique au rang de principe constitutionnel.
\end{savaistu}

\textbf{Les piliers du vivre ensemble au Cameroun.}\par
Le Cameroun est souvent décrit comme une « Afrique en miniature » du fait de sa diversité exceptionnelle. Cette diversité repose sur trois piliers majeurs :
\begin{itemize}
    \item \textbf{La diversité ethnique et culturelle :} plus de \textbf{250 groupes ethniques} (Bantous, Sémi-bantous, Soudanais, Peuls, etc.) cohabitent, chacun avec ses langues, coutumes, chefferies traditionnelles et modes de vie.
    \item \textbf{Le bilinguisme officiel français-anglais :} héritage de l'histoire coloniale (administration française et britannique sous mandat de la SDN puis tutelle de l'ONU), le pays a réuni en 1961 deux territoires aux cultures juridiques, éducatives et administratives distinctes (droit civil et common law, sous-systèmes éducatifs francophone et anglophone). Le bilinguisme est l'outil juridique de cette réunification.
    \item \textbf{Le pluralisme religieux :} christianisme (catholicisme, protestantisme, églises de réveil), islam (majoritaire dans le Grand Nord) et religions traditionnelles africaines coexistent généralement dans la paix, avec de nombreux mariages mixtes et fêtes partagées.
\end{itemize}

\begin{popfigure}
\begin{tikzpicture}[scale=0.9, every node/.style={font=\small}]
    % Central hub
    \node[circle, draw=popdark, thick, fill=popblueL, text width=2.2cm, align=center, minimum height=1.8cm] (center) at (0,0) {\textbf{VIVRE ENSEMBLE}\\(Cameroun)};
    % Pillars data: angle/label/color
    \foreach \angle/\titre/\couleur in {
        90/Tolérance/popgreen,
        30/Justice/popblue,
        -30/Dialogue/poporange,
        -90/Bilinguisme/poppurple,
        -150/Respect/poppink,
        150/Solidarité/popgold} {
        % Spoke line
        \draw[thick, \couleur!70!popdark] (center) -- (\angle:3.5cm);
        % Outer node
        \node[draw=\couleur!70!popdark, thick, fill=\couleur!20, rounded corners=4pt, text width=2cm, align=center, minimum height=1.2cm] at (\angle:4.5cm) {\textbf{\titre}};
        % Connection decoration
        \fill[\couleur!50] (\angle:3.5cm) circle (3pt);
    }
    % Arrows between pillars (cycle)
    \foreach \angle/\nextangle in {90/30, 30/-30, -30/-90, -90/-150, -150/150, 150/90} {
        \draw[->, >=Stealth, popmuted, thick, bend left=15] (\angle:4.5cm) to (\nextangle:4.5cm);
    }
\end{tikzpicture}
\legende{Schéma en roue : six valeurs interdépendantes du vivre ensemble au Cameroun. La flèche circulaire indique que chaque valeur renforce les autres (ex. : le bilinguisme favorise le dialogue, qui renforce la tolérance).}
\end{popfigure}

\textbf{L'importance stratégique du vivre ensemble.}\par
Ce n'est pas un luxe moral, c'est une \cle{nécessité vitale} pour l'État :
\begin{enumerate}
    \item \textbf{Cohésion sociale et paix intérieure :} prévenir les conflits intercommunautaires (ex. crise dans les régions du Nord-Ouest et du Sud-Ouest, tensions entre agriculteurs et éleveurs dans l'Adamaoua et le Nord).
    \item \textbf{Unité nationale et intégrité territoriale :} le vivre ensemble est le ciment de l'État unitaire décentralisé. Sans lui, le risque de fragmentation du territoire grandit.
    \item \textbf{Développement économique et social :} la paix est le premier investissement. Les investisseurs fuient l'instabilité ; la main-d'œuvre qualifiée émigre. Le vivre ensemble permet la libre circulation des personnes, des biens et des capitaux sur tout le territoire.
    \item \textbf{Rayonnement international :} le Cameroun est souvent cité comme un pôle de stabilité relative en Afrique centrale. Son rôle dans les organisations régionales (CEEAC, CEMAC, Commission du Bassin du Lac Tchad) repose aussi sur sa crédibilité interne.
\end{enumerate}

\textbf{Les obstacles majeurs (diagnostic lucide).}\par
Pour soigner, il faut nommer le mal. Les principaux freins identifiés par les rapports publics et la société civile sont :
\begin{itemize}
    \item \textbf{Le tribalisme et le népotisme :} préférence accordée à l'appartenance ethnique ou clanique au détriment du mérite et de l'intérêt général (recrutement, marchés publics, admission aux concours).
    \item \textbf{La discrimination linguistique :} sentiment d'exclusion de certains citoyens (documents administratifs rédigés dans une seule langue officielle, difficulté d'accès à la justice ou à l'école dans la langue de son choix).
    \item \textbf{Les discours de haine et les fausses informations :} propos déshumanisants sur les réseaux sociaux (Facebook, WhatsApp) ciblant une ethnie, une région ou une religion. Ils précèdent souvent les violences physiques.
    \item \textbf{Les frustrations régionales :} sentiment de marginalisation de certaines régions (infrastructures, représentation, application du statut spécial du Nord-Ouest et du Sud-Ouest).
\end{itemize}

\begin{attention}[Piège à éviter : confusion « bilinguisme individuel » et « bilinguisme institutionnel »]
Le bilinguisme officiel ne signifie \textbf{pas} que chaque Camerounais parle couramment français et anglais. C'est une obligation pour \textbf{l'État} (administration, justice, lois, signalisation, médias publics) de traiter les deux langues à égalité. Un citoyen a le \textbf{droit} de s'adresser à l'administration dans l'une ou l'autre langue officielle et d'obtenir réponse dans cette langue. C'est une distinction cruciale à l'examen : ne confondez pas \emph{droit du citoyen} et \emph{devoir individuel de bilinguisme parfait}.
\end{attention}

\courssub{La CNPBM : institution garante du bilinguisme et du multiculturalisme}

Face à la montée des tensions (revendications des avocats et des enseignants anglophones fin 2016, crise dans les régions du Nord-Ouest et du Sud-Ouest), l'État a créé une institution spécifique pour institutionnaliser le vivre ensemble.

\begin{exemplebox}[La CNPBM : naissance et mandat]
La \textbf{Commission Nationale pour la Promotion du Bilinguisme et du Multiculturalisme (CNPBM)} a été créée par le \textbf{décret n°2017/013 du 23 janvier 2017} portant organisation et fonctionnement, annoncée par le Président de la République, Paul Biya, dans son discours à la Nation du 31 décembre 2016. Elle est placée sous l'autorité du Président de la République. Son premier président est \textbf{Peter Mafany Musonge}, ancien Premier Ministre, personnalité originaire de la région du Sud-Ouest.
\end{exemplebox}

\begin{propriete}[Missions légales de la CNPBM (décret n°2017/013 du 23 janvier 2017)]
La CNPBM a pour missions principales de :
\begin{enumerate}
    \item \textbf{Promouvoir le bilinguisme :} veiller à l'égal usage du français et de l'anglais dans les services publics, les médias d'État, la signalisation, les concours et les actes officiels.
    \item \textbf{Promouvoir le multiculturalisme :} valoriser la diversité culturelle comme richesse nationale et lutter contre les discriminations culturelles.
    \item \textbf{Promouvoir le vivre ensemble :} sensibiliser, éduquer et diffuser les valeurs de tolérance, de dialogue et de paix.
    \item \textbf{Recevoir et traiter les plaintes :} toute personne physique ou morale peut saisir la CNPBM pour tout acte de \textbf{discrimination} fondé sur la langue, la culture ou l'origine. Elle instruit les dossiers et formule des recommandations aux autorités compétentes.
    \item \textbf{Proposer des mesures :} suggérer aux pouvoirs publics toute mesure de nature à renforcer le bilinguisme, le multiculturalisme et le vivre ensemble.
\end{enumerate}
\end{propriete}

\begin{methode}[Comment saisir la CNPBM ? (Démarche citoyenne)]
Si vous êtes victime ou témoin d'une discrimination linguistique ou culturelle (ex. : refus de service dans l'une des langues officielles à l'hôpital, concours sans épreuves dans les deux langues, signalisation publique monolingue) :
\begin{enumerate}
    \item \textbf{Identifier le fait :} date, lieu, service concerné, description précise de l'acte discriminatoire.
    \item \textbf{Rassembler les preuves :} photos, courriers, témoignages.
    \item \textbf{Saisir la Commission :} par lettre adressée au Président de la CNPBM (Yaoundé), ou par les voies de dépôt prévues par la Commission (secrétariat, antennes, canaux électroniques officiels).
    \item \textbf{Suivre le dossier :} la Commission accuse réception, instruit, peut entendre les parties, et rend un avis ou une recommandation aux autorités compétentes.
\end{enumerate}
\end{methode}

\begin{exoresolu}[Analyse d'une plainte type]
\textbf{Énoncé :} Un candidat anglophone à un concours national d'entrée dans une grande école publique reçoit sa convocation et l'intégralité des épreuves écrites uniquement en français, alors que la réglementation autorise la composition dans la langue officielle de son choix. Il saisit la CNPBM. Quelle sera la démarche de la Commission et sur quel fondement juridique ?
\tcblower
\textbf{Solution détaillée :}
\begin{enumerate}
    \item \textbf{Recevabilité :} la plainte est recevable (citoyen camerounais, acte émanant d'un service public, discrimination linguistique alléguée).
    \item \textbf{Fondement juridique :} article 1er, alinéa 3 de la Constitution (égalité de statut des deux langues officielles) ; loi n°98/004 du 14 avril 1998 d'orientation de l'éducation au Cameroun (qui consacre le bilinguisme dans l'enseignement) ; décret n°2017/013 (mission de recevoir les plaintes).
    \item \textbf{Instruction :} la CNPBM demande au ministère de tutelle et à l'établissement concerné des explications et les pièces du concours.
    \item \textbf{Médiation et recommandation :} si le fait est avéré, la CNPBM recommande par exemple : (a) la mise à disposition de sujets dans les deux langues officielles, (b) la correction des dysfonctionnements de traduction, (c) la mise en place d'un dispositif permanent de bilinguisme dans l'établissement.
    \item \textbf{Suivi :} la Commission suit l'exécution de ses recommandations et rend compte de son activité dans son rapport annuel.
\end{enumerate}
\end{exoresolu}

\courssub{Les autres leviers institutionnels et symboliques}

La CNPBM ne travaille pas en vase clos. D'autres mécanismes juridiques et symboliques consolident l'édifice du vivre ensemble.

\textbf{Le statut spécial du Nord-Ouest et du Sud-Ouest (loi n°2019/024 du 24 décembre 2019).}\par
Issue des recommandations du Grand Dialogue National (30 septembre -- 4 octobre 2019), cette loi portant code général des collectivités territoriales décentralisées accorde un statut spécial aux régions du Nord-Ouest et du Sud-Ouest, fondé notamment sur leurs spécificités linguistiques, éducatives (sous-système anglophone) et juridiques (common law). Il prévoit des assemblées régionales dotées de compétences élargies et la prise en compte de ces spécificités dans l'organisation de la justice et de l'enseignement. C'est une application concrète du multiculturalisme institutionnel.

\textbf{L'enseignement des deux langues officielles.}\par
\begin{itemize}
    \item \textbf{Programmes scolaires :} l'anglais est enseigné dans le sous-système francophone et le français dans le sous-système anglophone, du primaire au secondaire.
    \item \textbf{Examens officiels :} le Probatoire et le Baccalauréat technique comportent des épreuves dans les deux langues officielles (ex. : \emph{English Language} pour les candidats francophones, \emph{Langue française} pour les candidats anglophones).
    \item \textbf{Établissements bilingues :} développement d'écoles et de lycées bilingues où l'alternance des langues d'enseignement est effective.
\end{itemize}

\textbf{Les rites républicains et culturels.}\par
\begin{itemize}
    \item \textbf{La Fête de la Jeunesse (11 février) :} défilés mixtes, chants en français et en anglais, tenues traditionnelles diverses.
    \item \textbf{La Fête Nationale (20 mai) :} célébration de l'unité nationale (référendum du 20 mai 1972 instituant la République unie du Cameroun), avec un message présidentiel diffusé dans les deux langues officielles.
    \item \textbf{Journée mondiale de la diversité culturelle (21 mai) :} animations culturelles, colloques et concours valorisant le bilinguisme et les cultures nationales.
\end{itemize}

\begin{popfigure}
\begin{tikzpicture}[scale=0.85, every node/.style={font=\footnotesize}]
    % Simplified Cameroon outline (polygon approximation)
    \draw[popdark, thick, fill=popcream] 
        (0,0) -- (1,0.2) -- (2,0.5) -- (3,0.3) -- (4,0) -- % South coast
        (4.5,1) -- (5,2) -- (5.2,3) -- (5,4) -- (4.5,5) -- % East border
        (4,6) -- (3.5,6.5) -- (2.5,7) -- (1.5,7.2) -- (0.5,7) -- % North
        (-0.5,6) -- (-1,5) -- (-1.2,4) -- (-1,3) -- (-0.8,2) -- (-0.5,1) -- cycle; % West border
    % Regions labels (simplified positions)
    \node[fill=popblue!30, draw=popblue, rounded corners=2pt, inner sep=2pt] at (1.5, 1.5) {SUD};
    \node[fill=popblue!30, draw=popblue, rounded corners=2pt, inner sep=2pt] at (0.5, 3) {CENTRE};
    \node[fill=popblue!30, draw=popblue, rounded corners=2pt, inner sep=2pt] at (2.5, 3.5) {EST};
    \node[fill=popblue!30, draw=popblue, rounded corners=2pt, inner sep=2pt] at (3.5, 2.5) {LITTORAL};
    \node[fill=popblue!30, draw=popblue, rounded corners=2pt, inner sep=2pt] at (2.8, 5) {ADAMAOUA};
    \node[fill=popblue!30, draw=popblue, rounded corners=2pt, inner sep=2pt] at (1.2, 5.5) {OUEST};
    \node[fill=popblue!30, draw=popblue, rounded corners=2pt, inner sep=2pt] at (-0.2, 4.5) {NORD};
    \node[fill=popblue!30, draw=popblue, rounded corners=2pt, inner sep=2pt] at (0.8, 6.5) {EXTRÊME-NORD};
    % Anglophone zones (2 regions) - Highlighted
    \node[fill=poporange!40, draw=poporange, thick, rounded corners=3pt, inner sep=3pt, font=\bfseries, align=center] at (-0.5, 2.5){NORD-OUEST\\(anglophone)};
    \node[fill=poporange!40, draw=poporange, thick, rounded corners=3pt, inner sep=3pt, font=\bfseries, align=center] at (3.8, 1.2){SUD-OUEST\\(anglophone)};
    % Capital
    \node[circle, fill=popdark, inner sep=1.5pt] (ya) at (0.8, 2.8) {};
    \node[anchor=west] at (ya) {Yaoundé};
    % Legend
    \begin{scope}[shift={(6.5, 5)}]
        \fill[popblue!30] (0,0) rectangle (0.5,0.3); \node[anchor=west] at (0.6,0.15) {Régions majoritairement francophones (8)};
        \fill[poporange!40] (0,0.5) rectangle (0.5,0.8); \node[anchor=west] at (0.6,0.65) {Régions anglophones (Nord-Ouest, Sud-Ouest)};
        \draw[popdark] (0,1.1) rectangle (0.5,1.4); \node[anchor=west] at (0.6,1.25) {Frontière internationale};
    \end{scope}
    % North Arrow
    \draw[->, thick, popdark] (6.7,0.8) -- (6.7,1.4) node[above] {N};
\end{tikzpicture}
\legende{Carte linguistique très simplifiée du Cameroun (schéma indicatif, sans précision cartographique). Les 8 régions majoritairement francophones (bleu clair) et les 2 régions anglophones (orange : Nord-Ouest et Sud-Ouest) illustrent la base territoriale du bilinguisme officiel. La capitale Yaoundé est en zone francophone.}
\end{popfigure}

\courssub{Le rôle du citoyen et de l'élève : acteurs du quotidien}

Les institutions (CNPBM, lois, statuts) posent le cadre. Mais le vivre ensemble se vit (ou se défait) dans les interactions quotidiennes : à l'atelier, en classe, au marché, dans le bus, sur WhatsApp.

\begin{aretenir}[Le citoyen artisan de paix : 5 engagements concrets]
\begin{enumerate}
    \item \textbf{Pratiquer la tolérance active :} accepter que l'autre soit différent (nom, accent, tenue, foi, opinion) sans le juger inférieur. Refuser les blagues et les généralisations ethniques ou régionales.
    \item \textbf{Apprendre l'autre langue officielle :} faire l'effort d'acquérir des bases fonctionnelles dans la langue de l'autre (anglais pour les francophones, français pour les anglophones). C'est le premier signe de respect.
    \item \textbf{Refuser les préjugés et les stéréotypes :} vérifier l'information avant de partager un message haineux sur les réseaux sociaux. Appliquer la « règle des 3 questions » : est-ce vrai ? est-ce utile ? est-ce bienveillant ?
    \item \textbf{Participer à la vie communautaire :} s'investir dans les associations de quartier, les comités de développement, les clubs ECM, les mouvements de jeunesse. Le « faire ensemble » crée le lien social.
    \item \textbf{Saisir les mécanismes de recours :} ne pas se taire face à une discrimination. Utiliser la CNPBM, les services de l'État compétents ou les organisations de défense des droits humains. La plainte est un acte citoyen, pas une délation.
\end{enumerate}
\end{aretenir}

\begin{methode}[Réussir une dissertation ou une question de synthèse au Probatoire ECM (thème « Vivre ensemble »)]
\textbf{Sujet type :} « Le bilinguisme est-il un atout ou un obstacle pour l'unité nationale au Cameroun ? » ou « Comment l'élève peut-il contribuer au vivre ensemble dans son établissement ? »
\textbf{Démarche en 4 étapes (méthode « D.I.A.R. ») :}
\begin{enumerate}
    \item \textbf{Définir et délimiter :} définir les termes clés (bilinguisme officiel, unité nationale, vivre ensemble). Poser le paradoxe : la diversité est à la fois source de richesse et source possible de tension.
    \item \textbf{Illustrer par des faits précis (le « socle ») :} mobiliser obligatoirement des dates (1961 réunification, 1996 Constitution, 2017 CNPBM, 2019 statut spécial), des institutions (CNPBM, Assemblée nationale), des textes (art. 1er al. 3 de la Constitution, décret n°2017/013, loi n°2019/024) et des exemples concrets (concours bilingues, signalisation, plainte à la CNPBM).
    \item \textbf{Analyser et argumenter :} structurer en 2 ou 3 parties logiques (ex. : I. Un cadre juridique favorable ; II. Des défis persistants de mise en œuvre ; III. Les leviers d'action de l'élève et du citoyen). Utiliser les connecteurs : \emph{cependant, par conséquent, en outre, notamment}.
    \item \textbf{Répondre et conclure :} donner une réponse tranchée à la question posée, puis ouvrir (ex. : « Le bilinguisme n'est pas un fardeau, c'est un atout du Cameroun dans l'intégration régionale africaine »).
\end{enumerate}
\end{methode}

\begin{exercice}[Entraînement au Probatoire -- Durée 30 min]
\textbf{Sujet :} « Au Cameroun, la Commission Nationale pour la Promotion du Bilinguisme et du Multiculturalisme (CNPBM) est l'institution clé du vivre ensemble. » 
En vous appuyant sur des exemples précis, montrez comment la CNPBM contribue à la consolidation de l'unité nationale.
\textbf{Indications pour le plan :}
\begin{enumerate}
    \item Introduction : accroche (crise dans le Nord-Ouest et le Sud-Ouest), définition de la CNPBM, thèse (instrument de régulation et de promotion).
    \item I. Un cadre juridique et institutionnel robuste (création en 2017, missions fixées par le décret, composition pluraliste).
    \item II. Une action concrète sur deux registres : 
        \begin{itemize}
            \item \textbf{Promotion proactive :} rapports annuels, campagnes de sensibilisation, avis sur les projets de textes.
            \item \textbf{Régulation réactive :} traitement des plaintes (concours, signalisation, administration), recommandations aux autorités.
        \end{itemize}
    \item III. Les limites et défis (moyens, pouvoir de recommandation non contraignant, méconnaissance du public) et pistes d'amélioration.
    \item Conclusion : bilan nuancé, rôle complémentaire du citoyen.
\end{enumerate}
\end{exercice}

\begin{corrige}[Exercice n°1 -- Éléments de correction attendus]
\textbf{Introduction :} le Cameroun, « Afrique en miniature », gère plus de 250 groupes ethniques et 2 langues officielles. La crise anglophone (depuis 2016) a révélé la fragilité du vivre ensemble. La création de la CNPBM (décret n°2017/013) est une réponse institutionnelle forte. Thèse : la CNPBM est à la fois le \textbf{thermomètre et le régulateur} du bilinguisme et du multiculturalisme, indispensable mais perfectible.
\textbf{I. Un cadre robuste :} président Peter Mafany Musonge (personnalité originaire du Sud-Ouest) ; membres représentant les régions et la société civile ; missions : promotion, traitement des plaintes, propositions aux pouvoirs publics ; rattachement à la Présidence de la République.
\textbf{II. Une action concrète :}
\begin{itemize}
    \item \textit{Promotion :} rapports annuels d'activité, campagnes de sensibilisation au bilinguisme, avis sur les projets de textes relatifs aux langues officielles.
    \item \textit{Régulation :} plaintes types : sujets de concours monolingues, signalisation publique dans une seule langue, accueil des usagers dans les hôpitaux et préfectures. Recommandations : traduction des actes, formation des agents, correction des dysfonctionnements.
\end{itemize}
\textbf{III. Les limites :} pas de pouvoir de sanction directe (les recommandations ne sont pas juridiquement contraignantes) ; moyens limités pour couvrir tout le territoire ; méconnaissance de l'institution par une partie des populations ; délais de traitement parfois longs.
\textbf{Conclusion :} la CNPBM est une avancée majeure (institutionnalisation du dialogue). Mais le vivre ensemble ne se décrète pas, il se \textbf{pratique}. L'élève, futur technicien, doit être ambassadeur du bilinguisme dans son atelier ou son entreprise : respecter les deux langues officielles, refuser le tribalisme, saisir la CNPBM si ses droits sont violés. L'unité dans la diversité est une condition du développement national.
\end{corrige}

\begin{savaistu}[Le Cameroun, laboratoire du vivre ensemble pour l'Afrique ?]
Le Cameroun est l'un des rares pays africains à avoir constitutionnalisé l'égalité de statut de deux langues officielles héritées de la colonisation (français et anglais). L'expérience de la CNPBM est observée avec intérêt par les organisations internationales comme un modèle possible de \textbf{gouvernance de la diversité}. Si le pari réussit (application effective du statut spécial, bilinguisme réel dans l'administration), le Cameroun pourra valoriser son « savoir-faire » en gestion pacifique de la diversité linguistique et culturelle.
\end{savaistu}

\begin{attention}[Erreurs fréquentes en examen (Probatoire / Baccalauréat technique)]
\begin{itemize}
    \item \textbf{Confondre la CNPBM et la Commission des Droits de l'Homme (CDHC) :} la CDHC traite de tous les droits humains (détention, torture, élections). La CNPBM est \textbf{spécialisée} : bilinguisme, multiculturalisme, discriminations linguistiques et culturelles.
    \item \textbf{Oublier la date 2017 :} la création de la CNPBM est postérieure à la Constitution de 1996. C'est une réponse \textbf{conjoncturelle} à la crise de 2016. Ne pas écrire « la Constitution de 1996 a créé la CNPBM » (faux).
    \item \textbf{Réduire le multiculturalisme au folklore :} le multiculturalisme n'est pas « danser le Bikutsi et le Makossa ». C'est la \textbf{reconnaissance juridique} de la diversité (chefferies traditionnelles, droit coutumier, sous-systèmes éducatifs et juridiques distincts) dans l'unité républicaine.
    \item \textbf{Négliger le statut spécial (loi n°2019/024) :} c'est la traduction juridique la plus forte du multiculturalisme pour les régions du Nord-Ouest et du Sud-Ouest. L'ignorer montre une méconnaissance de l'actualité institutionnelle majeure.
\end{itemize}
\end{attention}

\newpage

\coursec{Activité d'intégration}

\coursec{Leçon N05 : La notion de paix — Activité d'intégration : « Mon établissement, territoire de paix »}

\courssub{Objectifs de l'activité}
\begin{itemize}
    \item Appliquer les notions de paix, de règlement pacifique des différends et de vivre ensemble à une situation concrète observée dans l'établissement ou le quartier.
    \item Mener une mini-enquête structurée, traiter des données qualitatives et quantitatives, et produire une représentation graphique (diagramme en barres).
    \item Identifier un conflit réel, analyser ses causes et proposer un mode de règlement pacifique adapté (négociation, médiation, arbitrage).
    \item Rédiger un plaidoyer argumenté d'une page, mobilisant au moins deux institutions étudiées (CNPBM, médiateur, comités de paix), et conclure par une position justifiée.
    \item Communiquer oralement les résultats et participer à un débat collectif respectueux.
\end{itemize}

\courssub{Situation}
Le Lycée Technique de Bafoussam, situé dans la région de l'Ouest, accueille 2\,450 élèves répartis en 48 classes (de la Seconde à la Terminale, toutes séries confondues : F1, F2, F3, STI, STPT, etc.). L'établissement est traversé par une diversité linguistique et culturelle marquée : environ 55\,\% des élèves sont francophones, 35\,\% anglophones et 10\,\% bilingues ou issus de minorités (Bamileke, Bassa, Haoussa, etc.). Depuis la rentrée 2025-2026, l'administration a recensé une hausse des tensions : 27 signalements de disputes verbales à connotation ethnique ou linguistique (dont 14 dans la cour de récréation, 8 aux abords des sanitaires, 5 en salle de permanence), 3 altercations physiques légères entre groupes d'élèves, et 12 plaintes pour discriminations à l'accès aux postes de délégués de classe ou aux responsabilités dans les clubs (journal scolaire, club UNESCO, association des jeunes entrepreneurs).

Parallèlement, le proviseur a mis en place un \emph{Comité de paix scolaire} (CPS) composé de 4 enseignants (dont le CPE), 6 élèves délégués (3 filles, 3 garçons), 2 parents d'élèves et le médiateur scolaire désigné par le MINESEC. La CNPBM (Commission nationale pour la promotion du bilinguisme et du multiculturalisme) a organisé en mars 2026 une journée de sensibilisation dans l'établissement, distribuant 300 affiches « Vivre ensemble : ma langue, ta langue, notre richesse » et 150 guides du « Règlement pacifique des différends en milieu scolaire ».

Dans ce contexte, votre classe de Première F3 (62 élèves) est chargée de réaliser l'activité « Mon établissement, territoire de paix » sur trois séances (2 h + 2 h + 1 h). L'objectif est de produire un diagnostic partagé et un plaidoyer écrit à destination du conseil d'administration et de la CNPBM délégation régionale.

\courssub{Consignes}
\begin{enumerate}
    \item \textbf{Constitution des groupes et préparation (15 min).} Former 6 groupes de 10-11 élèves. Chaque groupe désigne un enquêteur principal, un secrétaire, un analyste de données, un rédacteur et un porte-parole. L'enseignant remet le \emph{Questionnaire unique} (annexe 1) et la \emph{Grille d'analyse des conflits} (annexe 2).
    \item \textbf{Collecte des données (45 min).} Chaque groupe interroge au minimum 20 personnes (12 élèves de niveaux différents, 4 enseignants, 2 personnels administratifs, 2 parents ou riverains du quartier). Les réponses sont consignées sur papier ou tablette.
    \item \textbf{Traitement statistique et graphique (30 min).} À partir des réponses à la question Q3 (« Quel symbole représente le mieux la paix pour vous ? »), chaque groupe construit un diagramme en barres des symboles cités (drapeau national, hymne, sceau, colombe, main tendue, arbre de la paix, autre). Les effectifs sont reportés dans un tableau de fréquences.
    \item \textbf{Analyse d'un conflit réel (30 min).} Sur la base de la grille d'analyse, chaque groupe sélectionne \textbf{un} conflit observé (ex. : tension linguistique au club de débat, dispute pour l'usage du terrain de sport, discrimination à l'élection des délégués). Il en identifie la nature (intrapersonnel, interpersonnel, intergroupes), les acteurs, les causes profondes et les conséquences sur le climat scolaire.
    \item \textbf{Proposition de règlement pacifique (20 min).} Le groupe choisit le mode le plus adapté parmi : \emph{négociation directe}, \emph{médiation par un pair (délégué/élève médiateur)}, \emph{médiation par un adulte (enseignant, CPE, médiateur scolaire)}, \emph{arbitrage disciplinaire (conseil de discipline)}. Il justifie son choix en 5 lignes maximum.
    \item \textbf{Rédaction du plaidoyer (45 min).} Chaque groupe rédige un texte d'une page (300-350 mots) intitulé « Pour le vivre ensemble dans mon lycée ». Le texte doit : (a) présenter le diagnostic (chiffres clés), (b) exposer le conflit choisi et la solution proposée, (c) mobiliser \textbf{au moins deux} institutions parmi : CNPBM, Médiateur scolaire, Comité de paix scolaire, Conseil de discipline, Administration, Associations de parents d'élèves, (d) se terminer par une position justifiée (engagement personnel ou collectif).
    \item \textbf{Restitution orale et débat (55 min).} Chaque porte-parole présente le plaidoyer (3 min). Suit un débat modéré par l'enseignant (20 min) sur la faisabilité des propositions. Un vote à main levée désigne le plaidoyer qui sera transmis au proviseur et à la CNPBM.
\end{enumerate}

\courssub{Ressources à mobiliser}
\begin{aretenir}[Notions clés du programme]
\textbf{Paix :} absence de violence, présence de justice, respect des droits, coopération. \textbf{Symboles de la paix au Cameroun :} drapeau (étoile jaune = unité), hymne (O Cameroun, berceau de nos ancêtres), sceau de la République (épaules de la justice, fasces), colombe, arbre de la paix (palmier), main tendue. \textbf{Règlement pacifique :} négociation (partenaires directs), médiation (tierce partie neutre), arbitrage (décision contraignante), conciliation. \textbf{Institutions :} CNPBM (Loi n°2019/014 du 19 déc. 2019) — mission : promouvoir le bilinguisme, le multiculturalisme, recevoir les plaintes ; Médiateur scolaire (Circ. n°00000019/MINESEC/SG/DRH/SDSS du 12/01/2021) — écoute, médiation, prévention ; Comité de paix scolaire — instance locale de prévention et de gestion des conflits. \textbf{Vivre ensemble :} reconnaissance de la diversité, égalité des chances, participation citoyenne.
\end{aretenir}

\begin{attention}[Lien avec les dossiers du programme]
Cette activité met en œuvre le \textbf{Dossier 3} (CNPBM, bilinguisme, multiculturalisme) et prépare le \textbf{Dossier 2} (lutte contre l'extrémisme violent) : les tensions identitaires non traitées (exclusion linguistique, discriminations) constituent des facteurs de vulnérabilité au radicalisme. Le médiateur scolaire et le CPS sont les premiers remparts institutionnels pour détecter et désamorcer ces signaux faibles.
\end{attention}

\courssub{Démarche et corrigé commenté}
\begin{exoresolu}[Réalisation type — Groupe 3 (exemple guidé)]
\textbf{Étape 1 — Collecte (extrait).} Le groupe 3 a interrogé 22 personnes. Réponses Q3 (symbole de paix) : Drapeau 6, Hymne 2, Sceau 1, Colombe 8, Main tendue 3, Arbre de la paix 2, Autre 0. Total = 22.

\textbf{Étape 2 — Tableau de fréquences et diagramme en barres.}
\begin{center}
\begin{tabularx}{\linewidth}{|l|X|c|}\hline
Symbole & Effectif & Fréquence (\%)\\\hline
Drapeau national & 6 & 27,3\\ \hline
Hymne national & 2 & 9,1\\ \hline
Sceau de la République & 1 & 4,5\\ \hline
Colombe & 8 & 36,4\\ \hline
Main tendue & 3 & 13,6\\ \hline
Arbre de la paix (palmier) & 2 & 9,1\\ \hline
Autre & 0 & 0\\ \hline
\textbf{Total} & \textbf{22} & \textbf{100}\\\hline
\end{tabularx}
\end{center}
\begin{popfigure}
\begin{tikzpicture}
\begin{axis}[
    ybar,
    bar width=0.6cm,
    width=\linewidth,
    height=7cm,
    symbolic x coords={Drapeau,Hymne,Sceau,Colombe,Main,Arbre,Autre},
    xtick=data,
    x tick label style={rotate=45,anchor=east,font=\small},
    ylabel={Effectif},
    ymin=0,
    ymax=10,
    ymajorgrids=true,
    grid style=dashed,
    nodes near coords,
    nodes near coords align={vertical},
    every node near coord/.append style={font=\small,popdark},
    tick label style={font=\small},
    label style={font=\small},
]
\addplot[popblue,fill=popblue!60] coordinates {
    (Drapeau,6) (Hymne,2) (Sceau,1) (Colombe,8) (Main,3) (Arbre,2) (Autre,0)
};
\end{axis}
\end{tikzpicture}
\legende{Diagramme en barres des symboles de la paix cités par les 22 personnes enquêtées (Groupe 3). La colombe arrive en tête (36,4\,\%), suivie du drapeau (27,3\,\%).}
\end{popfigure}

\textbf{Étape 3 — Analyse du conflit choisi.} Conflit : \emph{Tensions linguistiques au club de débat}. Nature : intergroupes (francophones vs anglophones). Acteurs : 12 élèves (7 francophones, 5 anglophones), 1 enseignant encadrant. Causes : impression des anglophones que les sujets imposés favorisent la culture francophone ; absence de traduction simultanée ; moqueries sur l'accent. Conséquences : 3 anglophones ont quitté le club, climat de méfiance, plainte au CPE.

\textbf{Étape 4 — Mode de règlement proposé.} \textbf{Médiation par le médiateur scolaire assisté de deux élèves délégués bilingues.} Justification : le conflit oppose deux groupes structurés autour de l'identité linguistique ; un tiers neutre institutionnel (médiateur scolaire) garantit l'équité ; l'appui de pairs bilingues facilite la compréhension mutuelle et la reformulation ; la médiation préserve la relation (contrairement à l'arbitrage disciplinaire qui sanctionne) ; elle s'inscrit dans la mission du médiateur (Circ. 2021) et les recommandations de la CNPBM.

\textbf{Étape 5 — Plaidoyer rédigé (320 mots).}
\emph{Pour le vivre ensemble dans mon lycée}

Notre enquête auprès de 22 membres de la communauté scolaire révèle que la colombe (36\,\%) et le drapeau (27\,\%) sont les symboles les plus associés à la paix. Pourtant, le diagnostic montre des fractures : 27 signalements à connotation ethnolinguistique depuis septembre, dont un conflit ouvert au club de débat où des élèves anglophones se sentent exclus par le choix des thèmes et l'absence de traduction.

Face à ce conflit intergroupes, nous proposons la médiation conduite par le médiateur scolaire, appuyé par deux délégués bilingues. Cette voie, préconisée par la Circulaire n°00000019/MINESEC, permet d'instaurer un dialogue structuré, de reformuler les griefs sans jugement et de co-construire une charte du club (thèmes alternés, traduction, sanction des moqueries). Le Comité de paix scolaire (CPS) assurera le suivi mensuel.

Nous sollicitons l'appui de la CNPBM (Loi 2019/014) pour : (1) former 10 élèves médiateurs pairs d'ici fin mars ; (2) doter le CPS d'un budget « micro-projets vivre ensemble » (500\,000 FCFA) pour des activités interculturelles (journée gastronomique, dictée bilingue, tournoi mixte). L'administration s'engage à inscrire le bilinguisme effectif dans le règlement intérieur (affichage bilingue, procès-verbaux bilingues).

En tant qu'élèves de Première technique, nous nous engageons à être des sentinelles du vivre ensemble : signaler toute discrimination, participer aux médiations, valoriser nos langues maternelles comme richesses partagées. La paix n'est pas l'absence de conflits, c'est la capacité à les transformer par le dialogue et le droit. Signe : les 11 membres du Groupe 3, délégués de la paix.

\textbf{Étape 6 — Restitution et débat.} Le porte-parole présente en 3 min. Le débat porte sur : faisabilité du budget (recherche de partenaires ONG/mairie), formation des médiateurs pairs (calendrier), rôle des parents. Vote : 38 voix pour le plaidoyer du Groupe 3, 14 pour le Groupe 1, 10 abstentions. Le texte final est transmis au proviseur et à la CNPBM délégation régionale de l'Ouest.
\end{exoresolu}

\courssub{Bilan de l'activité}
Cette activité d'intégration a permis de :
\begin{itemize}
    \item \textbf{Ancrer les savoirs dans le réel :} les élèves ont manipulé des données chiffrées (effectifs, pourcentages) issues de leur propre environnement, donnant sens aux symboles de la paix et aux institutions (CNPBM, médiateur, CPS).
    \item \textbf{Développer des compétences transversales :} enquête de terrain (communication, éthique), traitement statistique (tableau, diagramme), analyse de conflit (grille structurée), argumentation écrite (plaidoyer) et orale (restitution, débat).
    \item \textbf{Mobiliser le cadre juridique camerounais :} la référence explicite à la Loi 2019/014 (CNPBM) et à la Circulaire 2021 (médiateur scolaire) ancre l'exercice dans la citoyenneté active.
    \item \textbf{Favoriser l'appropriation collective :} le vote final et la transmission aux autorités transforment l'exercice scolaire en acte citoyen concret.
    \item \textbf{Révéler des leviers d'action :} la médiation par les pairs, soutenue par l'institution, apparaît comme le levier le plus pertinent pour les conflits identitaires en milieu scolaire technique.
\end{itemize}
L'enseignant veillera à exploiter les plaidoiries non retenues (affichage au CDI, publication sur le journal scolaire) pour valoriser tous les travaux.

\newpage

\coursec{Méthodes et exercices résolus}

\coursec{La notion de paix}

\courssub{La notion de paix}

\begin{definition}[Paix]
La \cle{paix} est un état de calme, de sécurité, de justice et de concorde dans les relations entre les personnes, les groupes et les États. On distingue deux dimensions complémentaires :
\begin{itemize}
\item la \cle{paix négative} : absence de guerre, de violence physique ou de conflit ouvert ;
\item la \cle{paix positive} : présence active de justice, de dialogue, de respect des droits, d'équité et d'harmonie sociale.
\end{itemize}
\end{definition}

\begin{aretenir}[Essentiel sur la paix]
\begin{itemize}
\item La paix ne se décrète pas, elle se construit au quotidien par le respect de l'autre et le respect de la loi.
\item Au Cameroun, la paix est une valeur constitutionnelle (préambule : « le peuple camerounais proclame son attachement aux libertés fondamentales... ») et un pilier de la devise nationale « Paix – Travail – Patrie ».
\item La paix positive est le fondement du \cle{vivre ensemble} dans un État multiculturel et bilingue.
\end{itemize}
\end{aretenir}

\begin{methode}[Analyser une situation de vie courante avec la notion de paix]
Pour mobiliser la notion de \cle{paix} dans une situation concrète (atelier, entreprise, quartier, réseaux sociaux), suivre la démarche suivante :
\begin{enumerate}
\item \textbf{Lire et identifier} : repérer dans l'énoncé les personnages, le lieu, le conflit ou la tension décrite.
\item \textbf{Définir la notion} : rappeler que la paix est l'absence de guerre et de violence (paix négative), mais aussi un état de \cle{calme}, de \cle{sécurité}, de \cle{justice} et de concorde (paix positive).
\item \textbf{Qualifier la situation} : dire si la situation traduit la paix ou son absence (violence physique, verbale, injustice, discrimination).
\item \textbf{Proposer une attitude} : indiquer le comportement citoyen attendu (dialogue, tolérance, respect de l'autre, recours à la médiation).
\item \textbf{Conclure} par une phrase qui relie la situation à la notion : « Cette situation montre que... ».
\end{enumerate}
\textbf{Réflexe} : toujours citer un élément précis de l'énoncé pour justifier sa réponse ; ne jamais répondre de façon générale sans lien avec le cas proposé.
\end{methode}

\begin{exoresolu}[Atelier de menuiserie à Douala]
\textbf{Énoncé.} Dans un atelier de menuiserie à Douala, deux apprentis se disputent violemment l'usage d'une machine. Le chef d'atelier intervient, les invite à s'asseoir, écoute chacun et propose un planning d'utilisation de la machine accepté par les deux. Le calme revient dans l'atelier.
\begin{enumerate}
\item Définis la paix.
\item Montre, à partir de deux éléments du texte, comment la paix a été rétablie dans l'atelier.
\end{enumerate}
\tcblower
\textbf{Solution détaillée.}
\begin{enumerate}
\item \textbf{Définition.} La paix est un état de calme, de sécurité et de concorde dans les relations entre les personnes, les groupes et les États. Elle désigne d'abord l'absence de guerre et de violence (paix négative), mais aussi la présence de la justice, du dialogue et de l'harmonie sociale (paix positive).
\item \textbf{Application à la situation.}
\begin{itemize}
\item Premier élément : « Le chef d'atelier intervient, les invite à s'asseoir, écoute chacun. » Cela montre que le conflit est réglé par le \cle{dialogue} et l'écoute, et non par la violence : c'est une condition essentielle de la paix.
\item Deuxième élément : « propose un planning d'utilisation accepté par les deux. » Cela montre qu'une solution \cle{juste} et acceptée par les deux parties remplace l'affrontement : la paix positive (justice, entente) est rétablie.
\end{itemize}
\item \textbf{Conclusion.} Cette situation montre que la paix ne consiste pas seulement à faire taire un conflit, mais à le résoudre par le dialogue et une solution équitable, ce qui permet le retour durable du calme et de la concorde dans le groupe.
\end{enumerate}
\end{exoresolu}

\courssub{TD4 : Les symboles de la paix au Cameroun (enquêtes)}

\begin{definition}[Symbole de la paix]
Un \cle{symbole de la paix} est un signe, un objet, un geste ou une devise reconnu collectivement comme représentant l'idéal de non-violence, d'harmonie et de concorde. Au Cameroun, les principaux symboles sont : la \cle{colombe}, le \cle{drapeau blanc}, le \cle{rameau d'olivier}, la \cle{flamme de la paix}, les monuments aux morts et la devise nationale \cle{« Paix – Travail – Patrie »}.
\end{definition}

\begin{methode}[Réaliser une enquête sur les symboles de la paix au Cameroun]
Pour mener une enquête (travail dirigé) et présenter ses résultats :
\begin{enumerate}
\item \textbf{Définir l'objectif} : identifier les symboles de la paix connus dans ton environnement.
\item \textbf{Construire l'outil} : rédiger un questionnaire court (5 à 8 questions fermées et ouvertes) ou un guide d'interview.
\item \textbf{Choisir l'échantillon} : interroger des personnes variées (élèves, enseignants, commerçants, anciens) et noter les réponses.
\item \textbf{Dépouiller les résultats} : compter les réponses, calculer des pourcentages : $\text{part} = \dfrac{\text{effectif de la catégorie}}{\text{effectif total}} \times 100$.
\item \textbf{Présenter et interpréter} : tableau ou diagramme, puis conclusion sur le symbole le plus connu et sa signification.
\end{enumerate}
\textbf{Réflexe} : un pourcentage se calcule toujours sur l'effectif total des personnes interrogées ; vérifier que la somme des pourcentages vaut 100 \%.
\end{methode}

\begin{exoresolu}[Enquête dans un lycée technique de Yaoundé]
\textbf{Énoncé.} Dans le cadre du TD4, un élève de Première technique interroge 50 personnes de son quartier sur le symbole de la paix qu'elles connaissent le mieux. Résultats : 20 citent la colombe, 15 le drapeau blanc, 10 la devise « Paix-Travail-Patrie » et 5 le rameau d'olivier.
\begin{enumerate}
\item Calcule le pourcentage de personnes ayant cité chaque symbole.
\item Quel symbole est le plus connu ? Propose une explication.
\end{enumerate}
\tcblower
\textbf{Solution détaillée.}
\begin{enumerate}
\item \textbf{Calcul des pourcentages} (effectif total : 50 personnes) :
\begin{align*}
\text{Colombe} &= \frac{20}{50} \times 100 = 40\ \%\\
\text{Drapeau blanc} &= \frac{15}{50} \times 100 = 30\ \%\\
\text{Devise « Paix-Travail-Patrie »} &= \frac{10}{50} \times 100 = 20\ \%\\
\text{Rameau d'olivier} &= \frac{5}{50} \times 100 = 10\ \%
\end{align*}
Vérification : $40 + 30 + 20 + 10 = 100\ \%$. Le dépouillement est correct.
\item \textbf{Interprétation.} Le symbole le plus connu est la \cle{colombe} (40 \%). Explication possible : la colombe est un symbole universel de la paix, très présent dans les dessins, les cérémonies scolaires, les affiches et les médias ; elle est donc plus familière aux populations que les autres symboles.
\end{enumerate}
\textbf{Conclusion.} L'enquête montre que la colombe est le symbole de la paix le plus connu du quartier (40 \%), devant le drapeau blanc (30 \%). Cela illustre l'importance des symboles dans l'éducation à la paix au Cameroun.
\end{exoresolu}

\begin{aretenir}[Bilan du TD4]
\begin{itemize}
\item La colombe est le symbole le plus partagé ; la devise nationale « Paix-Travail-Patrie » ancre la paix dans l'identité camerounaise.
\item Savoir mener une enquête (outil, échantillon, dépouillement, interprétation) est une compétence citoyenne évaluée au Probatoire.
\end{itemize}
\end{aretenir}

\courssub{Les stratégies de promotion de la paix}

\begin{definition}[Stratégie de promotion de la paix]
Une \cle{stratégie de promotion de la paix} est un ensemble coordonné de moyens et d'actions mis en œuvre par les acteurs étatiques, la société civile, l'école ou les individus pour prévenir les conflits, résoudre les tensions et consolider la concorde durable.
\end{definition}

\begin{methode}[Construire une réponse sur les stratégies de promotion de la paix]
Pour répondre à une question du type « Quelles stratégies pour promouvoir la paix ? » :
\begin{enumerate}
\item \textbf{Définir} la stratégie de promotion de la paix.
\item \textbf{Classer les stratégies} selon leurs acteurs :
\begin{itemize}
\item \emph{État} : justice équitable, bonne gouvernance, lutte contre la corruption, éducation civique, désarmement, dialogue social ;
\item \emph{Collectivités et société civile} : médiation communautaire, comités de paix, campagnes de sensibilisation ;
\item \emph{Individus et jeunes} : tolérance, esprit de dialogue, refus des rumeurs et des discours de haine (notamment sur les réseaux sociaux), sport et culture ;
\item \emph{École} : clubs de paix, éducation à la citoyenneté, règlement non violent des conflits.
\end{itemize}
\item \textbf{Illustrer} chaque stratégie par un exemple camerounais concret.
\item \textbf{Conclure} : la promotion de la paix est l'affaire de tous, à tous les niveaux.
\end{enumerate}
\textbf{Réflexe} : au Bac, une stratégie non illustrée par un exemple perd des points ; toujours accompagner chaque idée d'un cas concret.
\end{methode}

\begin{exoresolu}[Conflit entre deux villages]
\textbf{Énoncé.} Deux villages de la région du Nord-Ouest s'opposent au sujet d'une terre agricole. Les jeunes des deux camps menacent d'en venir aux mains. Le sous-préfet décide d'agir.
\begin{enumerate}
\item Cite trois stratégies que le sous-préfet peut utiliser pour promouvoir la paix entre ces villages.
\item Pour chaque stratégie, explique comment elle contribue à rétablir la paix.
\end{enumerate}
\tcblower
\textbf{Solution détaillée.}
\begin{enumerate}
\item \textbf{Stratégie 1 : la médiation (dialogue intercommunautaire).} Le sous-préfet réunit les chefs des deux villages autour d'une table pour écouter les revendications de chacun. Le dialogue permet de désamorcer les menaces, car chaque partie se sent entendue ; c'est le premier pas vers un accord.
\item \textbf{Stratégie 2 : le recours à la justice (décision équitable).} Les titres fonciers et les bornes de la terre sont examinés par les services compétents afin de trancher selon le droit. Une décision juste et fondée sur la loi éteint durablement le conflit, car aucune partie ne se sent lésée arbitrairement.
\item \textbf{Stratégie 3 : la sensibilisation et l'éducation à la paix.} Des campagnes de sensibilisation (rencontres sportives et culturelles entre les jeunes des deux villages, messages de paix) rapprochent les populations et détruisent les préjugés. Le sport et la culture créent des liens qui rendent la violence moins acceptable.
\end{enumerate}
\textbf{Conclusion.} La combinaison du dialogue, de la justice et de la sensibilisation montre que la promotion de la paix repose à la fois sur les autorités, sur le droit et sur l'engagement des populations elles-mêmes.
\end{exoresolu}

\courssub{Le règlement pacifique des différends}

\begin{definition}[Règlement pacifique des différends]
Le \cle{règlement pacifique des différends} est l'ensemble des moyens non violents utilisés pour résoudre un conflit entre des personnes, des groupes ou des États, sans recourir à la force ni aux armes. Il repose sur le libre consentement des parties (sauf pour la voie judiciaire) et le respect du droit.
\end{definition}

\begin{aretenir}[Les cinq modes classiques (ordre croissant d'intervention du tiers)]
\begin{center}
\begin{tabularx}{\linewidth}{|l|X|X|}\hline
\textbf{Mode} & \textbf{Rôle du tiers} & \textbf{Nature de la solution} \\ \hline
\cle{Négociation} & Aucun (parties seules) & Accord mutuel \\ \hline
\cle{Médiation} & Facilite le dialogue & Accord mutuel (proposé par les parties) \\ \hline
\cle{Conciliation} & Propose une solution & Proposition non contraignante \\ \hline
\cle{Arbitrage} & Tranche (arbitre choisi) & Décision \emph{contraignante} (s'impose aux parties) \\ \hline
\cle{Justice (Tribunaux)} & Juge (institutionnel) & Décision de droit \emph{contraignante} \\ \hline
\end{tabularx}
\end{center}
\end{aretenir}

\begin{methode}[Présenter les modes de règlement pacifique des différends]
Pour traiter une question sur le \cle{règlement pacifique des différends} :
\begin{enumerate}
\item \textbf{Définir} : c'est l'ensemble des moyens utilisés pour résoudre un conflit sans recourir à la violence ni à la force.
\item \textbf{Énumérer les modes} dans l'ordre logique croissant : négociation, médiation, conciliation, arbitrage, recours aux tribunaux.
\item \textbf{Comparer} : préciser qui propose la solution (les parties, le tiers, le juge) — c'est le critère de distinction le plus souvent demandé.
\item \textbf{Appliquer} au cas de l'énoncé et conclure sur l'avantage du règlement pacifique (préserve les relations, évite les violences, coûte moins cher).
\end{enumerate}
\textbf{Réflexe} : ne pas confondre médiation (le tiers facilite) et arbitrage (le tiers décide) ; c'est la confusion la plus fréquente.
\end{methode}

\begin{exoresolu}[Différend commercial à Bafoussam]
\textbf{Énoncé.} À Bafoussam, un fournisseur reproche à un commerçant de ne pas avoir payé une livraison de marchandises. Plutôt que d'en venir aux insultes, ils demandent d'abord au président du syndicat des commerçants de les aider à discuter. N'y parvenant pas, ils acceptent de soumettre le litige à un arbitre dont ils s'engagent à respecter la décision.
\begin{enumerate}
\item Définis le règlement pacifique des différends.
\item Identifie dans le texte les deux modes de règlement utilisés et justifie.
\item Cite deux avantages du règlement pacifique des différends.
\end{enumerate}
\tcblower
\textbf{Solution détaillée.}
\begin{enumerate}
\item \textbf{Définition.} Le règlement pacifique des différends est l'ensemble des moyens non violents (négociation, médiation, conciliation, arbitrage, justice) utilisés pour résoudre un conflit entre des personnes ou des groupes.
\item \textbf{Identification des modes.}
\begin{itemize}
\item « Ils demandent au président du syndicat de les aider à discuter » : il s'agit de la \cle{médiation}, car un tiers neutre facilite le dialogue sans imposer de solution.
\item « Ils soumettent le litige à un arbitre dont ils s'engagent à respecter la décision » : il s'agit de l'\cle{arbitrage}, car le tiers tranche le différend et sa décision s'impose aux deux parties.
\end{itemize}
\item \textbf{Avantages.}
\begin{itemize}
\item Il évite les violences et préserve les relations entre les parties (le fournisseur et le commerçant pourront continuer à travailler ensemble).
\item Il est généralement plus rapide et moins coûteux qu'un long procès ou qu'un conflit ouvert.
\end{itemize}
\end{enumerate}
\textbf{Conclusion.} Ce cas montre que le règlement pacifique des différends, en passant de la médiation à l'arbitrage, permet de résoudre un conflit commercial dans le calme, l'équité et le respect mutuel.
\end{exoresolu}

\courssub{Dossier 2 : La lutte contre l'extrémisme violent au Cameroun}

\begin{definition}[Extrémisme violent]
L'\cle{extrémisme violent} est une attitude qui consiste à rejeter radicalement toute opinion différente, tout compromis et toute modération, et à recourir à la violence, à la terreur ou aux armes pour imposer ses idées, souvent en instrumentalisant la religion, l'ethnie ou l'idéologie. Au Cameroun, le cas emblématique est l'action de \cle{Boko Haram} dans la région de l'Extrême-Nord (depuis 2014) et les groupes armés dans les régions du Nord-Ouest et du Sud-Ouest.
\end{definition}

\begin{methode}[Traiter un dossier sur l'extrémisme violent]
Pour analyser un document ou une situation liée à l'\cle{extrémisme violent} :
\begin{enumerate}
\item \textbf{Définir les notions} : distinguer extrémisme (rejet radical) et extrémisme violent (passage à l'acte violent/terroriste).
\item \textbf{Identifier les causes} (facteurs de vulnérabilité) : pauvreté, chômage des jeunes, ignorance/manque d'éducation, manipulation par les discours de haine, repli identitaire, porosité des frontières.
\item \textbf{Présenter les moyens de lutte} au Cameroun (approche globale) :
\begin{itemize}
\item \emph{Moyens sécuritaires} : action des forces de défense et de sécurité (BIR, armée, police, gendarmerie), vigilance des populations (comités de vigilance), loi antiterroriste (2014).
\item \emph{Moyens économiques et sociaux} : Plan d'urgence pour le développement du Nord (Plan d'urgence triennal), éducation, formation professionnelle, emploi des jeunes, développement des zones menacées (écoles, routes, santé).
\item \emph{Moyens éducatifs, religieux et culturels} : prédication de la tolérance par les leaders religieux modérés (imams, pasteurs, chefs traditionnels), éducation à la citoyenneté (ECM), dénonciation des discours de haine, promotion du vivre ensemble.
\end{itemize}
\item \textbf{Dégager le rôle du citoyen} : vigilance, refus de la manipulation, dénonciation (numéros verts), solidarité avec les victimes, éducation des enfants à la paix.
\item \textbf{Conclure} : la lutte contre l'extrémisme violent exige à la fois la sécurité, le développement et l'éducation (approche holistique).
\end{enumerate}
\end{methode}

\begin{exoresolu}[Un jeune tenté par un discours de haine]
\textbf{Énoncé.} Dans une localité de l'Extrême-Nord, un jeune chômeur est contacté sur les réseaux sociaux par un groupe qui lui promet de l'argent s'il rejoint un mouvement prônant la violence au nom d'une idéologie. Il en parle à son ancien professeur d'ECM.
\begin{enumerate}
\item Définis l'extrémisme violent.
\item Relève dans le texte deux facteurs qui favorisent l'extrémisme violent.
\item Propose deux actions pour empêcher ce jeune de basculer dans la violence.
\end{enumerate}
\tcblower
\textbf{Solution détaillée.}
\begin{enumerate}
\item \textbf{Définition.} L'extrémisme violent est une attitude qui consiste à rejeter radicalement toute opinion différente et à recourir à la violence, à la terreur ou aux armes pour imposer ses idées, souvent en manipulant la religion ou l'identité.
\item \textbf{Facteurs relevés dans le texte.}
\begin{itemize}
\item « Un jeune \emph{chômeur} » : la pauvreté et le chômage rendent les jeunes vulnérables aux promesses d'argent ; c'est un terreau du recrutement.
\item « Contacté \emph{sur les réseaux sociaux} » par un groupe qui « promet de l'argent » : la manipulation et la propagande par les discours de haine, amplifiées par les réseaux sociaux, facilitent l'endoctrinement.
\end{itemize}
\item \textbf{Actions de prévention.}
\begin{itemize}
\item \emph{Action éducative} : le dialogue avec un adulte de confiance (professeur, chef religieux modéré) pour déconstruire le discours de haine et rappeler les valeurs de tolérance et de respect de la vie humaine.
\item \emph{Action socio-économique et citoyenne} : orienter le jeune vers une formation professionnelle ou un emploi, et signaler le groupe aux autorités compétentes afin de couper le réseau de recrutement.
\end{itemize}
\end{enumerate}
\textbf{Conclusion.} Ce cas montre que la lutte contre l'extrémisme violent au Cameroun ne repose pas seulement sur les forces de sécurité : l'éducation, l'insertion des jeunes et la vigilance citoyenne sont des armes tout aussi décisives.
\end{exoresolu}

\courssub{Le vivre ensemble et le bilinguisme (Dossier 3)}

\begin{definition}[Vivre ensemble]
Le \cle{vivre ensemble} est la capacité de personnes de cultures, de langues, de religions et d'origines différentes à cohabiter dans le respect mutuel, la tolérance, la solidarité et le partage des mêmes règles de droit, au sein d'une même communauté nationale.
\end{definition}

\begin{definition}[Bilinguisme et Multiculturalisme au Cameroun]
Le \cle{bilinguisme} officiel (français / anglais) et le \cle{multiculturalisme} (plus de 250 groupes ethniques, diversité religieuse) sont des réalités constitutionnelles du Cameroun (Loi n°96-06 du 18 janvier 1996, Art. 1er). Ils font du Cameroun une « Afrique en miniature ».
\end{definition}

\begin{definition}[CNPBM]
La \cle{Commission nationale pour la promotion du bilinguisme et du multiculturalisme (CNPBM)}, créée par le Décret n°2017/013 du 23 janvier 2017, est l'institution chargée de veiller au respect de l'égalité de statut du français et de l'anglais, de promouvoir le multiculturalisme et de recevoir les plaintes relatives aux violations de ces principes.
\end{definition}

\begin{methode}[Rédiger une conclusion argumentée sur le vivre ensemble et le bilinguisme (Dossier 3)]
Pour rédiger une conclusion ou un paragraphe argumenté au Bac :
\begin{enumerate}
\item \textbf{Annoncer la position} en une phrase claire (ex. : « Le bilinguisme est un pilier du vivre ensemble au Cameroun »).
\item \textbf{Définir les notions clés} : vivre ensemble, bilinguisme, multiculturalisme, CNPBM.
\item \textbf{Développer 2 ou 3 arguments}, chacun suivi d'un exemple camerounais :
\begin{itemize}
\item Unité dans la diversité (plus de 250 groupes ethniques) ;
\item Communication entre anglophones et francophones (école, administration, marché) ;
\item Rôle de la CNPBM (garant institutionnel de l'égalité linguistique et culturelle).
\end{itemize}
\item \textbf{Rejeter l'objection} éventuelle (ex. : « la diversité divise ») en montrant qu'elle est une richesse si elle est gérée par le dialogue et le droit.
\item \textbf{Conclure} par une phrase d'ouverture ou d'appel (engagement du citoyen, de l'élève, du jeune technicien).
\end{enumerate}
\textbf{Plan type d'un paragraphe} : idée $\rightarrow$ explication $\rightarrow$ exemple camerounais $\rightarrow$ lien avec la paix.
\end{methode}

\begin{exoresolu}[Paragraphe argumenté : le bilinguisme, facteur de paix]
\textbf{Énoncé.} Rédige un paragraphe argumenté d'une quinzaine de lignes montrant que le bilinguisme et le multiculturalisme contribuent au vivre ensemble au Cameroun. Tu préciseras le rôle de la Commission nationale pour la promotion du bilinguisme et du multiculturalisme.
\tcblower
\textbf{Solution détaillée (corrigé type).}

Le bilinguisme et le multiculturalisme constituent des piliers essentiels du vivre ensemble au Cameroun. Le vivre ensemble désigne la capacité de personnes de langues, de cultures et de religions différentes à cohabiter dans le respect mutuel et la concorde. Or le Cameroun est un pays profondément divers : il compte deux langues officielles, le français et l'anglais, héritées de son histoire, et plus de 250 groupes ethniques aux traditions variées. Cette diversité, loin d'être une faiblesse, devient une richesse lorsqu'elle est gérée par le dialogue et le respect.

D'abord, le bilinguisme permet la communication entre les citoyens des régions anglophones et francophones : un élève de Bamenda qui apprend le français et un élève de Douala qui apprend l'anglais peuvent se comprendre, travailler ensemble et tisser des liens d'amitié, ce qui renforce l'unité nationale. Ensuite, le multiculturalisme invite chaque citoyen à respecter les cultures des autres : les festivals, les langues locales et les traditions de chaque région sont autant de richesses partagées qui nourrissent la fierté commune d'être camerounais. Enfin, l'État accompagne cet effort : la Commission nationale pour la promotion du bilinguisme et du multiculturalisme (CNPBM), créée en 2017, veille à l'égalité de statut du français et de l'anglais dans l'administration et à la valorisation des cultures, et contribue ainsi à apaiser les tensions linguistiques.

Certains objectent que la diversité divise ; mais c'est le mépris de l'autre qui divise, non la différence. En conclusion, le bilinguisme et le multiculturalisme, promus par l'école, l'État et chaque citoyen, font de la diversité camerounaise un ciment du vivre ensemble et de la paix nationale.
\end{exoresolu}

\courssub{Entraînement au Probatoire / Baccalauréat Technique}

\begin{exercice}[Application : Conflit de voisinage à Ngaoundéré]
\textbf{Énoncé.} À Ngaoundéré, deux voisins se disputent bruyamment pour une question de limite de parcelle. L'un menace de détruire la clôture de l'autre. Le chef de quartier est alerté. Il reçoit les deux parties séparément, puis les réunit. Il leur rappelle la loi sur la propriété foncière et leur propose de faire appel à un géomètre pour borner le terrain. Les deux voisins acceptent la solution.
\begin{enumerate}
\item Identifie le mode de règlement pacifique utilisé par le chef de quartier lorsqu'il « reçoit les deux parties séparément, puis les réunit » et justifie.
\item Identifie le mode envisagé pour la suite (« faire appel à un géomètre ») et justifie.
\item Pourquoi le recours à la violence (détruire la clôture) aurait-il été une mauvaise solution ? Donne deux raisons.
\end{enumerate}
\end{exercice}

\begin{corrige}[Exercice n°1 : Conflit de voisinage à Ngaoundéré]
\begin{enumerate}
\item \textbf{Médiation.} Le chef de quartier agit comme un tiers neutre qui facilite le dialogue entre les parties (« reçoit les deux parties séparément, puis les réunit ») sans imposer de solution. Il les aide à trouver un accord par elles-mêmes.
\item \textbf{Conciliation (ou expertise amiable).} Le chef de quartier « propose de faire appel à un géomètre pour borner le terrain ». Le tiers (le géomètre, mandaté par le chef de quartier ou choisi d'un commun accord) propose une solution technique (le bornage) basée sur la loi. Cette proposition n'est contraignante que si les parties l'acceptent (ce qu'elles font).
\item \textbf{Raisons de rejeter la violence :}
\begin{itemize}
\item La violence aggrave le conflit (risque de blessures, de mort, de représailles) et ne règle pas le fond du problème (la limite exacte).
\item Elle expose l'auteur à des poursuites pénales (destruction de bien, menaces, violences) et détruit durablement le lien social de voisinage.
\end{itemize}
\end{enumerate}
\end{corrige}

\begin{exercice}[Analyse de document : Extrait de discours]
\textbf{Énoncé.} « Le Cameroun est un et indivisible. Sa devise "Paix - Travail - Patrie" nous oblige à construire la nation dans la diversité. Le bilinguisme n'est pas un obstacle, c'est une chance. La CNPBM veille à ce que nul ne soit discriminé pour sa langue ou sa culture. »
\begin{enumerate}
\item Cite les trois termes de la devise nationale et explique le lien entre le premier terme (« Paix ») et le vivre ensemble.
\item Quel est le sigle de l'institution chargée de promouvoir le bilinguisme et le multiculturalisme ? Donne sa date de création.
\item Selon le texte, pourquoi le bilinguisme est-il une « chance » ? Donne un argument concret lié à l'emploi ou à l'éducation.
\end{enumerate}
\end{exercice}

\begin{corrige}[Exercice n°2 : Analyse de document]
\begin{enumerate}
\item \textbf{Devise :} Paix – Travail – Patrie. \textbf{Lien :} La « Paix » (premier terme) est la condition sine qua non du vivre ensemble. Sans paix (sécurité, absence de violence, justice), la cohabitation entre les divers groupes culturels et linguistiques est impossible. Le vivre ensemble \emph{est} la mise en œuvre concrète de la paix positive au quotidien.
\item \textbf{Institution :} CNPBM (Commission nationale pour la promotion du bilinguisme et du multiculturalisme). \textbf{Création :} Janvier 2017 (Décret n°2017/013 du 23 janvier 2017).
\item \textbf{Argument :} Le bilinguisme (français/anglais) ouvre l'accès à deux marchés du travail, deux systèmes éducatifs, deux espaces culturels et économiques (zone CEMAC/CEEAC et zone Commonwealth/CEDEAO). Un jeune bilingue a plus d'opportunités d'emploi (administration, entreprises internationales, ONG) et peut étudier dans les deux sous-systèmes universitaires.
\end{enumerate}
\end{corrige}

\begin{attention}[Les 5 erreurs qui coûtent des points au Bac]
\begin{enumerate}
\item \textbf{Définir la paix uniquement comme « l'absence de guerre ».} C'est oublier la paix positive (justice, dialogue, harmonie sociale). Une définition complète mentionne les deux dimensions.
\item \textbf{Répondre de façon générale sans citer l'énoncé.} Dans les questions d'application, il faut toujours relever un élément précis du texte (« le texte dit que... ») pour justifier : une réponse hors du cas proposé perd la moitié des points.
\item \textbf{Confondre médiation et arbitrage.} Dans la médiation, le tiers \emph{aide à dialoguer} sans décider ; dans l'arbitrage, le tiers \emph{tranche} et sa décision s'impose. Cette confusion est sanctionnée dans les questions sur le règlement pacifique des différends.
\item \textbf{Se tromper dans le calcul des pourcentages d'enquête.} Le pourcentage se calcule toujours sur l'effectif \emph{total} des personnes interrogées : $\text{part} = \dfrac{\text{effectif de la catégorie}}{\text{effectif total}} \times 100$. Vérifier que la somme fait 100 \%.
\item \textbf{Présenter la diversité culturelle comme un problème ou ignorer les institutions.} Au Bac, il faut montrer que la diversité est une richesse gérée par le dialogue, et citer correctement les institutions : la \cle{Commission nationale pour la promotion du bilinguisme et du multiculturalisme} (CNPBM), et non un nom inventé ou approximatif.
\end{enumerate}
\end{attention}

\newpage

\coursec{Exercices}

\courssub{Je vérifie mes connaissances}

\begin{exercice}[QCM : Concepts fondamentaux de la paix]
Parmi les propositions suivantes, une seule est exacte. Recopie le numéro de l'exercice et la lettre de la bonne réponse.
\begin{enumerate}
\item La paix positive se définit par :
\begin{itemize}
\item[A.] L'absence totale de conflits armés.
\item[B.] La présence de justice, d'équité et de respect des droits humains.
\item[C.] La signature d'un traité de non-agression entre deux États.
\item[D.] Le silence imposé par une autorité forte.
\end{itemize}
\item Le symbole officiel de la paix au Cameroun, représenté sur le drapeau national, est :
\begin{itemize}
\item[A.] La colombe blanche.
\item[B.] L'étoile jaune sur fond rouge.
\item[C.] La branche d'olivier.
\item[D.] Le signe « V » de la victoire.
\end{itemize}
\item La Commission Nationale pour la Promotion du Bilingualisme et du Multiculturalisme (CNPBM) a pour mission principale de :
\begin{itemize}
\item[A.] Gérer les conflits frontaliers.
\item[B.] Promouvoir le vivre-ensemble par le respect de la diversité linguistique et culturelle.
\item[C.] Former les forces de l'ordre au maintien de la paix.
\item[D.] Surveiller les élections.
\end{itemize}
\item Le règlement pacifique des différends exclut :
\begin{itemize}
\item[A.] La médiation.
\item[B.] L'arbitrage.
\item[C.] Le recours à la force armée.
\item[D.] La conciliation.
\end{itemize}
\end{enumerate}
\end{exercice}

\begin{exercice}[Vrai ou Faux justifié]
Indique si chaque affirmation est Vraie (V) ou Fausse (F) et justifie ta réponse en une phrase en t'appuyant sur le cours.
\begin{enumerate}
\item L'extrémisme violent se caractérise uniquement par des actes terroristes commis par des groupes étrangers.
\item Le « vivre ensemble » au Cameroun repose sur l'assimilation de toutes les cultures minoritaires à la culture majoritaire.
\item La médiation est une procédure où un tiers impose sa solution aux parties en conflit.
\item L'arrêt de la Cour Internationale de Justice (CIJ) de 2002 sur l'affaire de la péninsule de Bakassi est un exemple de règlement judiciaire d'un différend territorial.
\end{enumerate}
\end{exercice}

\begin{exercice}[Définitions et distinctions]
Donne une définition précise (2 à 3 lignes) pour chacun des termes suivants :
\begin{enumerate}
\item Paix négative vs Paix positive.
\item Extrémisme violent.
\item Vivre ensemble (dans le contexte camerounais).
\item Médiation (comme mode de règlement pacifique).
\end{enumerate}
\end{exercice}

\begin{exercice}[Calcul et interprétation de données statistiques]
Lors d'une enquête menée auprès de 200 élèves d'un lycée technique de Yaoundé sur « Le symbole qui représente le mieux la paix au Cameroun », les résultats sont les suivants :
\begin{center}
\begin{tabular}{|l|c|}
\hline
Symbole & Effectif \\
\hline
Étoile jaune du drapeau & 80 \\
Colombe blanche & 50 \\
Main serrée (poignée de main) & 40 \\
Arbre de la paix (arbre à palabres) & 30 \\
\hline
Total & 200 \\
\hline
\end{tabular}
\end{center}
\begin{enumerate}
\item Calcule le pourcentage de chaque modalité.
\item Quelle conclusion peux-tu tirer sur la perception des symboles de la paix par ces jeunes ?
\end{enumerate}
\end{exercice}

\courssub{Je m'entraîne}

\begin{exercice}[Analyse de document : L'arrêt Bakassi (CIJ, 2002)]
On donne l'extrait suivant : « La Cour constate que le Traité de 1913 et les accords subséquents confèrent la souveraineté sur la péninsule de Bakassi au Cameroun. Elle invite les parties à respecter cet arrêt et à assurer le retrait ordonné des forces d'administration et de police nigérianes. »
\begin{enumerate}
\item Identifie les deux parties en présence dans ce différend.
\item Quel mode de règlement pacifique a été utilisé ici ? Justifie.
\item Dégage deux leçons que l'Afrique peut tirer de ce règlement pour la consolidation de la paix sur le continent.
\end{enumerate}
\end{exercice}

\begin{exercice}[Étude de cas : Conflit intercommunautaire fictif]
Dans le village de \emph{Mbalmayo-Nord}, un conflit oppose les éleveurs peuls (transhumants) et les agriculteurs bétis (autochtones) à cause de la destruction des champs par le bétail. Le chef de village a tenté d'arbitrer, mais les deux camps rejettent sa décision, l'accusant de partialité. La tension monte : menaces, blocage du marché hebdomadaire.
\begin{enumerate}
\item Propose une démarche de médiation en 4 étapes structurées pour sortir de cette crise.
\item Pour chaque étape, précise l'acteur principal, l'action menée et l'objectif visé.
\item Explique pourquoi la médiation est plus adaptée que l'arbitrage imposé dans ce cas précis.
\end{enumerate}
\end{exercice}

\begin{exercice}[Stratégies de lutte contre l'extrémisme violent au Cameroun]
À partir de tes connaissances sur le Dossier 2, complète le tableau suivant en associant chaque stratégie à son niveau d'action et un exemple concret observé au Cameroun.
\begin{center}
\begin{tabularx}{\linewidth}{|X|X|X|}
\hline
\textbf{Stratégie} & \textbf{Niveau d'action (Étatique / Communautaire / Individuel)} & \textbf{Exemple concret au Cameroun} \\
\hline
Opérations militaires (ex: Force Multinationale Mixte) & & \\
\hline
Programmes de déradicalisation et réinsertion (centres de désengagement) & & \\
\hline
Promotion de l'éducation à la citoyenneté et au vivre-ensemble à l'école & & \\
\hline
Veille citoyenne et comités de vigilance locaux & & \\
\hline
\end{tabularx}
\end{center}
\end{exercice}

\begin{exercice}[Travail pratique noté : Diagramme circulaire et commentaire]
Reprends les données de l'exercice 4 (partie « Je vérifie mes connaissances »).
\begin{enumerate}
\item Construis, sur papier millimétré ou à l'aide d'un tableur, un diagramme circulaire représentant la répartition des réponses. Respecte les règles de construction (titre, légende, calcul des angles, couleurs distinctes).
\item Rédige un commentaire analytique de 5 lignes maximum interprétant ce graphique (tendance majeure, surprise, lien avec le cours sur les symboles officiels vs symboles universels).
\end{enumerate}
\end{exercice}

\courssub{Je me prépare au Bac}

\begin{exercice}[Sujet de dissertation guidée type Probatoire]
\textbf{Sujet :} « La paix est-elle seulement l'absence de guerre ? »
Rédige une dissertation structurée en respectant le plan suivant :
\begin{enumerate}
\item \textbf{Introduction} : Accroche, définition des termes (paix négative / positive), annonce du plan en deux parties.
\item \textbf{Première partie : La paix comme absence de violence directe (Paix négative).} Argumenter avec 2 exemples camerounais ou internationaux (ex: cessez-le-feu, fin de guerre ouverte).
\item \textbf{Deuxième partie : La paix comme construction structurelle (Paix positive).} Argumenter avec 2 exemples montrant justice sociale, droits humains, vivre-ensemble (ex: CNPBM, décentralisation, éducation).
\item \textbf{Conclusion} : Bilan synthétique, ouverture sur le rôle de la jeunesse technique dans la consolidation de la paix positive.
\end{enumerate}
\end{exercice}

\begin{exercice}[Sujet de synthèse type Baccalauréat Technique]
\textbf{Sujet :} « Les stratégies de lutte contre l'extrémisme violent au Cameroun sont-elles suffisantes ? Argumenter à partir des documents fournis. »

\textbf{Document 1 (Extrait de discours officiel, 2023) :} « Notre réponse est globale : militaire, humanitaire, éducative et idéologique. Les centres de désengagement de Mérina et de Mora ont accueilli plus de 2000 ex-combattants. »
\textbf{Document 2 (Rapport ONG locale, 2024) :} « Les zones du Lac Tchad et du Nord-Ouest/Sud-Ouest restent fragiles. La pauvreté, le chômage des jeunes et le sentiment d'injustice alimentent le recrutement. La réinsertion socio-économique peine à suivre. »
\textbf{Document 3 (Témoignage d'un chef traditionnel, Extrême-Nord) :} « L'armée nous protège, mais la paix vraie revient quand les jeunes ont du travail et que la justice passe pour tous, pas seulement pour les puissants. »

\begin{enumerate}
\item Dégage la thèse principale de chaque document.
\item Identifie deux forces et deux limites de la stratégie camerounaise actuelle à la lumière des trois documents.
\item En t'appuyant sur l'analyse des documents et tes connaissances personnelles, réponds à la question du sujet en développant un raisonnement nuancé en deux arguments principaux (un pour le « oui », un pour le « non / mais »).
\item Conclus en proposant une piste concrète d'amélioration pour la prévention de l'extrémisme violent dans ton environnement scolaire ou professionnel.
\end{enumerate}
\end{exercice}

\begin{exercice}[Analyse de cas complexe : Le rôle de la CNPBM dans le vivre-ensemble]
\textbf{Contexte :} Suite à des tensions nées de propos haineux sur les réseaux sociaux ciblant une communauté linguistique lors d'une compétition sportive universitaire, la CNPBM est saisie par le Premier Ministre.
\begin{enumerate}
\item Rappelle la base légale de la CNPBM (loi de création, année) et sa composition (présidence, membres).
\item Détaille les 3 missions constitutionnelles de la Commission en lien avec cette saisine.
\item Imagine le communiqué de presse que la CNPBM pourrait publier : rédige-le en 10 lignes maximum (ton institutionnel, rappel de la loi, appel au calme, annonce d'actions concrètes : enquête, audience publique, recommandations).
\item Explique en quoi l'action de la CNPBM dans ce cas illustre la prévention des conflits par le dialogue institutionnel plutôt que par la répression policière.
\end{enumerate}
\end{exercice}

\newpage

\coursec{Corrigés des exercices}

\coursec{Corrigés des exercices d'application — La notion de paix}

\begin{corrige}[Exercice 1 — QCM : Concepts fondamentaux de la paix]
\begin{enumerate}
\item \textbf{Réponse B.} La paix positive ne se limite pas à l'absence de guerre : elle suppose la présence de la \cle{justice}, de l'\cle{équité} et du respect des \cle{droits humains}. Les propositions A et C décrivent la paix négative ; la proposition D décrit un ordre imposé par la contrainte, qui n'est pas la paix.
\item \textbf{Réponse B.} Sur le drapeau camerounais, l'\cle{étoile jaune} (dite « étoile de l'unité »), placée sur la bande rouge centrale, symbolise l'unité nationale et la paix. La colombe et la branche d'olivier sont des symboles universels, non officiels au Cameroun.
\item \textbf{Réponse B.} La \cle{CNPBM}, créée par la loi n\textsuperscript{o} 2017/014 du 12 juillet 2017, a pour mission de promouvoir le \cle{vivre-ensemble} par le respect de la diversité linguistique (bilinguisme) et culturelle (multiculturalisme).
\item \textbf{Réponse C.} Le règlement pacifique des différends repose sur la négociation, la médiation, la conciliation, l'arbitrage et le recours judiciaire. Le \cle{recours à la force armée} est par définition exclu des modes pacifiques.
\end{enumerate}
\textbf{Points de vigilance :} ne pas confondre symboles universels de la paix (colombe, rameau d'olivier) et symbole national camerounais (étoile de l'unité) ; distinguer toujours paix négative et paix positive.
\end{corrige}

\begin{corrige}[Exercice 2 — Vrai ou Faux justifié]
\begin{enumerate}
\item \textbf{Faux.} L'extrémisme violent peut être commis par des groupes étrangers (ex. Boko Haram) mais aussi par des acteurs internes radicalisés ; il se caractérise par l'usage de la violence au nom d'une idéologie, quelle que soit l'origine des auteurs.
\item \textbf{Faux.} Le vivre-ensemble camerounais repose sur le respect et la valorisation de la diversité (plus de 250 groupes ethniques, deux langues officielles), et non sur l'assimilation des minorités à une culture dominante.
\item \textbf{Faux.} Dans la médiation, le tiers (le médiateur) \emph{facilite} le dialogue et aide les parties à trouver elles-mêmes un accord ; c'est l'arbitre qui tranche et impose une solution.
\item \textbf{Vrai.} L'arrêt de la CIJ du 10 octobre 2002, attribuant la souveraineté sur la péninsule de Bakassi au Cameroun, est un exemple de règlement \cle{judiciaire} pacifique d'un différend territorial entre le Cameroun et le Nigeria.
\end{enumerate}
\textbf{Points de vigilance :} la justification doit mobiliser une notion précise du cours (médiation $\neq$ arbitrage ; vivre-ensemble $\neq$ assimilation).
\end{corrige}

\begin{corrige}[Exercice 3 — Définitions et distinctions]
\begin{enumerate}
\item \textbf{Paix négative :} simple absence de guerre ou de violence directe entre des groupes ou des États ; elle peut coexister avec des injustices latentes. \textbf{Paix positive :} état de société fondé sur la justice sociale, l'équité, le respect des droits humains et le dialogue, qui supprime les causes profondes des conflits.
\item \textbf{Extrémisme violent :} ensemble d'idéologies et de comportements qui rejettent le dialogue et recourent à la violence (attentats, enlèvements, destructions) pour imposer des convictions religieuses, politiques ou identitaires.
\item \textbf{Vivre ensemble (contexte camerounais) :} capacité des citoyens de cultures, de langues et de religions différentes à cohabiter harmonieusement, dans le respect mutuel, au sein d'une même nation unie et indivisible.
\item \textbf{Médiation :} mode de règlement pacifique dans lequel un tiers neutre et impartial aide les parties en conflit à dialoguer et à parvenir elles-mêmes à un accord acceptable par tous, sans leur imposer de solution.
\end{enumerate}
\end{corrige}

\begin{corrige}[Exercice 4 — Calcul et interprétation de données statistiques]
\begin{enumerate}
\item Le pourcentage d'une modalité se calcule par : $\text{pourcentage} = \dfrac{\text{effectif}}{\text{effectif total}} \times 100$.
\begin{align*}
\text{Étoile jaune} &: \frac{80}{200} \times 100 = 40\,\% \\
\text{Colombe blanche} &: \frac{50}{200} \times 100 = 25\,\% \\
\text{Main serrée} &: \frac{40}{200} \times 100 = 20\,\% \\
\text{Arbre à palabres} &: \frac{30}{200} \times 100 = 15\,\%
\end{align*}
Vérification : $40 + 25 + 20 + 15 = 100\,\%$. Le calcul est exact.
\item \textbf{Conclusion :} la majorité des élèves (40\,\%) identifie d'abord le symbole national officiel (l'étoile du drapeau), ce qui montre une bonne connaissance des symboles de l'État. Mais 60\,\% des réponses privilégient des symboles universels (colombe, poignée de main) ou locaux (arbre à palabres) : la paix est donc perçue à la fois comme une valeur nationale et comme une pratique quotidienne de dialogue et de réconciliation.
\end{enumerate}
\textbf{Points de vigilance :} toujours vérifier que la somme des pourcentages vaut 100\,\% ; citer les effectifs dans l'interprétation.
\end{corrige}

\begin{corrige}[Exercice 5 — Analyse de document : L'arrêt Bakassi (CIJ, 2002)]
\begin{enumerate}
\item Les deux parties en présence sont le \cle{Cameroun} et le \cle{Nigeria}, qui se disputaient la souveraineté sur la péninsule de Bakassi.
\item Il s'agit du \cle{règlement judiciaire} : le différend a été soumis à une juridiction internationale, la Cour Internationale de Justice, dont l'arrêt (10 octobre 2002) s'impose juridiquement aux deux États. Le texte le montre : « La Cour constate... confère la souveraineté... au Cameroun ».
\item Deux leçons pour l'Afrique :
\begin{itemize}
\item Les conflits frontaliers, hérités souvent de la colonisation, peuvent être réglés \emph{par le droit} et non par les armes : le recours à la justice internationale est une voie crédible.
\item Le respect effectif des décisions de justice (retrait ordonné des forces nigérianes, accords de Greentree de 2006) montre que la paix durable exige la bonne foi des États et l'accompagnement de la communauté internationale.
\end{itemize}
\end{enumerate}
\textbf{Points de vigilance :} bien nommer le mode de règlement (judiciaire, et non médiation) et le justifier par une citation du document.
\end{corrige}

\begin{corrige}[Exercice 6 — Étude de cas : Conflit intercommunautaire fictif]
\begin{enumerate}
\item Démarche de médiation en 4 étapes :
\begin{center}
\begin{tabularx}{\linewidth}{|c|X|X|X|}
\hline
\rowcolor{popblueL} \textbf{Étape} & \textbf{Acteur principal} & \textbf{Action menée} & \textbf{Objectif visé} \\
\hline
1. Désignation & Sous-préfet / autorité administrative & Nommer un médiateur neutre et accepté des deux camps (ex. notable extérieur au village) & Garantir l'impartialité, absente lors de l'arbitrage du chef \\
\hline
2. Écoute séparée & Médiateur & Rencontrer séparément éleveurs et agriculteurs pour recueillir griefs et besoins & Faire baisser la tension, comprendre les causes réelles \\
\hline
3. Négociation commune & Médiateur + délégués des deux communautés & Organiser une rencontre : délimitation des couloirs de transhumance, calendrier, indemnisation des dégâts & Construire un accord accepté par tous \\
\hline
4. Accord et suivi & Médiateur, chef, comité mixte & Rédiger un procès-verbal d'accord, créer un comité de suivi mixte, lever le blocage du marché & Rendre l'accord durable et prévenir les rechutes \\
\hline
\end{tabularx}
\end{center}
\item Voir tableau ci-dessus (acteur, action, objectif pour chaque étape).
\item La médiation est plus adaptée que l'arbitrage imposé car l'arbitrage du chef a été \emph{rejeté} pour partialité : une solution imposée ne serait pas respectée et raviverait le conflit. La médiation fait participer les deux parties à l'élaboration de la solution, ce qui restaure la confiance et garantit l'adhésion durable à l'accord.
\end{enumerate}
\textbf{Points de vigilance :} la neutralité du médiateur est la condition essentielle à citer ; l'accord final doit être un compromis, pas une victoire d'un camp.
\end{corrige}

\begin{corrige}[Exercice 7 — Stratégies de lutte contre l'extrémisme violent au Cameroun]
\begin{center}
\begin{tabularx}{\linewidth}{|X|X|X|}
\hline
\rowcolor{popblueL} \textbf{Stratégie} & \textbf{Niveau d'action} & \textbf{Exemple concret au Cameroun} \\
\hline
Opérations militaires (Force Multinationale Mixte) & Étatique (et coopération régionale) & Opérations contre Boko Haram dans l'Extrême-Nord, en coopération avec le Nigeria, le Tchad et le Niger \\
\hline
Programmes de déradicalisation et réinsertion & Étatique & Centres de désengagement de Mérina et de Mora accueillant d'ex-combattants repentis \\
\hline
Éducation à la citoyenneté et au vivre-ensemble & Individuel et communautaire & Cours d'ECM au lycée, clubs de paix, campagnes scolaires contre la radicalisation \\
\hline
Veille citoyenne et comités de vigilance & Communautaire & Comités de vigilance locaux dans l'Extrême-Nord appuyant les forces de défense et signalant les infiltrations \\
\hline
\end{tabularx}
\end{center}
\textbf{Points de vigilance :} ne pas réduire la lutte à la seule réponse militaire ; le cours insiste sur l'approche globale (sécuritaire, éducative, socio-économique).
\end{corrige}

\begin{corrige}[Exercice 8 — Travail pratique noté : Diagramme circulaire et commentaire]
\begin{enumerate}
\item \textbf{Calcul des angles} (angle $= \dfrac{\text{effectif}}{200} \times 360°$) :
\begin{align*}
\text{Étoile jaune} &: \frac{80}{200} \times 360 = 144° \\
\text{Colombe blanche} &: \frac{50}{200} \times 360 = 90° \\
\text{Main serrée} &: \frac{40}{200} \times 360 = 72° \\
\text{Arbre à palabres} &: \frac{30}{200} \times 360 = 54°
\end{align*}
Vérification : $144 + 90 + 72 + 54 = 360°$. Le calcul est exact.
\textbf{Construction attendue :} tracer un cercle, reporter au rapporteur les secteurs de 144°, 90°, 72° et 54° avec quatre couleurs distinctes ; ajouter un titre (« Symboles de la paix cités par 200 élèves d'un lycée technique de Yaoundé ») et une légende.
\item \textbf{Commentaire (5 lignes) :} Le diagramme montre la nette prédominance de l'étoile jaune du drapeau (40\,\%), symbole officiel de l'unité et de la paix nationales. La colombe blanche (25\,\%), symbole universel, arrive en deuxième position. La surprise vient du poids des symboles relationnels et locaux (poignée de main 20\,\%, arbre à palabres 15\,\%), qui rappellent que la paix se vit aussi au quotidien. Ainsi, les jeunes articulent symboles officiels, universels et traditionnels de la paix.
\end{enumerate}
\textbf{Barème indicatif (sur 10) :} calculs des angles exacts : 4 pts ; construction soignée (cercle, rapporteur, couleurs) : 3 pts ; titre et légende : 1 pt ; commentaire pertinent : 2 pts.
\textbf{Points de vigilance :} vérifier que la somme des angles fait 360° ; ne pas oublier titre et légende.
\end{corrige}

\begin{corrige}[Exercice 9 — Sujet de dissertation guidée type Probatoire]
\textbf{Corrigé structuré (éléments attendus) :}
\begin{enumerate}
\item \textbf{Introduction.} Accroche : les conflits armés font des ravages en Afrique, mais leur absence suffit-elle à parler de paix ? Définitions : la \cle{paix négative} est l'absence de guerre et de violence directe ; la \cle{paix positive} est la présence de justice, d'équité et de droits humains. Problématique : la paix est-elle seulement l'absence de guerre ? Annonce du plan : nous verrons d'abord que la paix est d'abord l'absence de violence directe, puis qu'elle est surtout une construction structurelle permanente.
\item \textbf{Première partie : la paix comme absence de violence directe.}
\begin{itemize}
\item Argument 1 : la fin d'une guerre ouverte est déjà un bien précieux. Exemple : les cessez-le-feu et la fin des affrontements ouverts liés à Boko Haram dans certaines localités de l'Extrême-Nord permettent le retour des populations.
\item Argument 2 : les traités et règlements judiciaires mettent fin aux guerres entre États. Exemple : l'arrêt de la CIJ de 2002 sur Bakassi a évité une guerre Cameroun--Nigeria.
\end{itemize}
\item \textbf{Deuxième partie : la paix comme construction structurelle.}
\begin{itemize}
\item Argument 1 : sans justice ni équité, l'absence de guerre cache des tensions prêtes à exploser. Exemple : la CNPBM (loi de 2017) promeut le bilinguisme et le multiculturalisme pour traiter les frustrations linguistiques et culturelles.
\item Argument 2 : la paix positive se construit par l'éducation et le développement. Exemple : l'ECM à l'école, la décentralisation et la lutte contre le chômage des jeunes tarissent les causes profondes des conflits.
\end{itemize}
\item \textbf{Conclusion.} Bilan : la paix n'est pas seulement l'absence de guerre ; elle est un édifice quotidien de justice et de vivre-ensemble. Ouverture : les jeunes des lycées techniques, futurs professionnels, consolident la paix positive par le travail, le respect de l'autre et le refus des discours de haine.
\end{enumerate}
\textbf{Barème indicatif (sur 20) :} introduction (accroche, définitions, plan) : 4 pts ; partie 1 (2 arguments + exemples) : 5 pts ; partie 2 (2 arguments + exemples) : 5 pts ; conclusion avec ouverture : 3 pts ; qualité de la langue et de l'organisation : 3 pts.
\textbf{Points de vigilance :} annoncer le plan en introduction ; illustrer chaque argument par un exemple précis ; ne pas opposer les deux parties mais montrer que la paix positive complète la paix négative.
\end{corrige}

\begin{corrige}[Exercice 10 — Sujet de synthèse type Baccalauréat Technique]
\begin{enumerate}
\item \textbf{Thèse principale de chaque document :}
\begin{itemize}
\item Document 1 : l'État affirme mener une réponse \emph{globale} (militaire, humanitaire, éducative, idéologique) avec des résultats concrets (plus de 2000 ex-combattants accueillis).
\item Document 2 : malgré ces efforts, les causes profondes (pauvreté, chômage, sentiment d'injustice) persistent et la réinsertion socio-économique est insuffisante.
\item Document 3 : la sécurité militaire est nécessaire mais insuffisante ; la « paix vraie » exige emploi des jeunes et justice pour tous.
\end{itemize}
\item \textbf{Forces :} (1) une approche globale combinant action militaire et déradicalisation (Doc. 1) ; (2) la protection effective des populations par l'armée (Doc. 3). \textbf{Limites :} (1) la réinsertion socio-économique des ex-combattants et des jeunes peine à suivre (Doc. 2) ; (2) la persistance de la pauvreté, du chômage et du sentiment d'injustice qui alimentent le recrutement (Doc. 2 et 3).
\item \textbf{Réponse nuancée :}
\begin{itemize}
\item \emph{Oui, en partie :} la stratégie est suffisamment \emph{large} dans ses moyens — réponse militaire (Force Multinationale Mixte), centres de désengagement de Mérina et Mora, volet humanitaire et éducatif — et elle a permis de reprendre du terrain face à Boko Haram.
\item \emph{Non, mais :} elle reste insuffisante sur le plan de la \emph{prévention structurelle} : tant que les jeunes n'ont ni travail ni confiance en la justice, le terreau du recrutement demeure. La paix positive, fondée sur l'équité et le développement, n'est pas encore acquise.
\end{itemize}
\item \textbf{Conclusion et piste concrète :} la stratégie camerounaise est nécessaire mais doit être complétée par la prévention locale. Dans un lycée technique, on peut créer un \emph{club de paix et de citoyenneté} : sensibilisation aux discours de haine sur les réseaux sociaux, tutorat des élèves en difficulté, journées du vivre-ensemble, formation à l'esprit critique face aux propagandes en ligne.
\end{enumerate}
\textbf{Barème indicatif (sur 20) :} thèses des documents : 4 pts ; forces et limites : 4 pts ; raisonnement nuancé en deux arguments : 6 pts ; conclusion et piste concrète : 3 pts ; expression : 3 pts.
\textbf{Points de vigilance :} appuyer chaque affirmation sur un document précis (citer « Doc. 1 », etc.) ; éviter la réponse manichéenne : l'examinateur attend un « oui, mais » argumenté.
\end{corrige}

\begin{corrige}[Exercice 11 — Analyse de cas complexe : Le rôle de la CNPBM dans le vivre-ensemble]
\begin{enumerate}
\item \textbf{Base légale et composition :} la CNPBM a été créée par la \cle{loi n\textsuperscript{o} 2017/014 du 12 juillet 2017}. Elle est placée sous l'autorité du Président de la République et comprend un président, un vice-président et treize membres (15 membres au total) nommés par décret, représentant les diverses sensibilités linguistiques et culturelles du pays.
\item \textbf{Trois missions en lien avec la saisine :}
\begin{itemize}
\item Promouvoir le bilinguisme et le multiculturalisme comme facteurs d'unité nationale et de paix ;
\item Recevoir les plaintes et mener des enquêtes sur les discriminations liées à la langue ou à la culture (ici, les propos haineux visant une communauté linguistique) ;
\item Formuler des recommandations au gouvernement et œuvrer à la prévention et au règlement pacifique des tensions communautaires.
\end{itemize}
\item \textbf{Communiqué de presse (proposition) :} « La Commission Nationale pour la Promotion du Bilinguisme et du Multiculturalisme, saisie par le Premier Ministre des incidents survenus lors de la compétition sportive universitaire, rappelle que le Cameroun est un État un et indivisible, fondé sur la richesse de sa diversité linguistique et culturelle, conformément à la Constitution et à la loi n\textsuperscript{o} 2017/014 du 12 juillet 2017. Elle condamne fermement les propos haineux diffusés sur les réseaux sociaux et appelle l'ensemble des étudiants et des citoyens au calme et à la retenue. Une enquête est ouverte ; une audience publique réunira les parties concernées. La Commission adressera au gouvernement des recommandations pour prévenir de tels agissements. Le dialogue reste la seule voie du vivre-ensemble. »
\item \textbf{Prévention par le dialogue plutôt que par la répression :} la CNPBM n'utilise ni arrestations ni force : elle enquête, écoute les parties (audience publique), rappelle la loi et propose des recommandations. En traitant les causes (frustrations linguistiques, discours de haine) et en restaurant le dialogue entre communautés, elle empêche l'escalade vers la violence : c'est l'illustration concrète de la \cle{paix positive} et du règlement pacifique des différends, là où la seule répression policière étoufferait les tensions sans les résoudre.
\end{enumerate}
\textbf{Barème indicatif (sur 20) :} base légale et composition : 4 pts ; missions : 5 pts ; communiqué (ton institutionnel, loi citée, appel au calme, actions) : 6 pts ; analyse dialogue/répression : 5 pts.
\textbf{Points de vigilance :} citer exactement la loi de création (2017) ; le communiqué doit rester institutionnel et apaisant, sans stigmatiser aucune communauté.
\end{corrige}

\newpage

\coursec{Fiche bilan}

\begin{popfigure}
\begin{tikzpicture}[
  mindmap,
  concept color=popblue,
  level 1 concept/.append style={level distance=4.5cm, sibling angle=360/7},
  level 2 concept/.append style={level distance=3cm},
  every node/.style={font=\small, text=white},
  grow cyclic,
  text width=2.8cm,
  align=center
]
\node[concept, minimum size=2.2cm] (centre) {La Paix}
  child[concept color=poporange] { node[concept, align=center]{1. Notion\\ \& Définition} }
  child[concept color=popgreen] { node[concept, align=center]{2. Symboles\\ au Cameroun\\(TD4)} }
  child[concept color=poppurple] { node[concept, align=center]{3. Stratégies\\ de promotion} }
  child[concept color=poppink] { node[concept, align=center]{4. Règlement\\ pacifique\\ différends} }
  child[concept color=popgold] { node[concept, align=center]{5. Lutte extrémisme\\ violent\\(Dossier 2)} }
  child[concept color=popblue] { node[concept, align=center]{6. Vivre ensemble\\ \& Bilinguisme\\(Dossier 3)} }
  child[concept color=popdark] { node[concept, align=center]{7. Citoyenneté\\ active} };
\end{tikzpicture}
\legende{Carte mentale de la leçon 05 : La notion de paix (Première Technique)}
\end{popfigure}

\begin{bilan}
\courssub{Définitions clés à connaître par cœur}
\begin{tabularx}{\linewidth}{|>{\bfseries}l|X|}
\hline
\rowcolor{popblueL}
\textbf{Terme} & \textbf{Définition / Contenu essentiel} \\
\hline
Paix & Absence de violence directe (guerre, conflit armé) \textbf{et} absence de violence structurelle (injustice, inégalité, pauvreté). Paix négative (cessez-le-feu) vs Paix positive (justice, développement, droits humains). \\
\hline
Culture de la paix & Ensemble de valeurs, attitudes, comportements reflétant le respect de la vie, la non-violence, la tolérance, le dialogue, la coopération (Résolution ONU 53/243). \\
\hline
Règlement pacifique & Modes : négociation, médiation, conciliation, arbitrage, règlement judiciaire (CJI, CEDH). Principe : consentement des parties, impartialité, respect du droit. \\
\hline
Extrémisme violent & Idéologie prônant l'usage de la violence pour imposer des vues politiques, religieuses, sociales. Recrutement, radicalisation, passage à l'acte terroriste. \\
\hline
Vivre ensemble & Construction quotidienne du lien social par le respect de la diversité (ethnique, religieuse, linguistique), la solidarité, le partage des valeurs républicaines. \\
\hline
Bilinguisme officiel & Français \& Anglais langues officielles (Art. 1(3) Constitution). Outil d'intégration nationale, d'accès à l'emploi, de cohésion. \\
\hline
Commission Nationale Bilingualisme \& Multiculturalisme (CNBBM) & Organe consultatif (Décret 2017/013). Missions : promouvoir le bilinguisme, le multiculturalisme, recevoir plaintes, proposer mesures. \\
\hline
\end{tabularx}

\courssub{Textes de référence incontournables}
\begin{itemize}
\item \textbf{Constitution du Cameroun (1996, révisée 2008)} : Préambule (droits, devoirs, bilinguisme, décentralisation), Art. 1(3) (bilinguisme), Art. 65 (décentralisation).
\item \textbf{Charte des Nations Unies} : Chap. VI (règlement pacifique), Résolution 53/243 (Culture de la paix).
\item \textbf{Déclaration universelle des droits de l'homme (1948)} : Art. 1, 3, 26(2) (éducation paix).
\item \textbf{Loi n°2014/028} : Code pénal (sanctions terrorisme, apologie, extrémisme).
\item \textbf{Décret n°2017/013} : Création, organisation, fonctionnement CNBBM.
\item \textbf{Plan d'urgence triennal (PLANUT)} \& \textbf{Stratégie nationale de lutte contre l'extrémisme violent}.
\end{itemize}

\courssub{Méthodes express (Bac)}

\textbf{Méthode 1 : Analyser un conflit pour proposer une issue pacifique}
\begin{enumerate}
\item \textbf{Identifier} les acteurs, causes (structurelles, conjoncturelles), enjeux.
\item \textbf{Qualifier} le type de violence (directe, structurelle, culturelle).
\item \textbf{Choisir} le mode de règlement adapté (négociation locale, médiation administrative, justice coutumière, tribunal).
\item \textbf{Argumenter} avec les textes (Constitution, Code pénal, Résolutions ONU) et les valeurs (dialogue, pardon, justice réparatrice).
\item \textbf{Conclure} sur la durabilité : paix négative (arrêt combats) vs paix positive (réconciliation, développement).
\end{enumerate}

\textbf{Méthode 2 : Exploiter un document (texte, image, statistique) sur l'extrémisme ou le vivre ensemble}
\begin{enumerate}
\item \textbf{Nature/Source} : Auteur, date, contexte (crise NOSO, Boko Haram, réseaux sociaux).
\item \textbf{Thème central} : Radicalisation, déscolarisation, discours de haine, rôle CNBBM, bilinguisme à l'école.
\item \textbf{Information clé} : Chiffre, citation, mesure concrète (ex: centres de déradicalisation, quotas bilinguisme concours).
\item \textbf{Analyse critique} : Efficacité, limites, lien avec les valeurs citoyennes (tolérance, patriotisme, respect loi).
\item \textbf{Proposition personnelle} : Action jeune (club paix, signalement, apprentissage L2, médiation pairs).
\end{enumerate}
\end{bilan}

\begin{aretenir}[Checklist avant le Bac]
$\square$ Définir la paix (négative/positive) et la culture de la paix (valeurs ONU).\\
$\square$ Citer 4 symboles de la paix au Cameroun (ex: colombe, main tendue, hymne, drapeau, monuments, journée 21 sept).\\
$\square$ Différencier négociation, médiation, conciliation, arbitrage, justice (tableau comparatif).\\
$\square$ Expliquer le processus de radicalisation (4 stades) et citer 3 facteurs de vulnérabilité (pauvreté, désœuvrement, discours haine, échec scolaire).\\
$\square$ Citer 3 missions de la CNBBM et 2 pouvoirs (enquête, recommandation, saisine).\\
$\square$ Expliquer le lien Bilinguisme = Vivre ensemble = Unité nationale (accès service public, concours, éducation).\\
$\square$ Identifier les piliers de la stratégie nationale anti-extrémisme (prévention, protection, poursuite, partenariat, réinsertion).\\
$\square$ Analyser un cas concret : proposer une médiation scolaire ou familiale (étapes, acteurs, issue gagnant-gagnant).\\
$\square$ Rédiger une argumentation structurée : « La paix est un devoir citoyen » (thèse, arguments juridiques/moraux, exemples Cameroun, conclusion).\\
$\square$ Connaître les numéros d'alerte (1501, 112, 116) et la conduite à tenir face à une suspicion de radicalisation (signaler, ne pas juger, protéger).\\
$\square$ Maîtriser le vocabulaire technique : résilience, cohésion sociale, dialogue intergénérationnel, justice transitionnelle, apologie du terrorisme.
\end{aretenir}

\findecours

\end{document}
